\documentclass{mathnote_cmyk}
\usepackage{mathnote_cmyk}

% 製本用：表紙の裏・裏表紙の裏は完全な白紙にする
\MathNoteVersoLabelfalse
% 製本用：総ページ数（表紙・裏表紙を含む）．4 の倍数にしておく
\newcommand{\BookTotalPages}{112}

\begin{document}

% タイトルページの呼び出し
\tituloum
  {Foundations of Mathematics}
  {Namachan}
  {Mathematics Notes}
  {An introduction to mathematical language, logic, proofs, sets, and functions}
  {\today}

%\include{Preambulo/Capa}

\cleardoublepage
\pagenumbering{arabic}
\pagecolor{white}

% 目次の表示
\localtableofcontents


\phantomsection
\part{用語の定義}

\begin{abstract}
  この章では，数学の学習を始める前に，定義・命題・公理・定理・補題・系などの基本的な用語について理解を深める．「『命題』という言葉は知っているけど，『定理』とはどう違うの？」という疑問を持ったり，「『定義』と『定理』の違いがあやふやなままで，よくわからない」と感じる学生は多い．まず，日常的な例と数学的な例をもとに「定義」の性質を確認する．次に，命題・定理・補題・系という用語を整理し，これらのラベルの使い分けについて注意を述べる．最後に，「かつ・または」「とる・おく・固定する」「十分大きい」など，数学に特有の言い回しを具体例とともに確認しよう．
\end{abstract}

\phantomsection
\section{定義・命題・公理・定理・補題・系}

\phantomsection
\subsection{定義}

数学用語には細かな意味がある．概要にも述べたように，「定義」と「定理」の違いなどをはっきりと認識しておくことは，
数学の学習を進める上で堅固たる「理解の基盤」となる．まずは「定義」の意味の確認からはじめよう：

\emph{定義}\index{ていぎ@定義}（definition）とは，用語の意味を明確に述べたものであり，Defと略記される．
要するに「定義」とは，「この数学的概念（用語）はこのような意味である」という取り決めである．
数学のお話にうつる前に，日常的な「定義」の例を見てみよう．

\begin{example}{アレニウスによる酸・塩基の定義}{アレニウスによる酸・塩基の定義}
  水溶液中で水素イオン\ce{H+}を放出する物質を酸，\ce{OH-}を放出する物質を塩基とする．
\end{example}

\begin{example}{「民主主義」の定義}{「民主主義」の定義}
  民主主義とは，国民が政治の意思決定に直接または間接的に参加する政治体制のことをいう．
\end{example}

どちらもわかりやすい例であり，「すでに知っていること」と思われる読者もいるかもしれない．
しかし，議論を行ううえで，定義が曖昧なままだと不都合が生じることがある．

\begin{example}{「速度」と「速さ」の違い}{「速度」と「速さ」の違い}
直線上を運動する物体の位置を$x$とすると，その速度$v$\footnote{空間内の運動では，速度を$\vec{v}$や$\bm{v}$などと表すこともある．一方，大学の数学や物理では，文脈から明らかな場合にはベクトルとスカラーを表記上区別しないことも多い．}は
\[
v=\frac{dx}{dt}
\]
と定義される．ここで，「速度」と「速さ」を同じ意味で用いたとしよう．

速度$v$には正負があり，その符号は物体がどちらの向きに運動しているかを表す．一方，速さは運動の方向を考えず，速度の大きさ
\[
\abs{v}
\]
で表される．

たとえば，速度が$v=-3$である物体の速さは$3$である．ここで速度と速さを混同して$v=3$としてしまうと，物体が実際とは反対向きに運動していることになってしまう．
このように，用語の意味を曖昧にすると必要な情報が失われ，まったく異なる結論に至ることがある．
\end{example}

\begin{example}{「平均」と「期待値」の違い}{「平均」と「期待値」の違い}
  ここで比較するのは，算術平均（標本平均）と期待値である．算術平均とは，実際に得られたデータをならした値である．
  たとえば，$n$個のデータ $x_1,x_2,\dots,x_n$ の算術平均は
  \[
    \frac{x_1+x_2+\cdots+x_n}{n}
  \]
  で与えられる．

  算術平均の定義に基づくと，$1,4,10$ という3個のデータの算術平均は
  \[
    \frac{1+4+10}{3}=5
  \]
  である．

  一方，期待値とは，各値にその値が出る確率をかけて足し合わせたものである．
  たとえば，ある確率変数 $X$ が $1,4,10$ をそれぞれ
  $\frac{1}{6},\frac{1}{6},\frac{2}{3}$ の確率でとるとき，
  その期待値は
  \[
    E[X]=1\cdot\frac{1}{6}+4\cdot\frac{1}{6}+10\cdot\frac{2}{3}=7.5
  \]
  である．

  このように，算術平均は実際のデータから求める値であり，
  期待値は確率を用いて求める値である\footnote{ただし，確率論では期待値のことを「平均」とよぶこともある．「平均」という語は文脈により指すものが異なるので，どの意味で用いられているかに注意したい．}．
\end{example}

\exref{「速度」と「速さ」の違い}や\exref{「平均」と「期待値」の違い}からもわかるように，
数学の学習において「定義」をあやふやにしてはならない．

しかし，同じ事柄について，2つ以上の定義の形式があることもある\footnote{「書籍によって採用している定義が違う」というのもよくあることである．}．次は数学における例を見てみよう：

\phantomsection
\begin{example}{「絶対値」の定義}{「絶対値」の定義}
  $ x \in \mathbb{R}$に対して，$x$の絶対値を$\abs{x}$とかき，次のように定義する：
  \[
    \abs{x} \coloneqq \max \{ x , -x \}
  \]
  この定義は以下のような形式で表現してもよい：
  \[
    \abs{x} \coloneqq  \sqrt{x^2}
  \]
\end{example}

つまり，$x\in\mathbb{R}$の絶対値$\abs{x}$は，$\abs{x} \coloneqq \max \{ x, -x \}$としてもよい一方，
$\abs{x} \coloneqq \sqrt{x^2}$と定義しても問題ないのである\footnote{もちろんこの他の定義を採用してもよい．たとえば$x\geqq 0$のとき$\abs{x} \coloneqq x$，$x<0$のとき$\abs{x} \coloneqq -x$とする定義がある．}

また，\exref{「絶対値」の定義}のように，同じ概念にふたつの定義の形式があるとき，
「どちらの形式を採用しても同じ概念が定まる」こと，すなわち「ふたつの定義が同値である」ということは，
証明を要する主張である\footnote{「証明」とはなにか，どのように書くのかについては\partref{証明}で詳しく扱う．
本章では「定義や既知の事実から，筋道を立てて主張が正しいことを確かめる手続き」という程度の理解でよく，
本章に現れる証明は書き方の雰囲気をつかむための例と思って読み進めてほしい．細部は\partref{証明}を読んだあとで改めて確認すればよい．}．
片方を定義として採用した場合には，もう片方は「証明すべき命題」という扱いになるのである．
実際に確かめてみよう．
 
\begin{prop}{絶対値のふたつの定義の同値性}{絶対値のふたつの定義の同値性}
  任意の$x \in \mathbb{R}$に対して，以下の等式が成り立つ：
  \[
    \max \{ x , -x \} = \sqrt{x^2}
  \]
  ただし，$a \geqq 0$に対して$\sqrt{a}$は「$2$乗すると$a$になる非負の実数」を表す\footnote{この記法が正当であるためには，「$2$乗すると$a$になる非負の実数」がただひとつ存在することが必要である．このことについては，のちの「一意性」の節でも触れる．}．
\end{prop}
 
\begin{leftbar}
  \begin{proof}
    $x$の符号で場合分けする．
      \proofstep{$x \geqq 0$の場合}
            $-x \leqq 0 \leqq x$により，$\max \{ x,-x \} = x$である．
            一方，$x \geqq 0$かつ$x^2 = x \cdot x$なので，$\sqrt{x^2}=x$である．
            よって両辺は一致する．
      \proofstep{$x < 0$の場合}
            $x < 0 < -x$により，$\max \{ x,-x \} = -x$である．
            一方，$-x > 0$かつ$x^2 = (-x)^2$なので，$\sqrt{x^2}=-x$である．
            よって両辺は一致する．
            
            \smallskip
    以上により，任意の$x \in \mathbb{R}$に対して主張の等式が成り立つ．これが証明すべきことであった．
  \end{proof}
\end{leftbar}

このように，\prref{絶対値のふたつの定義の同値性}に証明を与えたことで，実数$x$の絶対値を$\abs{x} \coloneqq \max \{x,-x\}$と定義しても，$\abs{x} \coloneqq \sqrt{x^2}$と定義しても，同じものが得られることがわかった．すなわち，どちらの定義を採用してもよい．


次に，具体例を通して「定義の性質」についてみていこう．具体的には「定義を設ける上での慣習」について確認をする．
線型代数の講義で学ぶことであるが，逆行列の定義は以下のようになる：

\begin{example}{「逆行列」の定義}{「逆行列」の定義}
  正方行列$A$に対して，$BA = AB =E$となるような$B$が存在するとき，このような$B$を$A$の逆行列という．
\end{example}


ここで注意するのは，数学では定義は最小限の情報にとどめることが慣習となっているということである．たとえば，\exref{「逆行列」の定義}から，以下の命題\footnote{「命題」については，のちほど詳しくみていくとする．ここでは「証明を与えるべきである主張」という理解でよい．}が成り立ち，その証明は容易である．

\phantomsection
\begin{prop}{逆行列の一意性}{逆行列の一意性}
  正方行列$A$に対して，$A$の逆行列は存在するとすればただひとつである．
\end{prop}

\begin{tleftbar}
  \begin{proof}
    $A$の逆行列が$B,C$であるとすると，
    \[
      AB=BA=E,\quad AC=CA=E
    \]
    が成り立つ．このとき，
    \[
      B=BE=B(AC)=(BA)C=EC=C
    \]
    が成り立つ．よって，$B=C$であり，逆行列の一意性が示された．これが証明すべきことであった．
  \end{proof}
\end{tleftbar}

このことから，\exref{「逆行列」の定義}でとりあげた逆行列の定義をもっと詳しく

\begin{quote}
  正方行列$A$に対して，$BA = AB =E$となるような$B$が存在するとき，このような$B$はただひとつで，$B$を$A$の逆行列という．
  $B$は一意に存在するので，これを$A^{-1}$とかく\footnote{一意に存在することがわかっていないと，$A^{-1}$のような記法で表すことはためらわれる．}．もちろん
  \[
    AA^{-1}=A^{-1}A=E
  \]
  が成り立つ．
\end{quote}
と，「逆行列の一意性を定義に含めてもいいのではないか．」と主張する人がいるかもしれない．
ただ，数学では「定義は最小限の情報にとどめ，そこから導かれる主張を命題として証明する」という慣習があり，
なにが「定義」で，なにが「証明すべきこと」であるかはっきりさせることが多い\footnote{ただ，実際はここで取り上げた逆行列の定義も情報過多である．線型代数で学ぶことになるが，$AB=E$と$BA=E$のどちらか片方の式のみで定義してよいからである．}．

もうひとつ，定義について注意しておきたいことがある．
それは，適切に与えられた定義は，通常，真偽を問う対象ではなく，理論全体をうまく構築できるように「設定」されるということである\footnote{「通常」と述べたのは，「対象を定められていない不適切な定義」がありうるからである．このことは，のちに述べるwell-definedの概念と関わる．}．
 
\phantomsection
\begin{example}{「$1$は素数でない」という約束}{「1は素数でない」という約束}
  素数は「$1$とその数自身以外に正の約数を持たない，$2$以上の自然数」と定義され，$1$は素数から除外される．
  これは「$1$を素数に含めるとなにか矛盾が生じるから」ではない．素因数分解の一意性
  \begin{quote}
    $2$以上の任意の自然数は，順序の違いを除いてただひととおりに素数の積として表される．
  \end{quote}
  という定理をきれいな形で述べるためである．もし$1$を素数に含めると，
  \[
    6 = 2 \cdot 3 = 1 \cdot 2 \cdot 3 = 1 \cdot 1 \cdot 2 \cdot 3 = \cdots
  \]
  となり，「ただひととおり」という主張が崩れてしまう\footnote{正確には，一意性の主張に「$1$の個数の違いを除いて」といった但し書きを毎回つける必要が生じる．}．
\end{example}

このように，「$1$は素数か」という問いに対する答えは，証明によって得られるのではなく，
「どのような定義を採用すると理論をうまく構築できるか」という観点から選ばれている．
$0\kaijou = 1$や$a^0 = 1$（$a \ne 0$）という約束も，既存の法則（$n\kaijou = n \cdot (n-1)\kaijou$や指数法則$a^{m+n} = a^m a^n$など）を例外なく記述できるようにする，という発想による．

また，新しい概念や関数を定義するときには，その定義がwell-defined\index{welldefined@well-defined}であることを確認する必要がある．
たとえば，有理数全体の集合を$\mathbb{Q}$とし，関数
\[
  f \colon \mathbb{Q} \to \mathbb{Z}
\]
を
\[
  f\left(\frac{a}{b}\right)=a+b
\]
によって定義したとしよう．このとき，
\[
  f\left(\frac12\right)=1+2=3
\]
である一方，
\[
  f\left(\frac24\right)=2+4=6
\]
となる．ところが，
\[
  \frac12=\frac24
\]
であるから，これは$f$が同じ有理数に対して異なる値を与えていることになる．一般に，
\[
  \frac12=\frac24=\frac36=\cdots
\]
のように，有理数は複数の表し方をもつ．同じ有理数をどのように表しても関数の値が一致しなければ，その規則は「有理数を定義域とする関数」を定義していることにはならない．
このような事例を踏まえ，対象の表し方によらず値が一意に定まることを，well-defined（よく定義されている）という．たとえば，
\[
g \colon \mathbb{Q} \to \mathbb{Q}
\]
を
\[
g\!\left(\frac{a}{b}\right)=\frac{2a}{b}
\]
と定めると，$g$はwell-definedである．

well-definedについては，\subsecref{well-defined}で再度確認をするので，ここではこの程度の説明に留めよう．

\phantomsection
\subsection{命題}

数学でいう\emph{命題}\index{めいだい@命題}（proposition）という語には，「広い意味」と「狭い意味」のふたつの用法がある．
広い意味では，真であるか偽であるかが定まる文を命題という\footnote{たとえば，「$1+1=2$」は真の命題であり，「$1+1=3$」は偽の命題である．真であっても偽であっても，真偽が定まってさえいれば命題とよぶことに注意しよう．}．
一方，狭い意味では，「定理」「補題」「系」などと同じように，正しい数学の主張につけるラベルのひとつとして「命題」という語を用いる\footnote{前項の\prref{絶対値のふたつの定義の同値性}や\prref{逆行列の一意性}に「命題」という見出しをつけたのは，この用法である．}．

\begin{figure}[h]
  \centering
  % 距離は A4 実寸で書き \mnsc で B5 版に換算する（transform shape を使うと
  % B5 版では本文側で縮んだ文字がさらに縮むため使わない）．
  % 「命題」の 2 つの箱をマゼンタで強調する．
  \begin{tikzpicture}[
    node distance=\mnsc{4.5mm} and \mnsc{9mm},
    flowbox/.append style={font=\footnotesize},
    flowq/.append style={font=\footnotesize},
    flowlab/.append style={font=\footnotesize},
    flowhl/.append style={font=\footnotesize\bfseries, draw=figaccent,
                          fill=figaccent!8, text=figaccent!80!black},
  ]
    \node[flowbox] (sent) {言葉・式};
    \node[flowq, below=of sent] (q0) {主張を表しているか？};
    \node[flowbox, right=\mnsc{10.8mm} of q0] (nosent)
      {命題ではない\\ \scriptsize（「$1+1$」「りんご」など）};

    \node[flowq, below=of q0] (q1) {真偽が定まるか？};
    \node[flowbox, right=\mnsc{10.8mm} of q1] (notprop)
      {命題ではない\\ \scriptsize（「この部屋は暑い」など）};

    \node[flowbox, flowhl, below=of q1] (broad) {命題（広い意味）};
    \node[flowq, below=of broad] (q2) {真か偽か？};

    \node[flowbox, left=\mnsc{9mm} of q2, yshift=\mnsc{-11.7mm}] (false)
      {偽の命題\\ \scriptsize（「$1+1=46$」など）};
    \node[flowbox, right=\mnsc{9mm} of q2, yshift=\mnsc{-11.7mm}] (true)
      {真の命題\\ \scriptsize（「$1+1=2$」など）};

    \node[flowq, below=of true] (q3) {証明を与え，\\見出しをつけるか？};

    \node[flowbox, flowhl, below=\mnsc{7.2mm} of q3, xshift=\mnsc{-12.6mm}] (narrow)
      {命題（狭い意味）};
    \node[flowlab, right=\mnsc{2.7mm} of narrow] (thm) {定理};
    \node[flowlab, right=\mnsc{2.7mm} of thm]    (lem) {補題};
    \node[flowlab, right=\mnsc{2.7mm} of lem]    (cor) {系};

    \node[
      draw=appleGray,
      dashed,
      rounded corners=3pt,
      inner sep=\mnsc{4.5pt},
      fit=(narrow)(thm)(lem)(cor),
      label={[flowtxt]below:
        正しいと分かっている主張につける見出し（ラベル）}
    ] (labels) {};

    \draw[flowarr] (sent) -- (q0);
    \draw[flowarr] (q0) -- node[flowtxt, above] {いいえ} (nosent);
    \draw[flowarr] (q0) -- node[flowtxt, right] {はい} (q1);

    \draw[flowarr] (q1) -- node[flowtxt, above] {いいえ} (notprop);
    \draw[flowarr] (q1) -- node[flowtxt, right] {はい} (broad);

    \draw[flowarr] (broad) -- (q2);
    \draw[flowarr] (q2) -|
      node[flowtxt, pos=0.25, above] {偽} (false);
    \draw[flowarr] (q2) -|
      node[flowtxt, pos=0.25, above] {真} (true);

    \draw[flowarr] (true) -- (q3);
    \draw[flowarr] (q3) --
      node[flowtxt, right] {はい} (labels.north -| q3);
  \end{tikzpicture}
  \caption{「命題」のふたつの意味}
  \label{fig:命題のふたつの意味}
\end{figure}

本項では「広い意味での命題」について解説する．そして，本書ではどちらの意味であるかが明らかな場合は，単に「命題」とかくと約束する．

「広い意味での命題」は「真の命題」と「偽の命題」のいずれかであり，このようにふたつの値を用いる論理体系を\emph{二値論理}\index{にちろんり@二値論理}とよぶ．
以下では，命題が真であることを$\mathrm{T}$，偽であることを$\mathrm{F}$と表す．

命題の定義にあてはまらないとされる主張には

\begin{enumerate}[(1)]
\item 用語や対象の定義が曖昧なもの
  \label{enu:定義が曖昧なもの}
\item 文の意味や解釈が曖昧なもの
  \label{enu:意味が曖昧なもの}
\item 真偽を問うことのできない文
  \label{enu:真偽を問えないもの}
\item 真偽が一意に定まらないもの
  \label{enu:真偽が定まらないもの}
\item パラドックス
  \label{enu:パラドックス}
\end{enumerate}

などがある．\ref{enu:定義が曖昧なもの}，\ref{enu:意味が曖昧なもの}，\ref{enu:真偽を問えないもの}，\ref{enu:真偽が定まらないもの}についてはのちほど説明するとして，ここでは\ref{enu:パラドックス}の例をあげよう．

\phantomsection
\begin{example}{自己言及のパラドックス}{自己言及のパラドックス}
次のような主張を考える：
\begin{quote}
この文は偽である．
\end{quote}
この文が真であるとすれば，その主張のとおり偽となってしまい，偽であるとすれば，主張どおりであるから真となってしまう．この文は嘘つきのパラドックスとして知られる自己言及のパラドックスであり，通常の二値論理の枠組みでは，そのまま命題として扱うことができない．
\end{example}

ここまで，二値論理では取り扱わない主張について述べてきたが，ここからは二値論理で取り扱う主張に戻ろう．

\phantomsection
\begin{example}{命題}{命題}
次に示す文は真の命題である：
\begin{enumerate}[(1)]
\item 「$1+1=2$」
\item 「$23 > 17$」
\item 「円周率は100未満である」
\item 「1日は24時間である」
\end{enumerate}
また，次に示す文は偽の命題である：
\begin{enumerate}[(a)]
\item 「$1+1=46$」
\item 「$23>26$」
\item 「円周率は3未満である」
\item 「1日は30時間である」
\end{enumerate}
\end{example}

\phantomsection
\begin{example}{命題でない文}{命題でない文}
次に示す文は命題ではない：
\begin{enumerate}[(A)]
\item 「$1+1$」（なにも主張しておらず，真か偽か判定できない）
\item 「この部屋は暑い」（「暑い」の基準が明確でなく，真か偽かが定まらない）
\item 「$10000$は大きい」（「大きい」の基準が明確でなく，真か偽かが定まらない）
\item 「この料理はおいしい」（「おいしい」という評価は人によって異なり，真か偽かが定まらない）
\item 「早く起きなさい」（命令文であって，真か偽か判定できない）
\item 「$x^2 > 4$」（$x$の値が定まっていないため，真か偽かが定まらない）\label{enu:命題でない文4}
\end{enumerate}
\end{example}


\phantomsection
\subsection{公理}

\emph{公理}\index{こうり@公理}（axiom）とは，証明を与えずに正しいと認め，議論の出発点として採用する主張のことである．
Axと略記される\footnote{歴史的には，ユークリッドの『原論』のように，「公理」（どの分野にも共通する前提）と「公準」（その分野に固有の前提）を区別することもあった．現代の数学ではこの区別をせず，まとめて「公理」とよぶのが普通である．}．

解析学の教科書でよく見かける「公理」として，$a,b \in \mathbb{R}$に対して
\[
a+b = b+a
\]
などとする「加法の交換律」があるが，我々は学習の早い段階において，この「公理」を証明無しに用いていることが多い．
それゆえ，「公理とは，誰でも知っている当たり前の事柄のことだろう」と思われる方もいるのではないか．
たしかに，古くは公理を「自明な真理」と考えていた時代もあった．
しかし，現代の数学における公理はそのような意味ではない．
公理とは，推論の出発点として「この主張は正しいものとする」と取り決めた命題であって，それが自明であるかどうかとは関係がない．
「定義」の項で述べた「定義は真偽を問う対象ではなく，理論をうまく構築できるように設定される」という事情と同じ構図である．
定義が用語についての取り決めであるのに対し，公理は主張についての取り決めなのである．

まず，公理が「当たり前の事柄」とは限らないことを示す歴史的な例を見てみよう．

\phantomsection
\begin{example}{平行線公準}{平行線公準}
  ユークリッド幾何学は，次の主張を出発点のひとつとしている\footnote{『原論』の第5公準を，のちの時代の同値な言い換え（プレイフェアの公理）の形で述べた．}：
  \begin{quote}
    直線$l$と，$l$上にない点$P$が与えられたとき，$P$を通り$l$と交わらない直線はただひとつ存在する．
  \end{quote}
  この主張は「ほかの公理から証明できるのではないか」と長らく考えられ，2000年以上にわたり多くの数学者が証明を試みた．
  しかし19世紀になって，この主張を「$P$を通り$l$と交わらない直線は2本以上存在する」という主張に取り替えても，矛盾のない幾何学（非ユークリッド幾何学）が構築できることがわかった．
  つまり平行線公準は，ほかの公理からは証明も反証もできない主張であり，「採用するかどうか」を取り決めるべき出発点だったのである．
\end{example}

次に，本書の主題にもっとも関わる例として，先にも述べた，「解析学の教科書」でよく見かける「実数の公理」を紹介しよう．
高校までは，実数を「数直線上の点」として扱い，その性質をとくに疑うことなく用いてきた．
一方，大学の解析学では，実数とは「次の公理をみたす集合$\mathbb{R}$である」と約束し，そこから議論を始める\footnote{要するに「$\mathbb{R}$は完備順序体である」という意味なのだが，各公理の正確な記述は解析学の教科書に譲り，ここでは概略のみを述べる．}．

\phantomsection
\begin{axiom}{実数の公理}{実数の公理}
  集合$\mathbb{R}$には加法$+$，乗法$\cdot$，および大小関係$\leqq$が定められていて，次をみたす：
  \begin{enumerate}[(I)]
    \item （四則演算の公理）加法と乗法について交換法則・結合法則・分配法則が成り立つ．
      加法の単位元$0$と乗法の単位元$1$（$0 \ne 1$）が存在する．
      任意の$x \in \mathbb{R}$に対して$x + (-x) = 0$となる$-x$が存在し，
      $0$でない任意の$x \in \mathbb{R}$に対して$x \cdot x^{-1} = 1$となる$x^{-1}$が存在する．\label{enu:四則演算の公理}
    \item （順序の公理）任意の$x, y \in \mathbb{R}$について，$x \leqq y$と$y \leqq x$の少なくとも一方が成り立ち，両方が成り立つのは$x = y$のときに限る．
      $x \leqq y$かつ$y \leqq z$ならば$x \leqq z$である．
      また，$x \leqq y$ならば任意の$z$について$x + z \leqq y + z$であり，さらに$z \geqq 0$ならば$xz \leqq yz$である．\label{enu:順序の公理}
    \item （連続の公理）上に有界な空でない部分集合$A \subset \mathbb{R}$は上限$\sup A$をもつ\footnote{「上に有界」「上限」の正確な意味はのちの章で述べる．ここでは「数直線に穴があいていない」ことを表す主張だと思っておけばよい．}．\label{enu:連続の公理}
  \end{enumerate}
\end{axiom}

\ref{enu:四則演算の公理}は「四則演算が自由にできる」，\ref{enu:順序の公理}は「大小が比較でき，それが四則演算と両立する」という，高校までに慣れ親しんだ性質である．
一方，\ref{enu:連続の公理}は，多くの読者にとって初めて目にする主張であろう．
これこそ，「公理は誰でも知っている事柄とは限らない」ことの典型例である．

\ref{enu:連続の公理}がなぜ必要なのかは，有理数全体$\mathbb{Q}$を見るとわかる．
$\mathbb{Q}$は\ref{enu:四則演算の公理}と\ref{enu:順序の公理}をみたすが，\ref{enu:連続の公理}をみたさない．
実際，$A = \{ x \in \mathbb{Q} \mid x^2 < 2 \}$は上に有界であるが，$\mathbb{Q}$の中には上限（$\sqrt{2}$となるべき数）が存在しない．
つまり，\ref{enu:四則演算の公理}と\ref{enu:順序の公理}だけでは$\mathbb{R}$と$\mathbb{Q}$を区別できず，「$\sqrt{2}$が存在する」ことすら導けないのである．
中間値の定理や最大値・最小値の定理\footnote{のちにロルの定理の証明で「成立を認める」として用いる定理である．}など，解析学の基本的な定理はすべて\ref{enu:連続の公理}を出発点として証明される．

また，\ref{enu:四則演算の公理}〜\ref{enu:連続の公理}をすべてみたす体系は，本質的に$\mathbb{R}$ただひとつであることが知られている\footnote{正確には「同型を除いてただひとつ」である．なお，この主張自体は公理ではなく，証明を要する定理である．}．
その意味で，\axref{実数の公理}は実数という対象を特徴づけている．

ここまでの内容を整理すると，数学の理論は
\begin{quote}
  定義（用語の取り決め）と公理（出発点となる主張の取り決め）から出発し，証明によって定理を積み上げていく
\end{quote}
という構造をもつ．
定義と公理はどちらも「取り決め」であって証明の対象ではなく，定理は証明の対象である．
公理は，議論全体を通して有効な「仮定」ともいえる\footnote{数学的帰納法や背理法で用いる一時的な「仮定」との違いについては，のちの「仮定する」の項で述べる．}．

ただし，定義と公理の境目は必ずしも明確ではない．
たとえば\ref{enu:四則演算の公理}の四則演算の公理は，代数学では「体の定義」とよばれる．
「これらの性質をみたす集合を体とよぶ」と読めば定義であり，「実数はこれらの性質をみたす」と読めば公理である．
現代の数学では，このように「公理系をみたす対象」として概念を定義する方法（公理的定義）が広く用いられる．


\phantomsection
\subsection{定理}

まず，前項で保留にしていた「狭い意味での命題」を確認しておこう．
数学では，証明によって正しいことが確かめられた主張に「命題」「定理」「補題」「系」といった見出し（ラベル）をつけて述べる．
「狭い意味での命題」とは，このラベルとしての「命題」のことであり，Propと略記される．
前項でみた「真偽が定まる文」という広い意味と異なり，こちらは「証明された（あるいは証明すべき）正しい数学の主張」を指す言葉である\footnote{したがって，狭い意味での命題は，広い意味では「真の命題」である．偽である主張に「命題」というラベルがつけられることはない．}．
両者の違いを\tabref{命題の比較}に整理した．

\begin{table}[htbp]
  \centering
  \begin{tabular}{>{\centering\arraybackslash}m{\mnsc{2.2cm}} L{4.9cm} L{4.9cm}}
    \toprule
      & \textbf{広い意味} & \textbf{狭い意味} \\
    \midrule
    なにを指すか & 真偽が定まっている文 & 証明された数学の主張につけるラベル \\
    真偽 & 真の場合も偽の場合もある & つねに真 \\
    対になる語 & 「命題でない文」 & 定理・補題・系 \\
    主な登場場面 & 論理（真理値表など） & 教科書の見出し \\
    例 & 「$1+1=2$」（真），「$1+1=46$」（偽） & \prref{逆行列の一意性} \\
    \bottomrule
  \end{tabular}
  \caption{広い意味・狭い意味の命題の比較}
  \label{tab:命題の比較}
\end{table}

そのうえで，\emph{定理}\index{ていり@定理}（theorem）とは，正しいと分かっている数学の主張（狭い意味での命題）の中でもとりわけ重要なものを指す．Thmと略記される．


定理とは，定義や公理，すでに得られた結果から論理的に証明された数学的主張であり，議論を進めるうえでの基礎となる．定理は\partref{証明}で扱う「証明」によって裏付けられているため，採用した定義・公理のもとで正しいことが保証されている．

定理が重要な主張であるがゆえに，「余弦定理」・「加法定理」・「ハイネ・ボレルの被覆定理」など，固有の名前が与えられているものも存在する．
その固有の名前は，定理の主張の詳細，もしくは発見者の名前にちなんで付けられることが多い．たとえば，「三平方の定理\footnote{「ピタゴラスの定理」とも呼ばれる．}」や「コーシー・シュワルツの不等式」である．

\phantomsection
\begin{example}{三平方の定理}{三平方の定理}
  三平方の定理の主張は以下のようになる：
  \begin{quotation}
    直角三角形の斜辺の長さを$c$とし，他の二辺の長さを$a,b$としたとき，
    \[
      a^2+b^2=c^2
    \]
    が成り立つ．
  \end{quotation}
\end{example}

\phantomsection
\begin{example}{コーシー・シュワルツの不等式}{コーシー・シュワルツの不等式}
  コーシー・シュワルツの不等式は，内積空間における基本的な不等式であり，以下のように表される：
  \begin{quotation}
    任意の内積空間において，ベクトル$u$と$v$に対して，
    \[
      \abs{\langle u, v \rangle} \leqq \norm{u} \cdot \norm{v}
    \]
    が成り立つ．
  \end{quotation}
  この不等式は解析学や線型代数学など，多くの分野で重要な役割を果たす．
\end{example}

\phantomsection
\begin{example}{微分積分学の基本定理}{微分積分学の基本定理}
  この定理は微分と積分の関係を明らかにするものであり，以下のように述べられる：
  \begin{quotation}
    $f$を区間$[a, b]$上で連続な実数値関数とする．このとき，
    \[
      F(x) = \int_a^x f(t) \, dt
    \]
    と定義すると，$F$は$[a, b]$上で連続，$(a, b)$上で微分可能であり，
    \[
      F'(x) = f(x) \qquad (a < x < b)
    \]
    が成り立つ．
  \end{quotation}
  この定理により，積分と微分が互いに逆操作であることが示される．
\end{example}
\phantomsection
\subsection{補題・系}

\emph{補題}\index{ほだい@補題}（lemma）とは，定理や命題を証明するために用いられる補助的な命題のことである．Lemと略記される．

一方，\emph{系}\index{けい@系}（corollary）とは，すでに証明された定理や命題から比較的容易に導かれる結果のことである．Corと略記される．

補題や系を説明するためには，関連するいくつかの命題や定理が必要である．以下に具体的な例を示す．


\begin{example}{（補題）ロルの定理}{（補題）ロルの定理}
  実数値関数$f$ が閉区間 $[a,b]$ で連続，開区間 $(a,b)$ で微分可能，かつ $f(a)=f(b)$ とする．
  このとき，
  \[
  f'(c)=0
  \]
  を満たす点 $c \in (a,b)$ が少なくとも1つ存在する．
\end{example}

\begin{tleftbar}
\begin{proof}
最大値・最小値の定理\footnote{ここでは，最大値・最小値の定理の成立を認めるものとする．}により，$f$は$[a,b]$ で最大値 $M$ と最小値 $m$ をとる．

\begin{enumerate}[(i)]
  \item \textbf{$f$ が定数の場合}：
  任意の $x \in [a,b]$ で $f'(x)=0$ なので，任意の $c \in(a,b)$ で
  $f'(c)=0$ となる．

  \item \textbf{$f$ が定数でない場合}：
  端点では $f(a)=f(b)$ だから，
  $M>f(a)$ または $m<f(a)$ が成り立つ．
  したがって，$f$ の最大値または最小値は内部点
  $c \in(a,b)$ にて達成される．
\end{enumerate}

$c$ で最大をとる場合を示す（最小のときは不等号が逆になる）．
$h>0$ とすると， $f(c+h) \leqq f(c)$だから
\[
\frac{f(c+h)-f(c)}{h} \leqq 0 ,\quad \therefore ~ f'_+(c) \coloneqq \lim_{h \to 0+} \frac{f(c+h)-f(c)}{h}\leqq 0
\]
である．
一方 $h<0$ では同様に
\[
\frac{f(c+h)-f(c)}{h}\geqq 0 , \quad \therefore~  f'_-(c)\coloneqq \lim_{h\to0-}\frac{f(c+h)-f(c)}{h}\geqq 0
\]
である．
$f$の点$c$における微分可能性により，$f'_+(c)=f'_-(c)=f'(c)$であるから，
\[
0\leqq f'(c)\leqq 0 , \quad \therefore ~f'(c)=0
\]
である．
このことから，直ちに主張が従う．
\end{proof}
\end{tleftbar}


\begin{example}{（定理）平均値の定理}{（定理）平均値の定理}
  実数値関数$f$ が $[a,b]$ で連続，$(a,b)$ で微分可能なら，
  ある $c \in(a,b)$ が存在して
  \[
    f'(c)=\frac{f(b)-f(a)}{b-a}
  \]
  が成り立つ．
\end{example}

\begin{tleftbar}
\begin{proof}
  関数$g$を
  \[
    g(x)=f(x)-\frac{f(b)-f(a)}{b-a}(x-a)
  \]
  と定義すると， $g(a)=g(b)$であり，なおかつ$g$ は $[a,b]$ で連続，$(a,b)$ で微分可能なので，
  ロルの定理より $(a,b)$ に $g'(c)=0$ を満たす $c$ がある．
  すると
  \[
    0=g'(c)=f'(c)-\frac{f(b)-f(a)}{b-a}
  \]
  が成り立つ．この式から，直ちに主張が従う．
\end{proof}
\end{tleftbar}

\begin{example}{（系1）$f'(x)=0$ なら $f$ は定数}{（系1）導関数ゼロなら定数}
  実数値関数$f$ が区間 $I$ で微分可能で $f'(x)=0$ （$x \in I$）とする．
  このとき $f$ は $I$ 上で一定である．
\end{example}

\begin{tleftbar}
\begin{proof}
  任意の $x<y$ に平均値の定理を適用すると，ある $c\in(x,y)$ があって
  \[
     \frac{f(y)-f(x)}{y-x}=f'(c)=0
  \]
  が成り立つ．よってこのとき$f(y)=f(x)$である．いま$x,y$ は任意なので $f$ は$I$上で一定である．
\end{proof}
\end{tleftbar}

\begin{example}{（系2）導関数の符号と単調性}{（系2）単調性の系}
  実数値関数$f$ が区間 $I$ で微分可能とする．
  \begin{enumerate}
    \item $f'(x)\geqq 0$（$x\in I$） なら $f$ は $I$ で単調増加関数である．
    \item $f'(x)\leqq 0$（$x\in I$） なら $f$ は $I$ で単調減少関数である．
  \end{enumerate}
\end{example}

\begin{tleftbar}
\begin{proof}
  (1) を示す．$x<y$ をとると平均値の定理により
  \[
     \frac{f(y)-f(x)}{y-x}=f'(c) \quad (c\in(x,y))
  \]
  となる．右辺は $0$以上，かつ $y-x>0$ だから $f(y)-f(x)\geqq 0$である．このことから直ちに主張が従う．
  (2) も同様である．
\end{proof}
\end{tleftbar}

ここまで補題・定理・系という用語を説明してきたが，最後にひとつ注意しておこう．
実は，これらのラベルの使い分けに数学的に厳密な基準はない．
どの主張を「定理」とよび，どの主張を「補題」とよぶかは，
著者がその文章の中で，その主張をどう位置づけたかを表しているにすぎない．
 
実際，本節では平均値の定理を証明するための補助として「ロルの定理」を紹介したが，
世間では「ロルの補題」ではなく「ロルの定理」の名で通っている．
逆に，「補題」の名を冠しながら，多くの定理よりも重要な主張もある\footnote{たとえば，集合論の「ツォルンの補題」は，それぞれの分野で中心的な役割を果たす主張であるが，歴史的な経緯から「補題」とよばれ続けている．}．
また，ある教科書で「定理」とされている主張が，別の教科書では「命題」とされていることも珍しくない．
 
したがって，「補題だから重要でない」「命題だから定理より格下だ」と機械的に判断するのではなく，
その文章の中でその主張が果たしている役割に注目して読むことが大切である．

\begin{comment}
\phantomsection
\begin{example}{（補題）ユークリッドの補題}{（補題）ユークリッドの補題}
  ユークリッドの補題は，整数論における基本的な結果であり，素因数分解の一意性を証明する際に重要な役割を果たす．その主張は以下の通りである：

  \begin{quote}
    素数$p$が整数$a$と$b$の積$ab$を割り切るならば，$p$は$a$と$b$の少なくとも一方を割り切る．
  \end{quote}

\end{example}

\begin{tleftbar}
  \begin{proof}
    素数$p$が$ab$を割り切るとする．もし$p$が$a$を割り切らないならば，$\gcd(p, a) = 1$である．
    このとき，ベズーの等式より，ある整数$s$と$t$が存在して，$sp + ta = 1$が成り立つ．
    両辺に$b$を掛けると，$spb + tab = b$となる．
    左辺の$spb$は$p$で割り切れるが，$tab$は$p$で割り切れるため，
    右辺の$b$も$p$で割り切れる．したがって，$p$は$b$を割り切る．
  \end{proof}
\end{tleftbar}

この補題を用いて，素因数分解の一意性（算術の基本定理）を証明することができる．

\phantomsection
\begin{example}{自然数の約数の個数}{自然数の約数の個数}
  自然数$n$の正の約数の個数$d(n)$は，$n$の素因数分解に基づいて以下の式で与えられる：

  \begin{equation}
    d(n) = (e_1 + 1)(e_2 + 1) \dotsm (e_k + 1)
  \end{equation}

  ただし，$n$を素因数分解して$n = p_1^{e_1} p_2^{e_2} \dotsm p_k^{e_k}$と表した．
\end{example}

\begin{tleftbar}
  \begin{proof}
    各素因数$p_i$について，指数$e_i$は$0$から$e_i$までの値を取り得る．
    したがって，各$p_i$に対する約数の取り得る指数の個数は$e_i + 1$個である．
    すべての素因数について独立に指数を選ぶことで，
    $n$のすべての約数を生成できるため，約数の総数は各$(e_i + 1)$の積となる．
  \end{proof}
\end{tleftbar}

この命題は，素因数分解の一意性から直接的に導かれる結果である．

そしてこの例から，次の系が導かれる：

\phantomsection
\begin{example}{（系）約数の個数が奇数となる条件}{（系）約数の個数が奇数となる条件}
  自然数$n$の約数の個数$d(n)$が奇数となるための必要十分条件は，$n$が完全平方数であることである．
\end{example}

\begin{tleftbar}
  \begin{proof}
    約数の個数$d(n) = (e_1 + 1)(e_2 + 1) \dotsm (e_k + 1)$である．
    各$(e_i + 1)$が奇数となるためには，$e_i$が偶数でなければならない．
    つまり，すべての素因数の指数$e_i$が偶数であるとき，$n$は各素因数の偶数乗の積であり，
    これは$n$が完全平方数であることと同値である．
    逆に，$n$が完全平方数であれば，各$e_i$は偶数であり，したがって$d(n)$は奇数となる．
  \end{proof}
\end{tleftbar}
\end{comment}


\phantomsection
\section{数学でよく使われる言い回し}

数学の文章では，日常語と同じ単語が，日常とは少しずれた意味や役割で使われることが多い．
「かつ・または」のような論理を表す言葉から，「嬉しい」「天下り」のような数学者特有の口ぐせまで，
数学書を読み始めた学生が引っかかりやすい言い回しを，具体例とともにひとつずつ確認していこう．

\phantomsection
\subsection{かつ・または}

数学において，「かつ」は日常とほとんど同じ意味で使われるが，「または」の使い方は日常とは異なる．以下に例を挙げよう．

\phantomsection
\begin{example}{「または」の使用例}{「または」の使用例}
  \begin{enumerate}[(A)]
    \item ランチメニューの主食として，米またはパンがついてくる\footnote{一般的には，この文を「米とパンの両方が食べられる」とは解釈しないであろう．}．\label{enu:米またはパン}
    \item 「運転免許を持っていない人」または「18歳未満の人」はレンタカーを借りることができない． \label{enu:運転免許または18歳未満}
  \end{enumerate}
  \ref{enu:米またはパン}は「どちらか片方のみ」の意味で「または」を使い，\ref{enu:運転免許または18歳未満}は「いずれかが」の意味で「または」を用いている．
\end{example}

数学では，「または」は「いずれかが」の意味で使われ，「どちらか片方のみ」の意味で使われることはない．

たとえば，$A,B$を集合とするとき，
\[
  x \in A \cup B
\]
は「$x$は$A$の元であるか，$B$の元であるか，あるいは両方である」という意味である\footnote{和集合$\cup$の定義は\partref{集合}で改めて述べる．}．
つまり，「$ x \in A$であり，$x \notin  B$である」あるいは「$ x \notin A$であり，$x \in B$である」といった状況のときにも，$x \in A \cup B$とかく．
%ここまで

\subsection{任意の・すべての}

数学において，「任意の\index{にんいの@任意の}」と「すべての\index{すべての@すべての}」はニュアンスが異なる．
英語では，「任意の」は``for any''，「すべての」は``for all''に相当する．
``for any''のニュアンスは「複数の対象があり，その中のどのひとつをとっても」という意味であり，
``for all''は「対象全体を見て，すべてのものが」という意味である．

たとえば，次の文はいずれも自然な日本語である．
\begin{itemize}
  \item 任意の実数$x$について，$x^2\ge0$である．
  \item すべての実数$x$について，$x^2\ge0$である．
\end{itemize}

一方，証明では，
\begin{quote}
任意に$x\in A$をとる．
\end{quote}
のような表現をよく用いる．
これは「集合$A$の中からどの元を選んでもよい」という意味であり，
証明の出発点として自然な表現である．

これに対して，
\begin{quote}
すべての$x\in A$をとる．
\end{quote}
という表現は，日本語として不自然である．
「すべての元」を同時にひとつずつ選ぶことはできないからである．

このように，「任意の」と「すべての」は論理的には同じ意味で用いられることが多いが，
日本語としては役割が異なるため，状況に応じて使い分ける必要がある．
例としてまとめよう：

\begin{example}{「任意の」と「すべての」の使用例}{「任意の」と「すべての」の使用例}
  \begin{enumerate}[(A)]
    \item 任意の実数$x$について，$x^2\ge0$である．
    \item すべての実数$x$について，$x^2\ge0$である．
    \item 任意に$x\in A$をとる．
  \end{enumerate}
\end{example}

\subsection{とる\index{とる}・おく\index{おく}・固定する\index{こていする@固定する}}

数学書の証明を読み始めると，
\begin{quote}
  任意に$x \in A$をとる．
\end{quote}
という表現に必ず出会う．初めて読むと「$x$をどこから取ってくるの？」と，かなり奇妙な日本語に感じられるかもしれない．
これは「集合$A$の元をひとつ選び，それを$x$と名づけて議論の対象にする」という意味であり，「選ぶ」に近い言葉である．
「任意に」がついている場合は，前項で述べたとおり「どの元を選んでも以下の議論が通用する」というニュアンスが込められている．
また，「十分大きな$n$をとる」のように，条件を満たす対象を選ぶ場面でも使われる．

これとよく似た言葉に「おく」がある．
\begin{quote}
  $s=x+y$とおく．
\end{quote}
のように，「おく」は既にある量や式に対して新しい記号を導入する（名前をつける）際に使う言葉である．
「とる」が対象を選ぶ操作であるのに対し，「おく」は記号の定義・命名の操作である点が異なる．

もうひとつ，「固定する」という言葉もある．
\begin{quote}
  $x \in A$をひとつ固定する．
\end{quote}
これは「$x$をひとつ選び，以後の議論ではそれを動かさない」という宣言である．
複数の文字が登場する議論で，「$x$は止めておいて，$y$だけを動かして考える」といった場面で威力を発揮する．
たとえば「$x$を固定するごとに，$y$についての関数$f(x,y)$を考える」のように使われる．

\begin{example}{「とる」「おく」「固定する」の使用例}{「とる」「おく」「固定する」の使用例}
  \begin{enumerate}[(A)]
    \item 任意に$x \in A$をとる．（$A$の元をひとつ選ぶ．どの元でも以下の議論が通用する）
    \item 十分大きな$n \in \mathbb{N}$をとる．（条件を満たす$n$を選ぶ）
    \item $t = \tan \frac{x}{2}$とおく．（新しい記号$t$を導入する）
    \item $\varepsilon > 0$をひとつ固定する．（以後の議論で$\varepsilon$は動かさない）
  \end{enumerate}
\end{example}

この3語はよく使われるので，以下で表として整理しよう：

\begin{table}[htbp]
  \centering
  \begin{tabular}{>{\centering\arraybackslash}m{\mnsc{1.8cm}} L{3.2cm} L{4.2cm} L{4.0cm}}
    \toprule
      & \textbf{操作の内容} & \textbf{典型的な文} & \textbf{心得} \\
    \midrule
    とる & 対象を選ぶ & 任意に$x \in A$をとる． & どの元でも以下の議論が通用する \\
    おく & 記号を名づける & $t = \tan(x/2)$とおく． & 新しい記号の導入（定義） \\
    固定する & 選んで動かさない & $\varepsilon > 0$をひとつ固定する． & 以後の議論でその文字は動かさない \\
    \bottomrule
  \end{tabular}
  \caption{「とる」「おく」「固定する」の使い分け}
  \label{tab:とるおく固定するの比較}
\end{table}

\subsection{成り立つ\index{なりたつ@成り立つ}}

数学において「成り立つ」は，等式・不等式・命題などが真であることを表す言葉である．
日常で「商売が成り立つ」「会話が成り立つ」というときの「うまく機能する」という意味とは少し違い，
数学では「その主張が正しい」ということを端的に表す．

\begin{example}{「成り立つ」の使用例}{「成り立つ」の使用例}
  \begin{enumerate}[(A)]
    \item 任意の実数$x$に対して，不等式$x^2 \geqq 0$が成り立つ．
    \item $n$が偶数のとき，$n^2$も偶数である，という命題が成り立つ．
  \end{enumerate}
\end{example}

主語になるのは「等式」「不等式」「命題」といった主張そのものであることに注意しよう．
本書でもすでに「$a^2+b^2=c^2$が成り立つ」のような形で繰り返し使ってきたが，
すべて「この式は真である」という意味である．
\enlargethispage{2mm}

\subsection{一般に\index{いっぱんに@一般に}}

「成り立つ」とセットでよく使われる言葉に「一般に」がある．
これは，肯定文と否定文とで読み方の注意点が異なる，少し曲者の言い回しである．
\newpage 

まず，肯定文とともに使われる場合を見よう．
\begin{quote}
  一般に，実数$x$に対して$x^2 \geqq 0$が成り立つ．
\end{quote}
このときの「一般に」は「すべての場合に」という意味であり，
「任意の実数$x$に対して成り立つ」という全称の主張と同じことを述べている．

要注意なのは，否定文とともに使われる場合である．
\begin{quote}
  $\sqrt{a^2} = a$は一般に成り立たない．
\end{quote}
これは「つねに成り立つとは限らない」，すなわち「成り立たない場合が少なくともひとつ存在する」という意味である．
決して「すべての場合に成り立たない」という意味ではない．
実際，$\sqrt{a^2}=a$は$a \geqq 0$のときには成り立ち，$a < 0$のときに成り立たない
（\exref{「絶対値」の定義}で見たとおり，正しくは$\sqrt{a^2}=\abs{a}$である）．
このように，成り立つ場合がいくらあっても，例外がひとつでもあれば「一般に成り立たない」というのである．
「一般には成り立たない」と，「は」を入れて書かれることもあるが，意味は同じである．

\begin{example}{「一般に」の使用例}{「一般に」の使用例}
  \begin{enumerate}[(A)]
    \item 一般に，$(a+b)^2 = a^2 + 2ab + b^2$が成り立つ．（任意の実数$a,b$で成り立つ）
    \item 正方行列$A,B$に対して，$AB=BA$は一般に成り立たない．（成り立つ組もあるが，成り立たない組が存在する）\label{enu:積の非可換}
  \end{enumerate}
\end{example}

「一般に成り立たない」ことを証明するには，成り立たない例，すなわち\emph{反例}\index{はんれい@反例}（counterexample）をひとつ挙げれば十分である．
たとえば\ref{enu:積の非可換}については，
\[
  A=\begin{pmatrix} 1 & 1 \\ 0 & 1 \end{pmatrix},\quad
  B=\begin{pmatrix} 1 & 0 \\ 1 & 1 \end{pmatrix}
\]
とすれば，$AB \ne BA$となることが計算により確かめられる\footnote{一方，$B=E$とすれば$AB=BA$である．「成り立つ例がある」ことと「一般に成り立つ」ことは別物である．}．
全称の主張（一般に成り立つ）を示すにはすべての場合を扱う議論が必要である一方，
その否定（一般に成り立たない）は反例ひとつで示せる．この非対称性は，数学の議論の基本として覚えておきたい．
なお，のちに述べる「一般性を失わず」の「一般性」も，この「すべての場合」という意味での「一般」を指している．

\subsection{示す\index{しめす@示す}}

数学において「示す」は「証明する」とほぼ同じ意味で使われる．
演習問題などで頻出の「〜を示せ」は「〜を証明せよ」という意味である．

また，証明の中では
\begin{quote}
  (1)を示す．
\end{quote}
のように，これから証明する目標を宣言するためにも使われる．
長い証明では「まず$\subseteq$を示す」「次に$\supseteq$を示す」のように，
証明全体をいくつかの目標に分割する際の目印となる．
証明を読むときは，「示す」と宣言された目標がどこで達成されたかを意識すると，議論の構造が見えやすくなる．

\begin{example}{「示す」の使用例}{「示す」の使用例}
  \begin{enumerate}[(A)]
    \item $\sqrt{2}$が無理数であることを示せ．（＝証明せよ）
    \item まず，数列$(x_n)_{n \in \mathbb{N}}$が上に有界であることを示す．（＝これから証明する）
  \end{enumerate}
\end{example}

\subsection{仮定する\index{かていする@仮定する}}

数学において「仮定する」は，ある主張を一時的に正しいものとみなして議論を進めることを表す．
真であると断定しているわけではないことに注意しよう．

「仮定する」が活躍する典型的な場面はふたつある．
ひとつは，数学的帰納法のように「$n=k$のとき成り立つと仮定する」と述べて，そこから$n=k+1$の場合を導く場面である．
もうひとつは背理法であり，「$\sqrt{2}$が有理数であると仮定する」のように，
最終的には否定したい主張をあえて一時的に認め，矛盾を導くために使われる．

\begin{example}{「仮定する」の使用例}{「仮定する」の使用例}
  \begin{enumerate}[(A)]
    \item $n=k$のとき主張が成り立つと仮定する．（数学的帰納法）
    \item $\sqrt{2}$が有理数であると仮定する．（背理法：のちに矛盾を導く）
  \end{enumerate}
\end{example}

いずれの場合も，「仮定」はその議論の中だけで有効な約束であり，議論が終われば効力を失う．
「仮定する」と書かれていたら，「この仮定はどの範囲で使われ，どこで役目を終えるのか」を意識して読むとよい．

\subsection{十分大きい\index{じゅうぶんおおきい@十分大きい}}

日常語の「十分大きい」と，数学における「十分大きい」には，かなりの差がある．
数学書で
\begin{quote}
  十分大きな$n$に対して，$a_n < \varepsilon$が成り立つ．
\end{quote}
と書かれていたら，これは「$n$がなんとなく大きければ」という曖昧な意味ではない．正確には
\begin{quote}
  ある$N \in \mathbb{N}$が存在して，$n \geqq N$を満たすすべての$n$について$a_n < \varepsilon$が成り立つ．
\end{quote}
という意味である．
つまり「十分大きい」は，「ある値（しきい値）以上ならすべて」という，存在と全称を組み合わせた厳密な主張の省略表現なのである．

\begin{figure}[htbp]
  \centering
  \begin{tikzpicture}[x=1cm]
    % n >= N の領域（帯）
    \fill[figaccent!10] (5,-0.3) rectangle (9.4,0.3);
    % 数直線
    \draw[thick, appleGrayDark, -{Stealth[length=2mm]}]
      (0,0) -- (9.8,0) node[below=2pt, font=\scriptsize, text=appleInk] {$n$};
    % N 未満：成り立たない n が混ざっていてもよい（有限個）
    \foreach \x in {0.7,1.5,2.3,3.2,4.1}
      \fill[appleGray] (\x,0) circle (1.6pt);
    % N 以上：すべて成り立つ
    \foreach \x in {5.5,6.2,6.9,7.6,8.3,9.0}
      \fill[figaccent] (\x,0) circle (1.6pt);
    % しきい値 N
    \draw[thick, appleGrayDark] (5,-0.3) -- (5,0.3)
      node[above=1pt, font=\scriptsize, text=appleInk] {しきい値$N$};
    % 注釈
    \node[flowtxt, align=center] at (2.5,-0.75)
      {$n < N$：成り立たない$n$が\\あってもよい（有限個）};
    \node[font=\scriptsize, text=figaccent, align=center] at (7.2,-0.75)
      {$n \geqq N$：すべての$n$で\\$a_n < \varepsilon$が成り立つ};
  \end{tikzpicture}
  \caption{「十分大きな$n$に対して成り立つ」のイメージ}
  \label{fig:十分大きいのイメージ}
\end{figure}

\begin{example}{「十分大きい」の使用例}{「十分大きい」の使用例}
  \begin{enumerate}[(A)]
    \item 十分大きな$n$に対して，$\dfrac{1}{n} < 0.01$が成り立つ．
    （実際，$N=101$とすれば，$n \geqq N$なるすべての$n$で成り立つ）
    \item 十分大きな$x$に対して，$x^2 > 100x$が成り立つ．
    （$x > 100$なるすべての$x$で成り立つ）
  \end{enumerate}
\end{example}

この言い回しは，のちに学ぶ数列の極限（$\varepsilon$-$N$論法）に直結する．
「十分大きい」を見たら「しきい値がどこかに存在する」と読み替える習慣をつけておくと，極限の定義がぐっと理解しやすくなる．

\subsection{簡単のため\index{かんたんのため@簡単のため}}

「簡単にするために」のほうがしっくりくる方もおられるかもしれないが，これは``For simplicity''の訳で，
その後の議論を簡単にするために「議論することが楽な仮定」を設ける際に使う言葉である．

\begin{example}{「簡単のため」の使用例}{「簡単のため」の使用例}
  \begin{enumerate}[(A)]
    \item 簡単のため，平行移動して頂点が原点にある形として，$f(x)=ax^2$を考える．
  \end{enumerate}
\end{example}

\subsection{一般性を失わず\index{いっぱんせいをうしなわず@一般性を失わず}}

数学書には
\begin{quote}
  一般性を失わず，$a \leqq b$としてよい．
\end{quote}
という表現がしばしば登場する．初めて読むと「勝手に$a \leqq b$と仮定していいの？」と戸惑うかもしれない．
これは，
\begin{quote}
  残りの場合（ここでは$a > b$の場合）も，文字を入れ替えるなどして同じ議論で処理できるので，一方の場合だけ考えても問題がない
\end{quote}
という意味である．英語では``without loss of generality''といい，WLOGと略記されることもある．

\begin{example}{「一般性を失わず」の使用例}{「一般性を失わず」の使用例}
  実数$a,b$に対して，$\max\{a,b\} + \min\{a,b\} = a+b$が成り立つことを示したい．
  このとき，
  \begin{quote}
    一般性を失わず，$a \leqq b$としてよい．すると$\max\{a,b\}=b$，$\min\{a,b\}=a$なので，
    \[
      \max\{a,b\} + \min\{a,b\} = b + a = a + b
    \]
    が成り立つ．
  \end{quote}
  と議論できる．$a > b$の場合は，$a$と$b$の役割を入れ替えれば同じ議論がそのまま通用するので，
  改めて書く必要がないのである（示したい等式が$a$と$b$について対称であることが効いている）．
\end{example}

「一般性を失わず」は，数学特有の省略の仕方のひとつである．
ただし，本当に残りの場合が同じ議論で処理できるのかは，読者・書き手の双方が確認すべきことである．
対称性がないのに「一般性を失わず」と書いてしまうと，証明に穴があくので注意しよう．

\subsection{同様に\index{どうように@同様に}}

「同様に」も，数学特有の省略の表現である．
本当に同じ論法がそのまま使える部分について，繰り返しを避けるために「同様に」と述べて詳細を省略する．

実際，本書でも「導関数の符号と単調性」の系の証明において，(1)を証明したあとに
\begin{quote}
  (2)も同様である．
\end{quote}
と述べて証明を終えた．これは「(1)の議論の不等号の向きを逆にすれば，そのまま(2)の証明になる」という意味である．

\begin{example}{「同様に」の使用例}{「同様に」の使用例}
  \begin{enumerate}[(A)]
    \item $c$で最大をとる場合を示す（最小のときは不等号が逆になるだけで，同様である）．
    \item $x \geqq 0$の場合は上で示した．$x<0$の場合も同様に示される．
  \end{enumerate}
\end{example}

「同様に」と書かれた箇所は，著者が「読者が自力で再現できる」と判断して省略した部分である．
数学書を読むときは，「同様に」の中身を実際に自分の手でなぞるのもよい．

\subsection{従う\index{したがう@従う}}

数学において「従う」は日常とはまた違った意味で使われる．日常だと「王に従う」や「上司に従う」など，「命令を受けてそれに従う」という意味で使われることが多いが，
数学では「$A$という事柄から，すぐ$B$という事柄がわかる」という場合に「$A$から$B$が従う」という．

\begin{example}{「従う」の使用例}{「従う」の使用例}
  \begin{enumerate}[(A)]
    \item $n$が偶数であることから，$n^2$が偶数であることが従う．\label{enu:偶数ならば偶数}
    \item $G$が群であることから，$G \ne \varnothing$が従う\footnote{群であれば単位元が存在することは群の定義からわかる．そのため群は空集合でない．}．\label{enu:群の例}
  \end{enumerate}
\end{example}
\enlargethispage{4mm}

\subsection{明らか\index{あきらか@明らか}・直ちに\index{ただちに@直ちに}}

数学書には「明らかである」「明らかに〜が成り立つ」という表現が頻出する．
文字どおり受け取ると「証明不要なほど自明」という意味に思えるが，実際には以下のような意味である：
\newpage 

\begin{quote}
  ここでは詳細を省略できる程度の主張である（本筋から外れるので，証明は読者に委ねる）
\end{quote}
つまり「明らか」は主張の真偽についての言葉というより，著者の省略の宣言なのである．
そのため，読者にとってはまったく明らかでないこともしばしばある\footnote{「明らか」と書かれた箇所でつまずくのは恥ずかしいことではない．むしろ，行間を自力で埋める良い練習問題だと考えるとよい．}．

一方，「直ちに」は，既知の事実からごく短い推論で結論が得られることを表す．
実際，本書でもロルの定理の証明などで
\begin{quote}
  このことから，直ちに主張が従う．
\end{quote}
という表現を用いた．これは「得られた式から，ほんの一，二歩の変形・言い換えで結論に到達する」という意味である．

\begin{example}{「明らか」「直ちに」の使用例}{「明らか」「直ちに」の使用例}
  \begin{enumerate}[(A)]
    \item $x > 0$は明らかである．（詳細は省略するが，各自で確認できる程度の主張である）
    \item $f'(c)=0$から，直ちに主張が従う．（短い推論で結論が得られる）
  \end{enumerate}
\end{example}

「明らか」「直ちに」に出会ったら，鵜呑みにせず，省略された推論を自分の手で復元してみることを勧める．
それが「本当に明らか」であれば数行で終わるし，そうでなければ，そこが理解の急所である．

\enlargethispage{4mm}
\subsection{〜によらない\index{によらない@〜によらない}}

数学において「〜によらない」は，選び方・表し方・添字などを変えても結果が変わらないことを表す．

\begin{example}{「〜によらない」の使用例}{「〜によらない」の使用例}
  \begin{enumerate}[(A)]
    \item 三角形の内角の和は，三角形の選び方によらず$180^\circ$である．
    \item 極限値$\lim_{n \to \infty} a_n$は，最初の有限項の値によらない．
  \end{enumerate}
\end{example}

この言葉は，本章の「定義」の項で述べたwell-definedの概念と深く関係している．
「$g\left(\frac{a}{b}\right)=\frac{2a}{b}$という定義は分数の表し方によらず値を定める」からこそ，$g$はwell-definedなのであった．
数学では，なにかを定義したときに「その定義が途中の選択によらないこと」の確認が要求される場面が多い．
「〜によらない」という表現を見たら，「なにを取り替えても結果が同じなのか」を意識して読むとよい．


\subsection{〜に関して\index{にかんして@〜に関して}}

数学において「〜に関して」は，どの文字（変数）に注目して操作・議論を行うのかを指定する言葉である．
\newpage 
\begin{example}{「〜に関して」の使用例}{「〜に関して」の使用例}
  \begin{enumerate}[(A)]
    \item 関数$f(x,y)=x^2 y$を$x$に関して微分すると，$2xy$である．
    （$y$は定数とみなし，$x$を変数として微分する）
    \item $f(x)=x^2$は，$x \geqq 0$の範囲で$x$に関して単調増加である．
    \item 等式$x+y=y+x$は，$x$と$y$に関して対称である．
  \end{enumerate}
\end{example}

複数の文字が登場する式では，「どの文字を動かし，どの文字を止めているのか」が曖昧だと議論が混乱する．
「$x$に関して」という指定は，先に述べた「固定する」と対になる表現であり，
「$y$は固定して，$x$に関して微分する」のように組み合わせて使われることも多い．

\subsection{嬉しい\index{うれしい@嬉しい}}

日常では，「嬉しい」という言葉は「喜びを感じる」という意味で使われることが多いが，数学の文脈では「都合がよい」，「議論が楽になる」という意味でたびたび使われる．

\begin{example}{「嬉しい」の使用例}{「嬉しい」の使用例}
  \[
    \int_{\alpha}^{\beta} (x-\alpha)(x-\beta)\, dx = -\frac{(\beta-\alpha)^3}{6}
  \]
  であることを示したい．このときに
  \[
    \int_{\alpha}^{\beta} (x-\alpha)(x-\beta)\, dx = \int_{\alpha}^{\beta}  (x-\alpha)\{(x-\alpha) - (\beta-\alpha)\} \, dx
  \]
  と変形してなにが嬉しいかというと，
  \begin{align*}
    \int_{\alpha}^{\beta}  (x-\alpha)\{(x-\alpha) - (\beta-\alpha)\} \, dx & = \int_{\alpha}^{\beta} \{  (x-\alpha)^2 -(\beta-\alpha)(x-\alpha) \} \, dx                        \\
                                                                           & = \Biggl [ \frac{(x-\alpha)^3}{3} - \frac{(\beta-\alpha)(x-\alpha)^2}{2} \Biggr ]_{\alpha}^{\beta} \\
                                                                           & = - \frac{1}{6} (\beta-\alpha)^3
  \end{align*}
  というふうに，簡単に計算できるからである．
\end{example}



\subsection{おさえる\index{おさえる}・評価\index{ひょうか@評価}する}

\begin{wrapfigure}[7]{r}{0.57\textwidth}
  \centering
  \begin{tikzpicture}[x=1.62cm, y=1.1cm]
    % 数直線
    \draw[thick, appleGrayDark, -{Stealth[length=2.2mm]}] (0.4,0) -- (4.8,0);
    \foreach \x in {1,2,3,4} {
      \draw[appleGrayDark] (\x,-0.07) -- (\x,0.07);
      \node[font=\scriptsize, text=appleInk, below=2.2pt] at (\x,0) {$\x$};
    }
    % 粗い評価 1 < M < 4
    \draw[line width=2.2pt, appleGray] (1,0.5) -- (4,0.5);
    \draw[appleGray] (1,0.38) -- (1,0.62);
    \draw[appleGray] (4,0.38) -- (4,0.62);
    \node[flowtxt, right=3.3pt] at (4.05,0.5) {$1 < M < 4$};

    % 厳しい評価 2 < M < 3
    \draw[line width=2.2pt, figaccent] (2,1.05) -- (3,1.05);
    \draw[figaccent] (2,0.93) -- (2,1.17);
    \draw[figaccent] (3,0.93) -- (3,1.17);
    \node[font=\scriptsize, text=figaccent, right=3.3pt] at (4.05,1.05)
      {$2 < M < 3$（評価が厳しい）};
  \end{tikzpicture}
  \caption{$M$のありかを絞り込むほど「評価が厳しい」}
  \label{fig:評価が厳しい}
\end{wrapfigure}
数学の文脈では「おさえる」という言葉は，ある数$K$を，別の数 $M$を用いて，$K<M$という形で表すことを指す．
より詳しく「上からおさえる」と表すこともある．

また，数学の文脈で「評価」という言葉は，ある数$M$を不等式で表し，その数がどの程度の大きさであるかを示すことを指す．
たとえば，$ 1 < M <4$と表すことは「$M$を$1$より大きく$4$より小さいと評価する」ということを指す．
さらに，上記の$M$を$ 2< M <3$と表すことができる場合は，後者の表現のほうが「評価が厳しい」ということができる．




\begin{example}{「おさえる」と「評価する」の使用例}{「おさえる」と「評価する」の使用例}
  \begin{enumerate}[(A)]
    \item 数列$(a_n)_{n \in \mathbb{N}}$を $ a_n <M$と上からおさえる．
    \item $e$の値を$2.7 < e < 2.8$と評価する．
  \end{enumerate}
\end{example}


\subsection{特徴づけ\index{とくちょうづけ@特徴づけ}}

数学の文脈では，「特徴づけ」（characterization）という言葉がよく使われる．
これは，ある対象（あるいはあるクラスに属すること）と同値な条件を与えることを指す．
「$X$であることと，条件$P$を満たすことは同値である」という形の主張が得られたとき，「条件$P$は$X$を特徴づける」というのである．

\begin{example}{「特徴づけ」の使用例}{「特徴づけ」の使用例}
  \begin{enumerate}[(A)]
    \item 偶数は，「ある整数$k$を用いて$n=2k$と表せる整数$n$」として特徴づけられる．
    （$n$が偶数であることと，$n=2k$と表せることは同値である）
    \item 正則行列は，「行列式が$0$でない正方行列」として特徴づけられる\footnote{これは線型代数で学ぶ重要な定理である．定義（逆行列を持つこと）とは別の，判定しやすい同値条件が得られている点に価値がある．}．
    \item 二等辺三角形は，「ふたつの内角が等しい三角形」として特徴づけられる．
  \end{enumerate}
\end{example}

特徴づけの価値は，定義とは別の切り口から同じ対象を捉え直せることにある．
たとえば偶数の定義を「$2$で割り切れる整数」としたとき，「$n=2k$と表せる」という特徴づけは，
偶数に関する等式の計算を進めるうえで扱いやすい形を与えてくれる．
定義・特徴づけのどちらを使って議論してもよいが，そのためには両者が同値であることの証明が必要である
（これは本章の冒頭で述べた，絶対値のふたつの定義の同値性の話と同じ構図である）．

\subsection{天下り\index{あまくだり@天下り}}

数学の文脈で，数式や概念を根拠を明示せずに導入することを「天下り」とよぶことがある．

\begin{example}{「天下り」の使用例}{「天下り」の使用例}
  \begin{enumerate}[(A)]
    \item $a,b$を$a<b$を満たす実数とするとき，$ a < c < b$となる$c$の存在を示すために，
    天下りではあるが，$c=(a+b)/2$と定める\label{enu:天下りの例}．
  \end{enumerate}
\end{example}

\ref{enu:天下りの例}においては，根拠を明示せずに$c=(a+b)/2$としたが，これは「$c$は$a$と$b$の相加平均である」という根拠に基づいて定めている．


\begin{table}[htbp]
  \centering
  \begin{tabular}{L{3.2cm} L{8.6cm}}
    \toprule
    \textbf{言い回し} & \textbf{ひとことで言うと} \\
    \midrule
    または & 両方であってもよい（「どちらか片方のみ」ではない） \\
    任意の・すべての & どのひとつをとっても／対象全体が \\
    とる・おく・固定する & 選ぶ／名づける／動かさないと宣言する \\
    成り立つ & その主張（等式・不等式・命題）は真である \\
    一般に & すべての場合に．否定文では「例外が存在する」 \\
    示す & 証明する．証明目標の宣言にも使う \\
    仮定する & 一時的に真とみなす（有効範囲つき） \\
    十分大きい & あるしきい値以上ならすべて \\
    簡単のため & 議論が楽になる仮定を設ける \\
    一般性を失わず & 残りの場合は同じ議論で処理できる \\
    同様に & 同じ論法の繰り返しを省略する \\
    従う & $A$から$B$がすぐわかる \\
    明らか・直ちに & 省略の宣言／ごく短い推論で結論に至る \\
    〜によらない & なにを取り替えても結果が変わらない \\
    〜に関して & どの文字に注目して操作するかの指定 \\
    嬉しい & 都合がよい・議論が楽になる \\
    おさえる・評価する & 不等式で大きさを見積もる \\
    特徴づけ & 定義と同値な別の条件を与える \\
    天下り & 根拠を明示せずに（後回しにして）導入する \\
    \bottomrule
  \end{tabular}
  \caption{数学でよく使われる言い回しの早見表}
  \label{tab:言い回し早見表}
\end{table}

\phantomsection
\part{論理}

\begin{abstract}
この章では，数学の議論を支える「論理」を整理する．まず命題論理の記号と真理値表を確認し，交換律・分配律・ド・モルガンの法則などの基本法則を述べる．含意の真偽をめぐる注意と，逆・裏・対偶の関係もここで確認する．続いて，必要条件・十分条件という言い回しと，証明を支える推論規則を恒真式として整理する．次に，変数を含む文を扱う述語論理を導入し，全称命題・存在命題，自由変数と束縛変数，および量化子の入れ子について注意しよう．最後に，量化子の否定を用いて「数列が収束しないこと」を書き下す．
\end{abstract}

\phantomsection
\section{命題論理\index{めいだいろんり@命題論理}}

命題論理では，命題を結合するために記号を用いる．代表的な論理記号を列挙してみよう：

\begin{table}[ht]
  \centering
  \caption{命題論理の記号}
  \begin{tabular}{c|c}
    \hline
    記号      & 意味                                     \\
    \hline
    $\land$   & かつ（論理積\index{ろんりせき@論理積}）   \\
    $\lor$    & または（論理和\index{ろんりわ@論理和}）  \\
    $\lnot$   & 否定\index{ひてい@否定}                  \\
    $\limp$   & ならば（含意\index{がんい@含意}）        \\
    $\lequiv$ & 同値\index{どうち@同値}                  \\
    \hline
  \end{tabular}
\end{table}

なお，含意を$\to$，同値を$\leftrightarrow$とかく文献も多い．本書でこれらを用いないのは，$\to$をのちに現れる写像$f \colon A \to B$や極限$\lim_{n \to \infty} a_n$の矢印に予約しておくためである．ひとつの記号にひとつの意味だけを担わせておくと，読み間違いが起こりにくい．

\begin{remark}[記号を使う場面と，日本語で書く場面]
  論理記号は，論理式そのものを組み立てて論じるための道具である．そこで本書では，これらを使う場面をつぎのように定めておく．
  \begin{itemize}
    \item 定義や命題の主張，証明の地の文は，原則として日本語で書く．すなわち「かつ」「または」「〜ならば」「任意の$x$に対して」「〜が存在して」と述べる．
    \item 論理式そのものを数式として書き下すときは，論理記号だけを用いる．このとき，数式のなかに日本語の接続詞を混ぜない．
  \end{itemize}
  たとえば「$A \subset B$かつ$B \subset A$」は日本語の文のなかでの言い方であり，同じ内容を論理式として書くときは$(A \subset B) \land (B \subset A)$とする．どちらの書き方も正しいが，ひとつの数式のなかで両者を混ぜると，同じ「かつ」が2通りの姿で現れて読みにくくなる．
\end{remark}

これらの論理記号を用いて，命題を結合し，複雑な論理式を構成することができる．
以下に，基本的な命題の真理値表を載せよう．

\begin{table}[htbp]
\centering
\small

%================================================
% 上段：3段の表
%================================================
\begin{minipage}[t]{0.45\linewidth}
\centering
\caption{否定命題の真理値表}
\label{tab:否定命題の真理値表}
\begin{tabular}{|c|c|} \hline
  $P$          & $\lnot P$    \\ \hline
  $\mathrm{T}$ & $\mathrm{F}$ \\ \hline
  $\mathrm{F}$ & $\mathrm{T}$ \\ \hline
\end{tabular}
\end{minipage}
\hfill
\begin{minipage}[t]{0.45\linewidth}
\centering
\caption{恒真命題と矛盾命題の真理値表}
\label{tab:恒真命題と矛盾命題の真理値表}
\begin{tabular}{|c|c|c|c|} \hline
  $P$          & $\lnot P$    & $P \lor \lnot P$ & $P \land \lnot P$ \\ \hline
  $\mathrm{T}$ & $\mathrm{F}$ & $\mathrm{T}$     & $\mathrm{F}$      \\ \hline
  $\mathrm{F}$ & $\mathrm{T}$ & $\mathrm{T}$     & $\mathrm{F}$      \\ \hline
\end{tabular}
\end{minipage}

\vspace{6mm}

%================================================
% 下段：5段の表
%================================================
\begin{minipage}[t]{0.32\linewidth}
\centering
\caption{論理和（積）の真理値表}
\label{tab:論理和（積）の真理値表}
\begin{tabular}{|c|c|c|c|} \hline
  $P$          & $Q$          & $P \land Q$  & $P \lor Q$   \\ \hline
  $\mathrm{T}$ & $\mathrm{T}$ & $\mathrm{T}$ & $\mathrm{T}$ \\ \hline
  $\mathrm{T}$ & $\mathrm{F}$ & $\mathrm{F}$ & $\mathrm{T}$ \\ \hline
  $\mathrm{F}$ & $\mathrm{T}$ & $\mathrm{F}$ & $\mathrm{T}$ \\ \hline
  $\mathrm{F}$ & $\mathrm{F}$ & $\mathrm{F}$ & $\mathrm{F}$ \\ \hline
\end{tabular}
\end{minipage}
\hfill
\begin{minipage}[t]{0.32\linewidth}
\centering
\caption{「含意」の真理値表}
\label{tab:含意の真理値表}
\begin{tabular}{|c|c|c|} \hline
  $P$          & $Q$          & $P \limp Q$  \\ \hline
  $\mathrm{T}$ & $\mathrm{T}$ & $\mathrm{T}$ \\ \hline
  $\mathrm{T}$ & $\mathrm{F}$ & $\mathrm{F}$ \\ \hline
  $\mathrm{F}$ & $\mathrm{T}$ & $\mathrm{T}$ \\ \hline
  $\mathrm{F}$ & $\mathrm{F}$ & $\mathrm{T}$ \\ \hline
\end{tabular}
\end{minipage}
\hfill
\begin{minipage}[t]{0.32\linewidth}
\centering
\caption{「同値」の真理値表}
\label{tab:同値の真理値表}
\begin{tabular}{|c|c|c|} \hline
  $P$          & $Q$          & $P \lequiv Q$ \\ \hline
  $\mathrm{T}$ & $\mathrm{T}$ & $\mathrm{T}$          \\ \hline
  $\mathrm{T}$ & $\mathrm{F}$ & $\mathrm{F}$          \\ \hline
  $\mathrm{F}$ & $\mathrm{T}$ & $\mathrm{F}$          \\ \hline
  $\mathrm{F}$ & $\mathrm{F}$ & $\mathrm{T}$          \\ \hline
\end{tabular}
\end{minipage}

\end{table}

上記の真理値表によって，命題論理の記号を定義する．

以下，ふたつの論理式$P$，$Q$について，それらに含まれる命題（命題変数）の真偽をどのように定めても$P$と$Q$の真理値がつねに一致するとき，$P$と$Q$は\emph{論理的に同値}\index{ろんりてきどうち@論理的同値}であるといい，$P \equiv Q$と表す\footnote{$P \lequiv Q$がそれ自体ひとつの論理式であるのに対し，$P \equiv Q$は「$P \lequiv Q$がつねに真である」という，論理式についての主張である点で，両者は区別される．}．

たとえば，のちに示す$P \limp Q \equiv \lnot P \lor Q$は，「$P$，$Q$の真偽の組合せ$4$通りのすべてにおいて，$P \limp Q$と$\lnot P \lor Q$の真理値が一致する」ことを主張している．ある特定の場合にたまたま真偽が一致するだけでは，論理的に同値であるとはいわないことに注意しよう．

\begin{prop}{}{}
  命題$P$，$Q$，$R$に対して，次の式が成り立つ：
  \begin{description}
    \item[交換律]
          \[
            P \land Q \equiv Q \land P,\quad
            P \lor Q \equiv Q \lor P
          \]
    \item[結合律]
          \[
            (P \land Q) \land R \equiv P \land (Q \land R),\quad
            (P \lor Q) \lor R \equiv P \lor (Q \lor R)
          \]
    \item[分配律]
          \[
            P \land (Q \lor R) \equiv (P \land Q) \lor (P \land R),\quad
            P \lor (Q \land R) \equiv (P \lor Q) \land (P \lor R)
          \]
    \item[ド・モルガンの法則]
          \[
            \lnot (P \land Q) \equiv \lnot P \lor \lnot Q,\quad
            \lnot (P \lor Q) \equiv \lnot P \land \lnot Q
          \]
  \end{description}
\end{prop}

\begin{prop}{}{}
  命題$P$，$Q$に対して，次の式が成り立つ：
  \begin{description}
    \item[含意の等価表現] 
          \[
            P \limp Q \equiv \lnot P \lor Q
          \]
    \item[同値の等価表現]
          \[
            P \lequiv Q \equiv (P \limp Q) \land (Q \limp P)
          \]
  \end{description}
\end{prop}

これらはいずれも，両辺の真理値表を書き，命題$P$，$Q$（，$R$）の真偽の組合せをすべて比べることで確かめられる．例として，ド・モルガンの法則の前半を確かめておこう．

\begin{tleftbar}
  \begin{proof}[ド・モルガンの法則の前半の証明]
    $P$，$Q$の真偽の組合せ$4$通りすべてについて，両辺の真偽を列挙する：
    \begin{center}
      \begin{tabular}{|c|c||c|c||c|c|c|} \hline
        $P$          & $Q$          & $P \land Q$  & $\lnot (P \land Q)$ & $\lnot P$    & $\lnot Q$    & $\lnot P \lor \lnot Q$ \\ \hline
        $\mathrm{T}$ & $\mathrm{T}$ & $\mathrm{T}$ & $\mathrm{F}$        & $\mathrm{F}$ & $\mathrm{F}$ & $\mathrm{F}$           \\ \hline
        $\mathrm{T}$ & $\mathrm{F}$ & $\mathrm{F}$ & $\mathrm{T}$        & $\mathrm{F}$ & $\mathrm{T}$ & $\mathrm{T}$           \\ \hline
        $\mathrm{F}$ & $\mathrm{T}$ & $\mathrm{F}$ & $\mathrm{T}$        & $\mathrm{T}$ & $\mathrm{F}$ & $\mathrm{T}$           \\ \hline
        $\mathrm{F}$ & $\mathrm{F}$ & $\mathrm{F}$ & $\mathrm{T}$        & $\mathrm{T}$ & $\mathrm{T}$ & $\mathrm{T}$           \\ \hline
      \end{tabular}
    \end{center}
    表より，$\lnot (P \land Q)$の列と$\lnot P \lor \lnot Q$の列は$4$通りすべての場合で一致している．よって$\lnot (P \land Q) \equiv \lnot P \lor \lnot Q$である．これが証明すべきことであった．
  \end{proof}
\end{tleftbar}

また，否定を$2$回とると元の命題に戻る．当たり前に見えることではあるが，のちの証明でたびたび用いるので，命題として述べておこう．

\phantomsection
\begin{prop}{二重否定}{二重否定}
  命題$P$に対して，
  \[
    \lnot \lnot P \equiv P
  \]
  である．
\end{prop}

\begin{tleftbar}
  \begin{proof}
    表\ref{tab:否定命題の真理値表}を$2$回用いると，$P$が真のとき$\lnot P$は偽，したがって$\lnot \lnot P$は真であり，$P$が偽のとき$\lnot P$は真，したがって$\lnot \lnot P$は偽である．よって$\lnot \lnot P$と$P$の真偽はつねに一致する．これが証明すべきことであった．
  \end{proof}
\end{tleftbar}

\phantomsection
\subsection{含意の真偽と逆・裏・対偶}

含意$P \limp Q$は，慣れないうちは真偽の判定にとまどうことが多い．そこで，表\ref{tab:含意の真理値表}にしたがって，$P$と$Q$の真偽の組合せ$4$通りをすべて列挙して確認しておこう．

\begin{center}
  \begin{tabular}{|c|c|c|l|} \hline
    $P$ & $Q$ & $P \limp Q$ & 備考                                   \\ \hline
    真  & 真  & 真          & 約束は守られている                     \\ \hline
    真  & 偽  & 偽 & 偽となるのは，この場合だけ    \\ \hline
    偽  & 真  & 真          & $P$が成り立たないので，約束は破られない \\ \hline
    偽  & 偽  & 真          & 同上                                   \\ \hline
  \end{tabular}
\end{center}

とくに注意したいのが3行目，すなわち「$P$でない，かつ$Q$である」（$\lnot P \land Q$）という場合である．このときにも$P \limp Q$は真とされる．
これは，含意$P \limp Q$が「$P$が成り立つ場合に，$Q$も成り立つ」ということのみを主張しており，$P$が成り立たない場合についてはなにも主張していないからである．
「約束が破られるのは，$P$が成り立ったのに$Q$が成り立たなかったときだけ」と考えるとよい\footnote{このように，仮定$P$が偽であるために$P \limp Q$が無条件に真となることを，空虚な真（vacuous truth）ということがある．}．

\phantomsection
\begin{example}{含意の真偽}{含意の真偽}
  実数$x$に関する命題「$x > 2$ならば$x > 0$」は真である．このことを，いくつかの$x$の値で確認してみよう．
  \begin{center}
    \begin{tabular}{|c|c|c|l|} \hline
      $x$  & $x > 2$（$P$） & $x > 0$（$Q$） & 約束は守られているか                     \\ \hline
      $3$  & 真             & 真             & 守られている                             \\ \hline
      $1$  & 偽             & 真             & 仮定が成り立たないので，破られてはいない \\ \hline
      $-1$ & 偽             & 偽             & 同上                                     \\ \hline
    \end{tabular}
  \end{center}
  $x = 1$の場合が「$P$でない，かつ$Q$である」場合にあたるが，$x > 2$が成り立っていない以上，約束が破られたわけではない．
  どの$x$についても「$x>2$が成り立つのに$x>0$が成り立たない」ということは起こらないので，この命題は真である．
\end{example}

次に，含意$P \limp Q$から作られる命題に名前をつけよう．

\phantomsection
\begin{definition}{否定・逆・裏・対偶}{否定・逆・裏・対偶}
  命題$P \limp Q$に対して，以下を定義する．
  \begin{center}
    \begin{tabular}{|l|c|} \hline
      名称                      & 命題                  \\ \hline
      否定\index{ひてい@否定}   & $\lnot ( P \limp Q )$     \\ \hline
      逆\index{ぎゃく@逆}       & $Q \limp P$               \\ \hline
      裏\index{うら@裏}         & $\lnot P \limp \lnot Q$   \\ \hline
      対偶\index{たいぐう@対偶} & $\lnot Q \limp \lnot P$   \\ \hline
    \end{tabular}
  \end{center}
\end{definition}

まず，含意の否定について次が成り立つ．

\begin{prop}{含意の否定}{含意の否定}
  命題$P$，$Q$に対して，
  \[
    \lnot ( P \limp Q ) \equiv P \land \lnot Q
  \]
  である．
\end{prop}

\begin{tleftbar}
  \begin{proof}
    含意の等価表現$P \limp Q \equiv \lnot P \lor Q$とド・モルガンの法則を用いると，
    \[
      \lnot ( P \limp Q )
      \equiv \lnot ( \lnot P \lor Q )
      \equiv \lnot \lnot P \land \lnot Q
      \equiv P \land \lnot Q
    \]
    である（最後の同値で\prref{二重否定}を用いた）．これが証明すべきことであった．
  \end{proof}
\end{tleftbar}

すなわち，「$P$ならば$Q$」の否定は「$P$であるのに$Q$でない」である．「$P$ならば$Q$でない」（$P \limp \lnot Q$）ではないことに注意しよう．

続いて，逆・裏・対偶と元の命題の関係を，真理値表を用いて調べよう．

\begin{prop}{対偶との同値}{対偶との同値}
  命題$P$，$Q$に対して，
  \[
    P \limp Q \equiv \lnot Q \limp \lnot P
  \]
  である．すなわち，命題とその対偶の真偽は一致する．同様に，逆$Q \limp P$と裏$\lnot P \limp \lnot Q$の真偽も一致する．

  一方，$P \limp Q$と逆$Q \limp P$の真偽は，一般には一致しない．
\end{prop}

\begin{tleftbar}
  \begin{proof}
    $P$，$Q$の真偽の組合せ$4$通りすべてについて，各命題の真偽を列挙する：
    \begin{center}
      \begin{tabular}{|c|c||c|c||c|c|c|c|} \hline
        $P$          & $Q$          & $\lnot P$    & $\lnot Q$    & $P \limp Q$  & $\lnot Q \limp \lnot P$ & $Q \limp P$  & $\lnot P \limp \lnot Q$ \\ \hline
        $\mathrm{T}$ & $\mathrm{T}$ & $\mathrm{F}$ & $\mathrm{F}$ & $\mathrm{T}$ & $\mathrm{T}$          & $\mathrm{T}$ & $\mathrm{T}$          \\ \hline
        $\mathrm{T}$ & $\mathrm{F}$ & $\mathrm{F}$ & $\mathrm{T}$ & $\mathrm{F}$ & $\mathrm{F}$          & $\mathrm{T}$ & $\mathrm{T}$          \\ \hline
        $\mathrm{F}$ & $\mathrm{T}$ & $\mathrm{T}$ & $\mathrm{F}$ & $\mathrm{T}$ & $\mathrm{T}$          & $\mathrm{F}$ & $\mathrm{F}$          \\ \hline
        $\mathrm{F}$ & $\mathrm{F}$ & $\mathrm{T}$ & $\mathrm{T}$ & $\mathrm{T}$ & $\mathrm{T}$          & $\mathrm{T}$ & $\mathrm{T}$          \\ \hline
      \end{tabular}
    \end{center}
    表より，$P \limp Q$の列と$\lnot Q \limp \lnot P$の列は$4$通りすべての場合で一致しているので，$P \limp Q \equiv \lnot Q \limp \lnot P$である．同様に，$Q \limp P$の列と$\lnot P \limp \lnot Q$の列も一致している．

    一方，$P \limp Q$の列と$Q \limp P$の列は，$2$行目・$3$行目で真偽が異なっている．よって$P \limp Q$と$Q \limp P$は同値でない．これが証明すべきことであった．
  \end{proof}
\end{tleftbar}

\begin{remark}[対偶と背理法]
  命題とその対偶の真偽が一致することから，$P \limp Q$を直接示す代わりに，対偶$\lnot Q \limp \lnot P$を示してもよい．これを\emph{対偶法}\index{たいぐうほう@対偶法}（対偶による証明）という．のちに紹介する背理法と並んで，重要な証明法のひとつである．
\end{remark}

\phantomsection
\begin{example}{逆・裏・対偶の例}{逆・裏・対偶の例}
  実数$x$に関する命題「$x > 2$ならば$x > 0$」（真）について，
  \begin{center}
    \begin{tabular}{|l|c|l|} \hline
      名称 & 命題                                  & 真偽                 \\ \hline
      逆   & 「$x > 0$ならば$x > 2$」              & 偽（反例：$x = 1$）  \\ \hline
      裏   & 「$x \leqq 2$ならば$x \leqq 0$」      & 偽（反例：$x = 1$）  \\ \hline
      対偶 & 「$x \leqq 0$ならば$x \leqq 2$」      & 真（元の命題と同じ） \\ \hline
    \end{tabular}
  \end{center}
  このように，元の命題が真であっても，その逆・裏は偽でありうる．
\end{example}

\begin{remark}[「逆」の誤用に注意]
  「逆」という言葉は，日常語では「向きを反対にすること」全般を指すため，数学の議論において誤用されやすい．たとえば，三角関数の加法定理
  \[
    \sin (\alpha + \beta) = \sin \alpha \cos \beta + \cos \alpha \sin \beta
  \]
  を右辺から左辺へ向けて用いて，$\sin \alpha \cos \beta + \cos \alpha \sin \beta$を$\sin (\alpha + \beta)$にまとめる変形を「加法定理の逆」とよぶのは誤りである．等式$A = B$と$B = A$は同じ内容を表しており，等式をどちら向きに用いても「加法定理そのもの」だからである．\deref{否定・逆・裏・対偶}で定義したとおり，「逆」とは含意$P \limp Q$の仮定と結論を入れ替えた$Q \limp P$のことをいうのであって，等式を用いる向きを指す言葉ではない．

  一方，加法定理を用いて得られる含意については，通常の意味で逆を考えることができる．たとえば，実数$\alpha$，$\beta$に関する命題
  \[
    \alpha + \beta = \pi
    \quad \limp \quad
    \sin \alpha \cos \beta + \cos \alpha \sin \beta = 0
  \]
  は加法定理から従う．実際，仮定$\alpha + \beta = \pi$のもとで
  \[
    \sin \alpha \cos \beta + \cos \alpha \sin \beta
    = \sin (\alpha + \beta)
    = \sin \pi
    = 0
  \]
  となる．この命題の逆
  \[
    \sin \alpha \cos \beta + \cos \alpha \sin \beta = 0
    \quad \limp \quad
    \alpha + \beta = \pi
  \]
  は偽である．実際，$\alpha = \beta = 0$とすれば，仮定の等式は成り立つが，$\alpha + \beta = 0 \ne \pi$である．
\end{remark}

\phantomsection
\subsection{必要条件と十分条件}

含意$P \limp Q$は，数学の文章では「ならば」以外の言い回しで述べられることも多い．その代表が「必要条件」と「十分条件」である．

\phantomsection
\begin{definition}{必要条件・十分条件}{必要条件・十分条件}
  命題$P \limp Q$が真であるとき，$Q$は$P$であるための\emph{必要条件}\index{ひつようじょうけん@必要条件}であるといい，$P$は$Q$であるための\emph{十分条件}\index{じゅうぶんじょうけん@十分条件}であるという．

  さらに，$P \lequiv Q$が真であるとき，$P$は$Q$であるための（$Q$は$P$であるための）\emph{必要十分条件}\index{ひつようじゅうぶんじょうけん@必要十分条件}であるという．
\end{definition}

\begin{wrapfigure}[8]{r}{\mnsc{5.6cm}}
  \centering
  \vspace{-0.4\baselineskip}
  \begin{tikzpicture}[font=\small, line join=round]
    \draw[draw=black!60, fill=black!4, line width=0.7pt]
      (0,0) ellipse [x radius=2.3, y radius=1.15];
    \draw[draw=figaccent!70, fill=figaccent!12, line width=0.6pt]
      (-0.75,-0.1) ellipse [x radius=1.05, y radius=0.68];
    \node[text=figaccent, align=center, font=\footnotesize]
      at (-0.75,-0.1) {$P$が\\成り立つ};
    \node[text=black!75, align=center, font=\footnotesize]
      at (1.3,0.05) {$Q$が\\成り立つ};
  \end{tikzpicture}
  \caption{$P \limp Q$が真であるとき}
  \label{fig:必要条件と十分条件}
\end{wrapfigure}
「十分」という名前は，$Q$を導くには$P$が成り立てば十分である，ということに由来する．
一方，「必要」という名前は，対偶を考えるとわかりやすい．\prref{対偶との同値}より，$P \limp Q$は$\lnot Q \limp \lnot P$と同値である．すなわち，$Q$が成り立たなければ$P$も成り立たないのであって，$P$が成り立つためには$Q$が成り立つことが欠かせない（必要である）．

\figref{必要条件と十分条件}のように，$P$が成り立つ範囲が$Q$が成り立つ範囲にすっぽり含まれている様子を思い浮かべるとよい．内側の$P$に入っていれば外側の$Q$にも必ず入る（十分）が，$Q$の外にあるものは$P$にも入れない（必要）．この見方は，\partref{集合}で集合の包含関係として正確に述べ直される．

\phantomsection
\begin{example}{必要条件と十分条件の例}{必要条件と十分条件の例}
  実数$x$について考える．
  \begin{enumerate}[(A)]
    \item 「$x > 2$ならば$x > 0$」は真であるから，$x > 0$は$x > 2$であるための必要条件であり，$x > 2$は$x > 0$であるための十分条件である．
          一方，逆「$x > 0$ならば$x > 2$」は偽である（反例：$x = 1$）から，$x > 2$は$x > 0$であるための必要条件ではない．
    \item 「$x = 1$ならば$x^2 = 1$」は真であるが，その逆は偽である（反例：$x = -1$）．よって$x^2 = 1$は$x = 1$であるための必要条件であるが，十分条件ではない．
    \item 「$x^2 = 1$」と「$x = 1$または$x = -1$」は同値であるから，一方は他方であるための必要十分条件である．
  \end{enumerate}
\end{example}

\begin{remark}[「〜であるための」の向き]
  「$P$は$Q$であるための十分条件である」と「$P$であるための十分条件は$Q$である」とでは，含意の向きが逆になる．前者は$P \limp Q$を，後者は$Q \limp P$を意味する．
  どちらの向きの含意を述べているのかは，「〜であるための」がどちらに付いているかで判断する．$P \limp Q$の形に書き直してから考えれば，取り違えることはない．

  なお，「$P$であるための必要十分条件は$Q$である」は$P \lequiv Q$を意味し，「$P$のとき，かつそのときに限り$Q$」や「$P$と$Q$は同値である」と言い換えられる．英語では``if and only if''といい，iffと略記されることもある．
\end{remark}

\phantomsection
\subsection{恒真式と推論規則}

表\ref{tab:恒真命題と矛盾命題の真理値表}で見た$P \lor \lnot P$は，$P$の真偽によらずつねに真であった．このような論理式に名前をつけておこう．

\phantomsection
\begin{definition}{恒真式・矛盾式}{恒真式・矛盾式}
  論理式に含まれる命題（命題変数）の真偽をどのように定めても，つねに真となる論理式を\emph{恒真式}\index{こうしんしき@恒真式}（tautology）という．また，つねに偽となる論理式を\emph{矛盾式}\index{むじゅんしき@矛盾式}という．
\end{definition}

たとえば$P \lor \lnot P$は恒真式であり，$P \land \lnot P$は矛盾式である．
また，論理式$P$，$Q$について$P \equiv Q$であることは，$P \lequiv Q$が恒真式であることにほかならない．

証明の中で用いられる推論の多くは，恒真式として表すことができる．代表的なものを挙げよう．

\phantomsection
\begin{prop}{推論規則}{推論規則}
  命題$P$，$Q$，$R$に対して，次の論理式はいずれも恒真式である：
  \begin{description}
    \item[三段論法\index{さんだんろんぽう@三段論法}]
          \[
            \bigl( (P \limp Q) \land (Q \limp R) \bigr) \limp (P \limp R)
          \]
    \item[前件肯定\index{ぜんけんこうてい@前件肯定}]
          \[
            \bigl( P \land (P \limp Q) \bigr) \limp Q
          \]
    \item[場合分け\index{ばあいわけ@場合分け}]
          \[
            \bigl( (P \lor Q) \land (P \limp R) \land (Q \limp R) \bigr) \limp R
          \]
  \end{description}
\end{prop}

\begin{tleftbar}
  \begin{proof}
    表\ref{tab:含意の真理値表}より，含意が偽となるのは，前件が真で後件が偽となる場合だけである．そこで，いずれの論理式についても，後件が偽であるときには前件が偽となることを確かめればよい．

    三段論法について，$P \limp R$が偽であるとすると，$P$は真，$R$は偽である．このとき，$Q$が真ならば$Q \limp R$が偽となり，$Q$が偽ならば$P \limp Q$が偽となる．いずれにせよ前件$(P \limp Q) \land (Q \limp R)$は偽である．

    前件肯定について，$Q$が偽であるとする．$P$が偽ならば$P \land (P \limp Q)$は偽であり，$P$が真ならば$P \limp Q$が偽となるから，やはり$P \land (P \limp Q)$は偽である．

    場合分けについて，$R$が偽であるとする．前件のうち$P \limp R$と$Q \limp R$がともに真であるためには，$P$と$Q$がともに偽でなければならないが，このとき$P \lor Q$は偽である．よって前件は偽である．

    以上より，$3$つの論理式はいずれも恒真式である．これが証明すべきことであった．
  \end{proof}
\end{tleftbar}

\begin{remark}[推論規則と証明]
  \prref{推論規則}の$3$つの論理式は，$P$，$Q$，$R$が具体的にどのような命題であるかによらず，その「形」だけから真となる．そのため，どの分野の証明においても安心して用いることができる．

  たとえば，仮定$P$から出発して$P \limp Q_1$，$Q_1 \limp Q_2$，\dots，$Q_k \limp Q$を順に示し，結論$Q$を得るという証明は，三段論法を繰り返し用いたものである．また，「$P$の場合」と「$Q$の場合」に分けてそれぞれ$R$を示し，$P \lor Q$がつねに成り立つことから$R$を結論するという証明は，場合分けの規則に基づいている．これらの証明法については，\partref{証明}で具体例とともに見ていく．
\end{remark}


\phantomsection
\section{述語論理\index{じゅつごろんり@述語論理}}

\exref{命題でない文}の\ref{enu:命題でない文4}において，「$x^2 >4$」という文について考えた．この文は命題ではないが，たとえば$x=3$を代入したときは真，$x=0$を代入したときは偽である．
つまり，$x$になにかしら値を代入すると，この文は真偽が決まり，命題となるのである．

このような「値が決まっていない変数を含み，変数になにかを代入すると真偽が判定できる」という性質を持つ文を\emph{述語}\index{じゅつご@述語}という．
そして，述語を用いて命題を表現する論理を\emph{述語論理}\index{じゅつごろんり@述語論理}という．


述語論理においてはいくつかの記号を用いる．次の表で確認しよう：

\begin{table}[ht]
  \centering
  \caption{述語論理の記号}
  \begin{tabular}{c|c}
    \hline
    記号        & 意味                              \\
    \hline
    $\forall$ & 任意の（全称量化\index{ぜんしょうりょうか@全称量化}） \\
    $\exists$ & 存在する（存在量化\index{そんざいりょうか@存在量化}） \\
    \hline
  \end{tabular}
\end{table}

これらの記号を用いて，述語を命題にすることができる．$\forall$と$\exists$をまとめて\emph{量化子}\index{りょうかし@量化子}とよぶ．

\phantomsection
\begin{definition}{全称命題}{全称命題}
  $P(x)$を述語とする．
  \[
    \lall{x} \lst P(x)
  \]
  は，「すべての$x$について，$P(x)$が成り立つ」という意味であり，これを\emph{全称命題}\index{ぜんしょうめいだい@全称命題}という．
\end{definition}

\phantomsection
\begin{definition}{存在命題}{存在命題}
  $Q(x)$を述語とする．
  \[
    \lex{x} \lst Q(x)
  \]
  は，「$Q(x)$が成り立つような$x$が存在する」という意味であり，これを\emph{存在命題}\index{そんざいめいだい@存在命題}という．
\end{definition}

\phantomsection
\begin{example}{量化記号の例}{量化記号の例}
  \begin{itemize}
    \item $\lall{x \in \mathbb{R}} \lst x^2 \geqq 0$\\
          「任意の実数$x$について，$x^2$は$0$以上である．」
    \item $\lex{x \in \mathbb{Z}} \lst x^2 = 4$\\
          「ある整数$x$が存在して，$x^2 = 4$となる．」
  \end{itemize}
\end{example}

\begin{remark}[量化子つき記法の意味]
  本書で用いる
  \[
    \lall{x \in \mathbb{R}} \lst P(x)
  \]
  という記法は，厳密には
  \[
    \lall{x} \bigl( x \in \mathbb{R} \limp P(x) \bigr)
  \]
  の略記である．すなわち，「どの$x$についても，それが実数であるならば$P(x)$が成り立つ」という含意の形をしている．一方，
  \[
    \lex{x \in \mathbb{R}} \lst P(x)
  \]
  は
  \[
    \lex{x} \bigl( x \in \mathbb{R} \land P(x) \bigr)
  \]
  の略記であり，こちらは含意ではなく「かつ」であることに注意したい．全称は$\limp$と，存在は$\land$と組になる，というこの対応は，のちに述べる量化子の否定の規則とも整合的である\footnote{実際，$\lnot \lall{x} ( x \in \mathbb{R} \limp P(x) ) \equiv \lex{x} \lnot ( x \in \mathbb{R} \limp P(x) ) \equiv \lex{x} ( x \in \mathbb{R} \land \lnot P(x) )$となり，「$P(x)$を満たさない実数$x$が存在する」という正しい否定が得られる（含意の否定$\lnot(P \limp Q) \equiv P \land \lnot Q$を用いた）．}．
\end{remark}

\phantomsection
\subsection{自由変数と束縛変数}

述語$P(x) : x^2 > 4$における$x$のように，値が定まっておらず，なにを代入するかが自由に残されている変数を\emph{自由変数}\index{じゆうへんすう@自由変数}という．自由変数を含む文は，そのままでは真偽が定まらない．

一方，
\[
  \lall{x \in \mathbb{R}} \lst x^2 \geqq 0
\]
における$x$は，量化子$\forall$によって束縛されているといい，このような変数を\emph{束縛変数}\index{そくばくへんすう@束縛変数}という．束縛変数を含む文は，変数に値を代入する余地がすでにないので，全体として真偽の定まった命題となりうる．

束縛変数は，文字を付け替えても命題の意味が変わらない．たとえば，
\[
  \lall{x \in \mathbb{R}} \lst x^2 \geqq 0
  \qquad \text{と} \qquad
  \lall{y \in \mathbb{R}} \lst y^2 \geqq 0
\]
はまったく同じ命題である\footnote{同様の現象は，定積分$\int_0^1 x^2 \, dx = \int_0^1 t^2 \, dt$や総和$\sum_{k=1}^{n} k = \sum_{j=1}^{n} j$の変数にも見られる．これらの変数も「束縛」されており，ダミー変数とよばれることがある．}．

述語に含まれる変数のうち，どれが自由でどれが束縛されているかを意識することは，複雑な論理式を正確に読み書きするうえで重要である．とくに，量化子が2つ以上並ぶ場合（次項）には，この区別が議論の鍵となる．


\subsection{入れ子の量化子\index{いれこのりょうかし@入れ子の量化子}}
 
ここまでは変数がひとつの述語を扱ってきたが，$P(x, y)$のように変数を2つ以上含む述語も考えられる．
このとき，命題を作るにはすべての変数を量化する必要があり，量化子が
\[
  \lall{x} \lsep \lex{y} \lst P(x, y)
\]
のように2つ以上並ぶことになる．これを量化子の\emph{入れ子}\index{いれこのりょうかし@入れ子の量化子}という．
 
入れ子になった量化子は，左から順に読む．
すなわち，上の例は「任意の$x$に対して，ある$y$が存在して，$P(x, y)$が成り立つ」と読む．
このとき重要なのは，内側（右側）の変数は，外側（左側）の変数に依存してよいということである．
上の例でいえば，$y$は$x$に応じて選び直してよい．
 
\phantomsection
\subsubsection{同種の量化子の入れ子}
 
まず，同じ種類の量化子が並ぶ場合を考えよう．この場合，量化子の順序を入れ替えても命題の意味は変わらない：
 
\begin{prop}{同種の量化子の交換}{同種の量化子の交換}
  $P(x, y)$を2変数の述語とする．このとき，次が成り立つ：
  \begin{align*}
    \lall{x} \lsep \lall{y} \lst P(x, y) & \equiv \lall{y} \lsep \lall{x} \lst P(x, y) \\
    \lex{x} \lsep \lex{y} \lst P(x, y) & \equiv \lex{y} \lsep \lex{x} \lst P(x, y)
  \end{align*}
\end{prop}
 
直観的には，$\lall{x} \lsep \lall{y}$は「どの$x$とどの$y$の組合せについても」，
$\lex{x} \lsep \lex{y}$は「うまい$x$と$y$の組がひとつでもあれば」という意味であり，
いずれも$x$と$y$のどちらを先に述べるかは関係ない．
 
\phantomsection
\begin{example}{同種の量化子の例}{同種の量化子の例}
  \begin{itemize}
    \item $\lall{x \in \mathbb{R}} \lsep \lall{y \in \mathbb{R}} \lst x^2 + y^2 \geqq 2xy$\\
          「任意の実数$x$と任意の実数$y$について，$x^2 + y^2 \geqq 2xy$が成り立つ．」\\
          これは$x^2 + y^2 - 2xy = (x - y)^2 \geqq 0$からわかる．
          $x$と$y$の役割を入れ替えても主張は変わらない．
    \item $\lex{x \in \mathbb{Z}} \lsep \lex{y \in \mathbb{Z}} \lst x^2 + y^2 = 5$\\
          「ある整数$x$とある整数$y$が存在して，$x^2 + y^2 = 5$となる．」\\
          実際，$(x, y) = (1, 2)$がその一例である．
          こちらも，$x$と$y$のどちらを先に見つけるかは問題にならない．
  \end{itemize}
\end{example}
 
\phantomsection
\subsubsection{異種の量化子の入れ子}
 
一方，$\forall$と$\exists$が混在する場合には，順序が本質的に重要になる．
\[
  \lall{x} \lsep \lex{y} \lst P(x, y)
\]
は「任意の$x$に対して，（$x$に応じた）$y$が存在する」という意味であり，$y$は$x$ごとに選び直してよい．
これに対して，
\[
  \lex{y} \lsep \lall{x} \lst P(x, y)
\]
は「ある$y$が存在して，その$y$は任意の$x$に対して$P(x, y)$を満たす」という意味であり，
$y$は$x$によらずひとつに固定されていなければならない．
後者のほうが強い主張であり，一般に次の一方向の含意のみが成り立つ：
 
\begin{prop}{異種の量化子の順序}{異種の量化子の順序}
  $P(x, y)$を2変数の述語とする．このとき，次が成り立つ：
  \[
    \bigl( \lex{y} \lsep \lall{x} \lst P(x, y) \bigr)
    \limp
    \bigl( \lall{x} \lsep \lex{y} \lst P(x, y) \bigr)
  \]
  この逆は，一般には成り立たない．
\end{prop}
 
すなわち，「全員に共通のうまい$y$がある」ならば「各人に応じたうまい$y$がある」のは当然だが，
その逆は保証されない．逆が成り立たない例として，次のアルキメデスの原理を見てみよう．
 
学部1年の微分積分学の講義で，アルキメデスの原理
\[
  \lall{a>0} \lsep \lex{n\in\mathbb{N}} \lst n>a
\]
を学ぶことがある．この命題を学習する際には，$\forall$と$\exists$が入れ子になっていることに注意しよう．
アルキメデスの原理の主張の解釈は「任意の正の実数$a$に対して，それより大きい自然数$n$が存在する」
である．これを
\[
  \lex{n\in\mathbb{N}} \lsep \lall{a>0} \lst n>a
\]
とかき換えると，意味は大きく変わってしまう．
こちらは「ある自然数$n$が存在して，すべての正の実数$a$より大きい」
という意味になるが，そのような自然数は存在しないので偽である．
 
このように，量化記号の順序を入れ替えると，
命題の意味だけでなく真偽まで変わることがある．
$\varepsilon$-$\delta$論法では，
\[
  \lall{\varepsilon>0} \lsep \lex{\delta>0}
\]
という形で量化記号が入れ子になっており，
$\delta$は選んだ$\varepsilon$に応じて決めてよいことを表している．
 
\phantomsection
\subsubsection{入れ子の量化子の否定}
 
最後に，量化子を含む命題の否定について述べる．
量化子の否定は，$\lnot$が量化子を通過するたびに$\forall$と$\exists$が入れ替わる，という規則に従う：
 
\begin{prop}{量化子の否定}{量化子の否定}
  $P(x)$を述語とする．このとき，次が成り立つ：
  \begin{align*}
    \lnot \bigl( \lall{x} \lst P(x) \bigr) & \equiv \lex{x} \lst \lnot P(x) \\
    \lnot \bigl( \lex{x} \lst P(x) \bigr) & \equiv \lall{x} \lst \lnot P(x)
  \end{align*}
\end{prop}
 
これはド・モルガンの法則の量化子版とみなせる．
実際，「すべての$x$で$P(x)$が成り立つ」の否定は「$P(x)$が成り立たない$x$が（少なくともひとつ）存在する」であり，
「$P(x)$が成り立つ$x$が存在する」の否定は「すべての$x$で$P(x)$が成り立たない」である．
 
この規則を入れ子の量化子に対して外側から順に適用すれば，入れ子になった命題の否定も機械的に作ることができる：
\begin{align*}
  \lnot \bigl( \lall{x} \lsep \lex{y} \lst P(x, y) \bigr)
  & \equiv \lex{x} \lsep \lnot \bigl( \lex{y} \lst P(x, y) \bigr) \\
  & \equiv \lex{x} \lsep \lall{y} \lst \lnot P(x, y)
\end{align*}
すなわち，量化子の列$\lall{x} \lsep \lex{y}$を$\lex{x} \lsep \lall{y}$にすべて反転させ，
最後に内側の述語を否定すればよい．
この操作は，$\varepsilon$-$\delta$論法で「収束しないこと」を定式化する際などに繰り返し現れる．
実際に，数列の収束を例として確かめてみよう．

\phantomsection
\begin{example}{数列が収束しないこと}{数列が収束しないこと}
  数列$(a_n)_{n \in \mathbb{N}}$が実数$\alpha$に収束するとは，任意の$\varepsilon > 0$に対して，十分大きな$n$について$\abs{a_n - \alpha} < \varepsilon$が成り立つこと，すなわち
  \[
    \lall{\varepsilon > 0} \lsep \lex{N \in \mathbb{N}} \lsep \lall{n \geqq N} \lst \abs{a_n - \alpha} < \varepsilon
  \]
  が成り立つことをいう\footnote{ここで$\lall{n \geqq N}$は，「$n \geqq N$を満たす任意の$n \in \mathbb{N}$に対して」の略記である．収束の定義については，微分積分学で詳しく学ぶ．}．
  量化子の列を外側から順に反転させ，最後に述語を否定すると，「$(a_n)$は$\alpha$に収束しない」ことは
  \[
    \lex{\varepsilon > 0} \lsep \lall{N \in \mathbb{N}} \lsep \lex{n \geqq N} \lst \abs{a_n - \alpha} \geqq \varepsilon
  \]
  と表される．
  すなわち，「ある$\varepsilon > 0$があって，どんなに大きな$N$をとっても，それ以降に$a_n$が$\alpha$から$\varepsilon$以上離れるような$n$が現れる」ということである．

  たとえば，$a_n = (-1)^n$で定まる数列は$1$に収束しない．
  実際，$\varepsilon = 1$とし，$N \in \mathbb{N}$を任意にとる．$n = 2N + 1$とおけば，$n \geqq N$であり，$n$は奇数であるから$a_n = -1$である．よって
  \[
    \abs{a_n - 1} = 2 \geqq 1 = \varepsilon
  \]
  となり，上の論理式が成り立つ．
\end{example}

\begin{remark}[否定をとるときの注意]
  \exref{数列が収束しないこと}の否定の操作について，いくつか注意しておこう．
  \begin{itemize}
    \item 不等式$\abs{a_n - \alpha} < \varepsilon$の否定は$\abs{a_n - \alpha} \geqq \varepsilon$であって，$\abs{a_n - \alpha} > \varepsilon$ではない．等号の扱いを誤りやすいので注意したい．
    \item $\lall{\varepsilon > 0}$を否定して得られるのは$\lex{\varepsilon > 0}$であり，「$\varepsilon > 0$」という範囲の指定そのものは否定しない．これは，量化子つき記法の意味に関する注意で述べたとおり，$\lnot \lall{x} ( x \in \mathbb{R} \limp P(x) ) \equiv \lex{x} ( x \in \mathbb{R} \land \lnot P(x) )$となることによる．$\lall{n \geqq N}$についても同様である．
  \end{itemize}
\end{remark}

\phantomsection
\part{証明}

\begin{abstract}
この章では，前章で整理した論理を踏まえて，数学の議論の中心である「証明」を扱う．まず証明とはなにかを確認し，帰納法・背理法・直接証明・対偶法・場合分けや反例による反証といったさまざまな証明法を具体例とともに見ていく．数学的帰納法については，その原理と累積帰納法もあわせて述べる．そのうえで，数学的主張の骨格をなす「存在」と「一意性」について，「数学では何を証明しなければならないのか」という観点から具体例をもとに整理し，構成的・非構成的な存在証明や一意性の示し方にも触れよう．
\end{abstract}

\phantomsection
\section{証明とは}
数学において，主張が正しいことを示すプロセスを\emph{証明}\index{しょうめい@証明}（proof）という．

\begin{remark}[定義は「証明すべきこと」ではない]
  \partref{用語の定義}で見たように，命題・定理・補題・系は証明されるべきものであるが，定義は一般に「証明すべきこと」でないことに注意したい．
  ただし，たとえば「\exref{「絶対値」の定義}において$\abs{x} \coloneqq \max \{ x, -x \}$を仮定して$\abs{x} = \sqrt{x^2}$を導く」など，異なる定義の同値性を示す場面では，定義が「証明すべきもの」となる．このことから，どの定義を採用しているか明確にすることが重要であることがわかる．
\end{remark}

以下で，証明の例をいくつか見てみよう：


\begin{example}{証明①（余りの計算）}{証明①（余りの計算）}
  「任意の奇数$n$について，$n^2$を$4$で割った余りは$1$である」という命題を証明しよう：

  奇数$n$は，ある整数$k$を用いて$n = 2k + 1$と表すことができる．このとき，
\[
  n^2 = (2k + 1)^2 = 4k^2 + 4k + 1 = 4(k^2 + k) + 1
\]
である．したがって，$n^2$を$4$で割った余りは$1$である．これが証明すべきことであった．\qed
\end{example}

このように証明することで，「任意の奇数$n$について，$n^2$を$4$で割った余りは$1$である」という命題が真であることが示され，したがってこの命題の主張は「疑いようのない真実」であることがわかる．

証明はひとつの命題にいくつか存在する場合がほとんどであり，たとえば，\exref{三平方の定理}の証明は100通り以上も存在することが知られている．
数学の学習を進めるにあたって，さまざまな証明を見てみることは，論理的な思考力を養ううえで非常に有益である．

\phantomsection
\begin{example}[
    sidebyside,
    righthand width=\mnsc{5cm},
    sidebyside gap=\mnsc{4mm},
  ]{証明（三平方の定理）}{証明（三平方の定理）}
  辺の長さが$a$と$b$の直角三角形を$4$つ用意し，
  これらを組み合わせて，辺の長さが$a + b$の正方形を作る．
  この正方形は，中央の辺の長さが$c$の正方形と$4$つの三角形からなる．

  したがって，面積を2通りに計算すると
  \[
    (a + b)^2
    = c^2 + 4\left(\frac{1}{2}ab\right)
    = c^2 + 2ab
  \]
  である．これを展開して
  \[
    a^2 + 2ab + b^2 = c^2 + 2ab,
    \quad \therefore\quad a^2 + b^2 = c^2
  \]
  を得る．
  これが証明すべきことであった．\qed

  \tcblower

  \centering
  \begin{tikzpicture}[scale=0.58]
    \def\a{3}
    \def\b{4}

    % 中央の正方形（面積 c^2）を塗る
    \fill[figaccent!12] (\b,0) -- (\a+\b,\b) -- (\a,\a+\b) -- (0,\a) -- cycle;
    \node[text=figaccent] at ({(\a+\b)/2},{(\a+\b)/2}) {$c^2$};

    % 大きな正方形
    \draw (0,0) rectangle (\a+\b,\a+\b);

    % 4つの直角三角形
    \draw (0,0) -- (\b,0) -- (0,\a) -- cycle;
    \draw (0,\a) -- node[midway,above] {$c$} (\b,0);

    \draw (\b,0) -- (\a+\b,0) -- (\a+\b,\b) -- cycle;
    \draw (\a+\b,\b) -- node[midway,above] {$c$} (\b,0);

    \draw (\a+\b,\b) -- (\a+\b,\a+\b) -- (\a,\a+\b) -- cycle;
    \draw (\a,\a+\b) -- node[midway,above] {$c$} (\a+\b,\b);

    \draw (\a,\a+\b) -- (0,\a+\b) -- (0,\a) -- cycle;
    \draw (0,\a) -- node[midway,above] {$c$} (\a,\a+\b);

    % 中央の正方形
    \draw[figaccent, thick] (\b,0) -- (\a+\b,\b) -- (\a,\a+\b) -- (0,\a) -- cycle;

    % 寸法
    \draw[<->] (-0.5,0) -- (-0.5,\a) node[midway,left] {$a$};
    \draw[<->] (-0.5,\a) -- (-0.5,\a+\b) node[midway,left] {$b$};
    \draw[<->] (0,-0.5) -- (\b,-0.5) node[midway,below] {$b$};
    \draw[<->] (\b,-0.5) -- (\a+\b,-0.5) node[midway,below] {$a$};
  \end{tikzpicture}
\end{example}


\begin{mycolumn}
  証明を記述する際には，文末に$\square$や$\blacksquare$をかくことが一般的である．そのほかに``Q.E.D.''\footnote{``quod erat demonstrandum''の略で，ラテン語で「証明されるべきであったこと」という意味である．}などの表現も用いられる．
\end{mycolumn}

\phantomsection
\section{さまざまな証明法}

命題が真であることを示す方法にも，偽であることを示す方法にも，さまざまなものがある．代表的なものを挙げてみよう．

\begin{description}[descnext]
  \item[帰納法\index{きのうほう@帰納法}]
        たとえば数学的帰納法は，数学的な主張が自然数全体に対して成り立つことを示すための証明法である．
  \item[背理法\index{はいりほう@背理法}]
        「$P$が成り立たないと仮定すると矛盾が生じる」という論理的な構造を用いて証明を行う方法である．
  \item[直接証明\index{ちょくせつしょうめい@直接証明}]
        仮定から出発し，定義やすでに知られている事実を用いた推論を積み重ねて，結論を直接導く方法である．とくに含意$P \limp Q$の形の命題については，$P$を仮定し，その仮定と既知の事実から直接$Q$を導く．
  \item[対偶法\index{たいぐうほう@対偶法}]
        含意$P \limp Q$を示す代わりに，それと同値な対偶$\lnot Q \limp \lnot P$を示す方法である（\prref{対偶との同値}）．
  \item[場合分け\index{ばあいわけ@場合分け}]
        \prref{推論規則}の場合分けの規則に基づき，考えるべき状況をいくつかの場合に分けて，それぞれの場合に結論を示す方法である．
  \item[反例による反証\index{はんれい@反例}]
        全称命題「すべての$x$について$P(x)$が成り立つ」が偽であることを示すには，$P(x)$が成り立たない$x$，すなわち反例をひとつ挙げれば十分である．
\end{description}

以下では，まず「$n \geqq 2$ を満たす任意の自然数について，$2^n > n$である」という命題を3通りのやり方で証明してみよう．

\phantomsection
\begin{example}{帰納法}{帰納法}
  \begin{enumerate}[(I)]
    \item $n=2$のとき，
          \[
            2^2 = 4 > 2
          \]

          となり，与えられた不等式は成立する．
    \item 自然数$k$を$k \geqq 2$を満たすように任意にとる．$2^k > k$ が成り立つと仮定する．このとき，
          \[
            2^{k+1} = 2 \cdot 2^k > 2k
          \]
          である．
          ここで，$k \geqq 2$であるから，$2k - (k + 1) = k - 1 \geqq 1$であるがゆえに，$2k > k + 1$である．

          したがって，

          \[
            2^{k+1} > 2k > k + 1
          \]
          である．
          ゆえに，$n=k+1$のときも$2^{k+1} > k + 1$である．
  \end{enumerate}

  以上の考察と数学的帰納法により，$2^n > n$ は $n \geqq 2$を満たすすべての自然数について成り立つ．これが証明すべきことであった．\qed
\end{example}

\phantomsection
\begin{example}{背理法}{背理法}
  $n \geqq 2$を満たす自然数で $2^n \leqq n$ となるものが存在すると仮定する．
  $2^n$ の正の約数は $1, 2, 2^2, \dots, 2^n$ であるから，その個数はちょうど $n + 1$ 個である．
  しかし，仮定より $2^n \leqq n$ なので，$2^n$ は $n$ 以下の自然数である．
  ここで，$n$ 以下の自然数 $m$ の正の約数は，すべて $1$ 以上 $m$ 以下，したがって $1$ 以上 $n$ 以下の自然数であり，しかも互いに異なるから，その個数は高々 $n$ 個である．

  これは，$2^n$ が $n + 1$ 個の正の約数を持つことと矛盾する．

  したがって先の仮定は誤りであり，背理法により$2^n > n$ が成り立つ．これが証明すべきことであった．\qed
\end{example}

\phantomsection
\begin{example}{直接証明}{直接証明}
  \[
    2^n = n \int_{1}^{2} x^{n-1} \, dx + 1
  \]
  となることはよい．さて，$x \in [1, 2]$ かつ $n \geqq 2$ なので，
  \[
    x^{n-1} \geqq x^{1} \geqq 1
  \]
  が成り立つ．したがって，
  \begin{align*}
    2^n & = n \int_{1}^{2} x^{n-1} \, dx + 1 \\
        & \geqq  n \int_{1}^{2} \, dx +1     \\
        & = n+1
  \end{align*}
  である．
  これと$n+1 >n$であることを併せると，$n \geqq 2$のとき$2^n > n$である．これが証明すべきことであった．\qed
\end{example}

最後に，反例による反証の例を見てみよう．全称命題は「すべての場合において成り立つ」ことを主張しているため，ひとつでも例外があれば命題全体が偽となる．したがって，命題が偽であることの証明としては，反例をひとつ示すだけで十分なのである．

\phantomsection
\begin{example}{反例による反証}{反例による反証}
  次の命題を考える：
  \[
    \lall{n \in \mathbb{N}} \lst \frac{n^3}{24} + 84 \geqq n^2
  \]
  この命題が成り立たないことを示すためには，ある$n$について不等式が成立しないことを示せばよい．実際に，$n = 15$のとき，
  \[
    \frac{15^3}{24} + 84 = \frac{3375}{24} + 84 = 140.625 + 84 = 224.625
  \]
  である．
  一方，
  \[
    15^2 = 225
  \]
  である．
  したがって，
  \[
    \frac{15^3}{24} + 84 = 224.625 < 225 = 15^2
  \]
  が成り立つ．
  このように，$n = 15$において不等式が成り立たないため，元の命題は偽であることがわかる．
\end{example}

また，全称命題が偽であることを証明する際には，「反例をどう見つけたか」ということは明記しない場合が多い．
たとえば\exref{反例による反証}の場合に$n=15$を見つけるには，微分法を用いて議論したりさまざまな方法があるが，そのような「反例を見つけるまでの過程」を説明する必要はなく，「$n=15$の場合に命題が成り立たないこと」のみを証明として書けばよい．

\phantomsection
\subsection{数学的帰納法}

\exref{帰納法}では，(I)と(II)を示したのち，「数学的帰納法により」と結論した．この推論が正しいことは，なにによって保証されているのだろうか．
その答えは，自然数全体の集合$\mathbb{N}$をどのように定義するかにかかっている．ここでは\cite{sugiura1}にならい，\axref{実数の公理}をみたす実数全体の集合$\mathbb{R}$を出発点として，$\mathbb{N}$をその部分集合として定義しよう．

\phantomsection
\begin{definition}{継承的集合}{継承的集合}
  $\mathbb{R}$の部分集合$S$が次の$2$つを満たすとき，$S$は\emph{継承的}\index{けいしょうてきしゅうごう@継承的集合}であるという：
  \begin{enumerate}[label=(\roman*),leftmargin=3.2em]
    \item $1 \in S$である．
    \item 任意の$x \in S$に対して，$x + 1 \in S$である．
  \end{enumerate}
  継承的な集合を\emph{連鎖}\index{れんさ@連鎖}ともいう\cite{akasetsuya}．
\end{definition}

たとえば，$\mathbb{R}$自身や区間$[1, \infty)$は継承的である．一方，区間$[1, 5]$は$5$を含むが$6$を含まないので，継承的ではない．
継承的集合は「$1$を含み，$1$を足す操作で閉じている」集合であるから，$1, 1 + 1, 1 + 1 + 1, \ldots$をすべて含んでいなければならない．そこで，これらの数「だけ」からなる集合を，すべての継承的集合に共通する元の集合として定義する．

\phantomsection
\begin{definition}{自然数}{自然数}
  \[
    \mathbb{N} \coloneqq \{ x \in \mathbb{R} \mid \text{$\mathbb{R}$の任意の継承的部分集合$S$に対して$x \in S$} \}
  \]
  と定め，$\mathbb{N}$の元を\emph{自然数}\index{しぜんすう@自然数}という．
\end{definition}

\phantomsection
\begin{prop}{最小の継承的集合}{最小の継承的集合}
  $\mathbb{N}$は次の性質をもつ：
  \begin{enumerate}[label=(\roman*),leftmargin=3.2em]
    \item $\mathbb{N}$は継承的である．
    \item $S$が$\mathbb{R}$の継承的部分集合ならば，$\mathbb{N} \subset S$である．
  \end{enumerate}
  すなわち，$\mathbb{N}$は最小の継承的集合である．
\end{prop}

\begin{tleftbar}
  \begin{proof}
    (i)について，任意の継承的集合$S$は$1$を含むから，$1 \in \mathbb{N}$である．
    次に，$x \in \mathbb{N}$とし，$S$を任意の継承的集合とする．$\mathbb{N}$の定義より$x \in S$であり，$S$は継承的であるから$x + 1 \in S$である．$S$は任意であったから，$x + 1 \in \mathbb{N}$である．よって$\mathbb{N}$は継承的である．

    (ii)は，$\mathbb{N}$の定義そのものである．実際，$x \in \mathbb{N}$ならば，$x$は継承的集合$S$に属する．これが証明すべきことであった．
  \end{proof}
\end{tleftbar}

この命題から，数学的帰納法の原理が直ちに従う．

\phantomsection
\begin{theorem}{数学的帰納法の原理}{数学的帰納法の原理}
  $H$を$\mathbb{N}$の部分集合とする．$H$が継承的ならば，$H = \mathbb{N}$である．
\end{theorem}

\begin{tleftbar}
  \begin{proof}
    $H$は継承的であるから，\prref{最小の継承的集合}(ii)より$\mathbb{N} \subset H$である．一方，仮定より$H \subset \mathbb{N}$である．両方向の包含が成り立つので，$H = \mathbb{N}$である．これが証明すべきことであった．
  \end{proof}
\end{tleftbar}

\thref{数学的帰納法の原理}を述語の言葉で述べ直したものが，見慣れた形の数学的帰納法である．

\phantomsection
\begin{cor}{数学的帰納法}{数学的帰納法}
  $P(n)$を自然数$n$に関する述語とする．次の$2$つが成り立つならば，任意の$n \in \mathbb{N}$に対して$P(n)$が成り立つ：
  \begin{enumerate}[label=(\roman*),leftmargin=3.2em]
    \item $P(1)$が成り立つ．
    \item 任意の$k \in \mathbb{N}$に対して，$P(k)$が成り立つならば$P(k+1)$が成り立つ．
  \end{enumerate}
  論理式で書けば，
  \[
    \Bigl( P(1) \land \bigl( \lall{k \in \mathbb{N}} \lst ( P(k) \limp P(k+1) ) \bigr) \Bigr)
    \limp
    \lall{n \in \mathbb{N}} \lst P(n)
  \]
  である．
\end{cor}

\begin{tleftbar}
  \begin{proof}
    $H \coloneqq \{ n \in \mathbb{N} \mid P(n) \}$とおく．(i)より$1 \in H$である．また，$k \in H$ならば，(ii)より$P(k+1)$が成り立ち，\prref{最小の継承的集合}(i)より$k + 1 \in \mathbb{N}$であるから，$k + 1 \in H$である．よって$H$は継承的であり，\thref{数学的帰納法の原理}より$H = \mathbb{N}$である．すなわち，任意の$n \in \mathbb{N}$に対して$P(n)$が成り立つ．これが証明すべきことであった．
  \end{proof}
\end{tleftbar}

(i)と(ii)が成り立てば，(i)より$P(1)$が成り立ち，(ii)で$k = 1$とすれば$P(2)$が，さらに$k = 2$とすれば$P(3)$が成り立つ．このようにドミノ倒しの要領で，どの自然数$n$についても，有限回の推論で$P(n)$にたどりつく．\coref{数学的帰納法}は，この「有限回の推論の繰り返し」を，すべての自然数について一度に結論してよいことを保証しているのである．

\begin{remark}[$\mathbb{N} = \{1, 2, 3, \ldots\}$との関係]
  自然数全体の集合を$\{1, 2, 3, \ldots\}$とかくことは多いが，「$\ldots$」は「以下同様」という直観に頼った表記であり，それ自体は定義になっていない．\deref{自然数}は，「$1$から出発して$1$を足し続けて得られる数全体」という直観を，「$\ldots$」を用いずに表したものである．
  たとえば，$1/2 \notin \mathbb{N}$であることは，$1/2$を含まない継承的集合$[1, \infty)$があることからわかる．同様に，$3/2 \notin \mathbb{N}$であることは，継承的集合$\{1\} \cup [2, \infty)$を考えればわかる．このように，$\mathbb{N}$に属さない実数は，それを含まない継承的集合によって「ふるい落とされる」のである．
\end{remark}

\begin{remark}[出発点を変えた帰納法]
  \exref{帰納法}では，$n = 1$ではなく$n = 2$から出発した．一般に，整数$m$を固定したとき，
  \begin{enumerate}[label=(\roman*),leftmargin=3.2em]
    \item $P(m)$が成り立つ，
    \item $k \geqq m$を満たす任意の整数$k$に対して，$P(k)$が成り立つならば$P(k+1)$が成り立つ
  \end{enumerate}
  の$2$つが示されれば，$n \geqq m$を満たす任意の整数$n$に対して$P(n)$が成り立つ．
  これは，自然数$n$に関する述語$Q(n)$を$Q(n) \coloneqq P(n + m - 1)$で定め，$Q$に\coref{数学的帰納法}を適用すれば得られる．
\end{remark}

帰納法の仮定として，$P(k)$だけでなく$P(1), P(2), \dots, P(k)$をすべて用いてよい形の帰納法もある．

\phantomsection
\begin{prop}{累積帰納法}{累積帰納法}
  $P(n)$を自然数$n$に関する述語とする．次の$2$つが成り立つならば，任意の$n \in \mathbb{N}$に対して$P(n)$が成り立つ：
  \begin{enumerate}[label=(\roman*),leftmargin=3.2em]
    \item $P(1)$が成り立つ．
    \item 任意の$k \in \mathbb{N}$に対して，$P(1), P(2), \dots, P(k)$がすべて成り立つならば$P(k+1)$が成り立つ．
  \end{enumerate}
\end{prop}

\begin{tleftbar}
  \begin{proof}
    自然数$n$に関する述語$Q(n)$を
    \[
      Q(n) \coloneqq \text{「$P(1), P(2), \dots, P(n)$がすべて成り立つ」}
    \]
    で定め，$Q(n)$に\coref{数学的帰納法}を適用する．
    まず，$Q(1)$は$P(1)$そのものであるから，(i)より成り立つ．
    次に，$k \in \mathbb{N}$を任意にとり，$Q(k)$が成り立つと仮定する．すなわち$P(1), \dots, P(k)$がすべて成り立つから，(ii)より$P(k+1)$も成り立ち，したがって$Q(k+1)$が成り立つ．
    よって数学的帰納法により，任意の$n \in \mathbb{N}$に対して$Q(n)$が成り立ち，とくに$P(n)$が成り立つ．これが証明すべきことであった．
  \end{proof}
\end{tleftbar}

\prref{累積帰納法}の形の帰納法を\emph{累積帰納法}\index{るいせききのうほう@累積帰納法}（強い帰納法）という．
\prref{累積帰納法}も，出発点を$n = 1$以外に変えて用いることができる．
累積帰納法が威力を発揮する例として，素因数分解の存在を示そう．ここで，$2$以上の自然数であって，$1$とそれ自身以外に正の約数をもたないものを\emph{素数}\index{そすう@素数}という．

\phantomsection
\begin{prop}{素因数分解の存在}{素因数分解の存在}
  $2$以上の任意の自然数は，有限個の素数の積として表せる（素数$1$個からなる積も含める）．
\end{prop}

\begin{tleftbar}
  \begin{proof}
    $n \geqq 2$を満たす自然数$n$に関する述語$P(n)$を「$n$は有限個の素数の積として表せる」とし，$n = 2$を出発点とする累積帰納法で示す．

    $n = 2$のとき，$2$は素数であるから$P(2)$は成り立つ．

    次に，$k \geqq 2$を満たす自然数$k$を任意にとり，$2 \leqq m \leqq k$を満たすすべての自然数$m$について$P(m)$が成り立つと仮定する．$k + 1$が素数ならば，$P(k+1)$は成り立つ．$k + 1$が素数でないならば，$1$と$k + 1$以外の正の約数$a$をもつので，
    \[
      k + 1 = ab, \qquad 2 \leqq a \leqq k,\quad 2 \leqq b \leqq k
    \]
    と表せる．仮定より$a$と$b$はそれぞれ有限個の素数の積として表せるから，それらをかけあわせれば，$k + 1$も有限個の素数の積として表せる．よって$P(k+1)$が成り立つ．

    以上と累積帰納法により，$2$以上の任意の自然数$n$について$P(n)$が成り立つ．これが証明すべきことであった．
  \end{proof}
\end{tleftbar}

この証明では，$k + 1$を分解して現れる$a,b$が「$k$」であるとは限らないので，$P(k)$だけを仮定する通常の帰納法ではうまくいかない．$k$以下のすべての自然数についての仮定が必要になるところに，累積帰納法の意義がある．

\begin{mycolumn}
  本書では，実数の性質を既知として，自然数を$\mathbb{R}$の最小の継承的部分集合として定義した．これとは別に，実数を用いずに自然数を直接公理的に定める方法もある．イタリアの数学者ペアノ（G.~Peano）は，自然数全体の集合を「$1$から出発し，『次の数』をとる操作を繰り返して得られるもの」として公理的に定式化した．このペアノの公理の中には，数学的帰納法が公理としてそのまま含まれている．

  また，数学的帰納法は，\emph{整列性}\index{せいれつせい@整列性}とよばれる性質「$\mathbb{N}$の空でない部分集合には，最小の元が存在する」と同値であることが知られている．整列性を用いると，「$P(n)$が成り立たない$n$が存在したとすれば，そのような$n$のうち最小のものがとれる」として矛盾を導く，という形の証明（最小反例による証明）が可能になる．これは，背理法と帰納法を組み合わせた論法とみることができる．
\end{mycolumn}

\phantomsection
\subsection{対偶法と背理法}

直接示すことが難しい命題でも，対偶を考えると見通しがよくなることがある．

\phantomsection
\begin{prop}{平方数の偶奇}{平方数の偶奇}
  整数$n$について，$n^2$が偶数ならば，$n$は偶数である．
\end{prop}

\begin{tleftbar}
  \begin{proof}
    対偶「$n$が奇数ならば，$n^2$は奇数である」を示す．$n$が奇数であるとすると，ある整数$k$を用いて$n = 2k + 1$と表せる．このとき，
    \[
      n^2 = 4k^2 + 4k + 1 = 2(2k^2 + 2k) + 1
    \]
    であり，$2k^2 + 2k$は整数であるから，$n^2$は奇数である．よって対偶が示されたので，\prref{対偶との同値}より元の命題も成り立つ．これが証明すべきことであった．
  \end{proof}
\end{tleftbar}

仮定「$n^2$は偶数」から直接出発すると，$n^2 = 2m$という式から$n$の形を取り出すことは難しい．一方，対偶を考えれば，$n$の形から出発して$n^2$を計算するだけで済む．
この命題を用いて，背理法の古典的な例を$2$つ紹介しよう．

\phantomsection
\begin{prop}{$\sqrt{2}$の無理性}{根号2の無理性}
  $\sqrt{2}$は無理数である．
\end{prop}

\begin{tleftbar}
  \begin{proof}
    $\sqrt{2}$が有理数であると仮定する．すると，互いに素な\footnote{$1$以外に正の公約数をもたないことをいう．分数は，約分を繰り返せばいつでもこの形にできる．}整数$a$と自然数$b$を用いて$\sqrt{2} = a/b$と表せる．両辺を$2$乗して分母を払うと，
    \[
      a^2 = 2b^2
    \]
    である．よって$a^2$は偶数であり，\prref{平方数の偶奇}より$a$も偶数である．そこで$a = 2c$（$c$は整数）とおくと，$4c^2 = 2b^2$，すなわち$b^2 = 2c^2$となり，同様にして$b$も偶数である．
    すると$a$と$b$はともに$2$を約数にもつことになり，$a$と$b$が互いに素であることに矛盾する．
    したがって$\sqrt{2}$は有理数ではない．これが証明すべきことであった．
  \end{proof}
\end{tleftbar}

\phantomsection
\begin{prop}{素数の無限性}{素数の無限性}
  素数は無限に存在する．
\end{prop}

\begin{tleftbar}
  \begin{proof}
    素数が有限個しか存在しないと仮定し，それらを$p_1, p_2, \dots, p_k$とする．
    \[
      N \coloneqq p_1 p_2 \dotsm p_k + 1
    \]
    とおくと，$N \geqq 2$であるから，\prref{素因数分解の存在}より，$N$はある素数$p$で割り切れる．仮定より$p$は$p_1, \dots, p_k$のいずれかであるから，$p$は$p_1 p_2 \dotsm p_k$も割り切る．
    よって$p$は差$N - p_1 p_2 \dotsm p_k = 1$を割り切ることになるが，$p \geqq 2$であるから，これは矛盾である．
    したがって，素数は無限に存在する．これが証明すべきことであった．
  \end{proof}
\end{tleftbar}

\begin{remark}[背理法と対偶法の違い]
  含意$P \limp Q$を示すとき，対偶法では$\lnot Q$を仮定して$\lnot P$を導く．一方，背理法では$P$と$\lnot Q$をともに仮定し，なんらかの矛盾を導く．
  対偶法による証明は，「$\lnot P$が導かれ，仮定$P$と矛盾する」と読めば背理法による証明ともみなせるが，背理法では導く矛盾が$\lnot P$である必要はないので，背理法のほうが適用できる範囲は広い．

  ただし，背理法の証明の途中に現れる主張は，「最終的には否定される仮定」のもとでの主張であり，それ自体は正しいとは限らない．そのため，背理法の証明は読み手にとって見通しが悪くなりがちである．直接証明や対偶法で示せる場合には，そちらを選ぶほうがよいことが多い．
\end{remark}

\phantomsection
\subsection{場合分けによる証明}

\prref{推論規則}の場合分けの規則によれば，$P \lor Q$がつねに成り立つとき，「$P$ならば$R$」と「$Q$ならば$R$」をそれぞれ示せば$R$が結論できる．場合が$3$つ以上あっても同様である．

\phantomsection
\begin{prop}{平方数を$3$で割った余り}{平方数を3で割った余り}
  任意の整数$n$について，$n^2$を$3$で割った余りは$0$または$1$である．
\end{prop}

\begin{tleftbar}
  \begin{proof}
    任意の整数$n$は，ある整数$k$を用いて$n = 3k$，$3k + 1$，$3k + 2$のいずれかの形に表せる．それぞれの場合について，
    \begin{align*}
      (3k)^2     & = 3 \cdot 3k^2,                \\
      (3k + 1)^2 & = 3(3k^2 + 2k) + 1,            \\
      (3k + 2)^2 & = 9k^2 + 12k + 4 = 3(3k^2 + 4k + 1) + 1
    \end{align*}
    であるから，$n^2$を$3$で割った余りは，$n = 3k$のとき$0$，それ以外のとき$1$である．これが証明すべきことであった．
  \end{proof}
\end{tleftbar}

\begin{remark}[場合分けの網羅性]
  場合分けによる証明で最も注意すべきことは，分けた場合が考えるべき状況をすべて覆っていること（網羅性）である．\prref{平方数を3で割った余り}では，「任意の整数は$3k,3k+1,3k+2$のいずれかの形に表せる」ことがこれを保証している．
  一方，場合どうしが重なっていても，証明としては誤りにならない．たとえば実数$x$について「$x \geqq 0$の場合」と「$x \leqq 0$の場合」に分けてもよい．漏れは致命的であるが，重なりは許されるのである．

  なお，\partref{用語の定義}で述べた「一般性を失わず」は，対称性によって同じ議論で済む場合を省略するための言い回しであり，これも場合分けの一種である．
\end{remark}

\phantomsection
\section{数学における存在と一意性}

ここまでで，証明とはなにか，そしてどのような証明法があるかを確認した．
本節では，数学的な主張そのものの論理構造に一歩踏み込み，「存在する」という主張と「ただひとつである」という主張について考える．
これらは単なる言い回しではなく，「数学では何を証明しなければならないのか」を決める骨格である．

\phantomsection
\subsection{存在}

数学の証明において，「存在\index{そんざい@存在}」は重要な概念である．といっても，あまりこのことを意識したことのない読者の方もいると思うので，この場を借りて具体例をもとに説明を試みることとする．

まず，以下に記す「最大値・最小値の定理」を考える：
\begin{quotation}
  $[a,b]$で連続な関数$f$に対して，$f$は$[a,b]$上で最大値と最小値を持つ．
\end{quotation}

「最大値・最小値が存在するなんて当たり前だ」と思われる読者もいるかもしれない．しかし，本当にそれは自明なのか．
実際には，関数の連続性や区間の閉有界性といった条件が揃って初めて，最大値や最小値の「存在」を保証できる．この定理を証明するにあたっては，厳密な数学的議論が必要となる．

「存在」が決して当たり前でないことがわかる例として，最大値が存在しない場合を実際に見てみよう．

\phantomsection
\begin{example}{最大値が存在しない例}{最大値が存在しない例}
  開区間$(0,1)$上で定められた関数
  \[
    f(x) = x \qquad (0 < x < 1)
  \]
  を考える．$f$の値はいくらでも$1$に近づくが，$f$は$(0,1)$上で最大値を持たない．

  実際，任意の$c \in (0,1)$をとると，
  \[
    c < \frac{c+1}{2} < 1
  \]
  であるから，$\dfrac{c+1}{2}$も$(0,1)$に属し，
  \[
    f\!\left(\frac{c+1}{2}\right) = \frac{c+1}{2} > c = f(c)
  \]
  である．すなわち，どの$c \in (0,1)$を選んでも，$f(c)$より大きい値をとる点が$(0,1)$の中に存在する．したがって，「$f(c)$が最大値である」ような$c \in (0,1)$は存在しない．

  一方，考える区間を閉区間$[0,1]$に変えれば，関数$f(x)=x$は$x=1$で最大値$1$をとる．
\end{example}

このように，「閉区間」という仮定をわずかに弱めるだけで，最大値の存在は崩れてしまう．最大値・最小値の定理が主張する「存在」は，決して当たり前のことではなく，仮定と証明によってはじめて保証されるものなのである．

次に，「存在」を確認しないまま議論を進めると，誤った結論に到達しうることを見てみよう．


以下のような漸化式で定められた数列$(y_n)_{n \in \mathbb{N}}$を考える：
\[
  \begin{cases}
    y_1 =1 , \\
    y_{n+1}= 2y_n +2
  \end{cases}
\]
この数列の極限値が存在すると仮定して，それを$y$とおこう．
\begin{tleftbar}
  与えられた漸化式により，
  \begin{align*}
     & y = 2y +2,           \\
     & \therefore ~ y = -2.
  \end{align*}
  しかし，
  \[
    \lim_{n \to \infty} y_n = \infty
  \]
  なので\footnote{実際，$y_1 = 1$と漸化式から，数学的帰納法により任意の$n \in \mathbb{N}$で$y_n \geqq 2^{n-1}$が示される．よって$y_n \to \infty$である．}，$y=-2$は誤りである．
\end{tleftbar}

このように，数列の極限が存在することを確認しないと，誤った結論に到達することがある．この例からわかるように，存在を確認することは重要なことなのである．

なお，数列の極限の存在を保証する定理として，「上に有界な単調増加数列は収束する」という\emph{単調収束定理}\index{たんちょうしゅうそくていり@単調収束定理}がある．このような存在定理については，微分積分学で本格的に学ぶことになる．「極限を求める計算」は，「極限の存在」の保証とセットになってはじめて意味を持つのである．

\phantomsection
\subsection{構成的な存在証明と非構成的な存在証明}

存在を示す証明には，大きく分けて$2$つの型がある．
ひとつは，条件を満たす対象を具体的に作ってみせる\emph{構成的}\index{こうせいてきなそんざいしょうめい@構成的な存在証明}な証明である．たとえば「任意の有理数$a < b$に対して，$a < c < b$を満たす有理数$c$が存在する」ことは，$c = (a + b)/2$とおけば示される．
もうひとつは，対象を具体的に示さずに，存在することだけを結論する\emph{非構成的}\index{ひこうせいてきなそんざいしょうめい@非構成的な存在証明}な証明である．次の例は，非構成的な存在証明の典型である．

\phantomsection
\begin{prop}{無理数の無理数乗}{無理数の無理数乗}
  無理数$a,b$であって，$a^b$が有理数となるものが存在する．
\end{prop}

\begin{tleftbar}
  \begin{proof}
    $\sqrt{2}^{\sqrt{2}}$が有理数であるか否かで場合を分ける．

    $\sqrt{2}^{\sqrt{2}}$が有理数である場合は，\prref{根号2の無理性}より$\sqrt{2}$は無理数であるから，$a = b = \sqrt{2}$とおけばよい．

    $\sqrt{2}^{\sqrt{2}}$が無理数である場合は，$a = \sqrt{2}^{\sqrt{2}}$，$b = \sqrt{2}$とおけば，$a,b$はともに無理数であり，
    \[
      a^b = \bigl( \sqrt{2}^{\sqrt{2}} \bigr)^{\sqrt{2}} = \sqrt{2}^{\sqrt{2} \cdot \sqrt{2}} = \sqrt{2}^{\,2} = 2
    \]
    は有理数である．

    いずれの場合にも条件を満たす$a,b$が存在する．これが証明すべきことであった．
  \end{proof}
\end{tleftbar}

\begin{remark}[非構成的な証明からわかること]
  \prref{無理数の無理数乗}の証明は，条件を満たす$a,b$が存在することを確かに示しているが，$2$つの場合のどちらが実際に起こるのかは教えてくれない\footnote{実際には$\sqrt{2}^{\sqrt{2}}$は無理数であることが知られている（ゲルフォント--シュナイダーの定理）が，その証明はずっと難しい．}．
  このように，非構成的な証明は「存在すること」と「それがなにであるか」を切り離す．先に述べた最大値・最小値の定理も，最大値が存在することを保証するだけで，それがどの点でとられるかまでは教えてくれない．存在の主張を読むときには，それが具体的な対象を与えているのかどうかを意識しておくとよい．
\end{remark}

一方，\prref{無理数の無理数乗}は構成的に示すこともできる．条件を満たす$a,b$を具体的に与える例を見てみよう．

\phantomsection
\begin{example}{無理数の無理数乗の構成的な例}{無理数の無理数乗の構成的な例}
  $a = \sqrt{2}$，$b = \log_2 9$とおくと，$a,b$はともに無理数であり，$a^b$は有理数である．

  実際，$a = \sqrt{2}$が無理数であることは\prref{根号2の無理性}による．

  次に，$b = \log_2 9$が無理数であることを背理法で示す．$\log_2 9 > 0$であるから，$\log_2 9$が有理数であると仮定すると，正の整数$p,q$を用いて$\log_2 9 = p/q$と表せる．このとき$2^{p/q} = 9$であるから，両辺を$q$乗して
  \[
    2^p = 9^q
  \]
  を得る．$p \geqq 1$より左辺は偶数であり，右辺は奇数であるから，これは矛盾である．よって$b = \log_2 9$は無理数である．

  最後に，$a^b$を計算すると，
  \[
    a^b = \sqrt{2}^{\,\log_2 9} = \bigl( 2^{1/2} \bigr)^{\log_2 9} = 2^{(\log_2 9)/2} = 2^{\log_2 3} = 3
  \]
  となり，$a^b$は有理数である．
\end{example}

この例では，条件を満たす$a,b$が具体的に与えられており，「どれがそうなのか」まで明らかになっている．同じ命題であっても，証明の方法によって得られる情報が異なるのである．


\phantomsection
\subsection{一意性}

\begin{wrapfigure}[10]{r}{0.425\textwidth}
  %\centering
  \begin{tikzpicture}[scale=0.78]
    \begin{axis}[
        axis lines=middle,
        xmin=-1.2, xmax=1.2,
        ymin=-1.5, ymax=1.5,
        xtick=\empty, ytick=\empty, % 目盛を消す
        xlabel={}, ylabel={},       % ラベルを消す
        samples=200,
        domain=-1.2:1.2,
        width=10cm,
        height=8cm,
      ]
      % 関数 f(x) = x^3 のグラフ
      \addplot [black, thick] {x^3};

      % 点 (-1, -1) と (1, 1) を結ぶ割線
      \addplot [figaccent, thick] coordinates {(-1, -1) (1, 1)};

      % 平均値の定理を満たす点 x = -1/√3 の接線
      \def\xone{-1/sqrt(3)}
      \def\yone{(\xone)^3}
      \addplot [black!55, dashed] {1*(x - \xone) + \yone};

      % 平均値の定理を満たす点 x = 1/√3 の接線
      \def\xtwo{1/sqrt(3)}
      \def\ytwo{(\xtwo)^3}
      \addplot [black!55, dashed] {1*(x - \xtwo) + \ytwo};

      % 重要な点をマーク
      \addplot [only marks, mark=*, mark options={fill=black}] coordinates {(\xone, \yone) (\xtwo, \ytwo)};
      \addplot [only marks, mark=*, mark options={draw=figaccent, fill=figaccent}] coordinates {(-1, -1) (1, 1)}; % 割線の端点
    \end{axis}
  \end{tikzpicture}
\end{wrapfigure}

読者の中には線型代数の講義で「逆行列の一意性」などに触れた方もいると思われる．線型代数以外の分野でも，数学の学習を進める際に，「一意性\index{いちいせい@一意性}」が重要である場面は多い．なぜ「一意性」が重要であるのか，ひとつ例を挙げて考えてみることにしよう．

微分積分の講義で習う定理に「平均値の定理」というものがある．
その主張は
\begin{quotation}
  $[a,b]$で連続，$(a,b)$で微分可能な関数$f$に対して，
  \[
    \frac{f(b)-f(a)}{b-a}=f'(c)
  \]
  を満たす$ c \in (a,b)$が存在する．
\end{quotation}

というものである．証明は微分積分学に譲るとして，この定理の主張を少し変更してみよう：

\begin{quotation}
  $[a,b]$で連続，$(a,b)$で微分可能な関数$f$に対して，
  \[
    \frac{f(b)-f(a)}{b-a}=f'(c)
  \]
  を満たす$ c \in (a,b)$がただひとつ存在する．
\end{quotation}


ここで「$c \in (a,b)$が存在する」という主張を「$c \in (a,b)$がただひとつ存在する」というより強い主張に変更した．この主張は偽である．このことが問題になる状況を挙げよう．
$f(x)=x^3$という関数を考える．この関数は$[-1,1]$で連続，$(-1,1)$で微分可能である．
このとき，図のように，条件を満たす$c \in (-1,1)$は複数存在するので，「ただひとつ存在する」という主張は偽である．
さらに言えば，2本より多くこのような接線を引ける場合もある\footnote{各自で確認してみてほしい．}．

このことから，安易に「ただひとつ存在する」などと強い主張をすることは避けるべきであることがわかる．このことは平均値の定理に限らず，中間値の定理なども同様である．

\phantomsection
\subsection{存在と一意性の違い}

ここまでの考察からわかるように，「存在する」という主張と「ただひとつである」という主張は，別々に考える必要がある．
とくに，一意性を示すことは，それ自体では存在を保証しない．
典型的な例を整理すると，次のようになる．

\begin{center}
  \begin{tabular}{|c|p{\mnsc{9.5cm}}|}
    \hline
    \textbf{状況} & \textbf{例} \\ \hline

    存在し，一意である
    &
    正則行列\footnote{逆行列を持つ正方行列のことを正則行列という．}の逆行列，
    方程式$x^3 = 8$の実数解など．
    \\ \hline

    存在するが，一意ではない
    &
    方程式$x^2 = 4$の実数解など．
    また，平均値の定理によって存在が保証される$c$も，一般には一意とは限らない．
    \\ \hline

    存在しない
    &
    方程式$x^2 = -1$の実数解など．
    \\ \hline

    一意性だけが示されている
    &
    存在するかどうかとは別に，
    「もし存在するなら，それはただひとつである」とだけ示されている場合である．
    \prref{逆行列の一意性}がこの形の主張である．
    \\ \hline
  \end{tabular}
\end{center}

したがって，ある対象が「ただひとつ存在する」ことを示すには，
存在性と一意性の両方を示さなければならない．

\begin{remark}[「存在」と「一意性」は別の課題]
  とくに，「一意性だけが示されている」場合に注意しよう．
  \prref{逆行列の一意性}をあらためて見返すと，
  「存在するとすれば」とあるように，この命題は逆行列の存在を一切主張していない．
  実際，逆行列を持たない正方行列は存在する\footnote{たとえば零行列．}．
  このように，一意性の証明は存在の証明と切り離して，ときには存在の証明よりも先に行うことができる．
  「存在」と「一意性」は，証明においても別々に扱われるべきふたつの課題なのである．
\end{remark}

\phantomsection
\subsection{一意性の示し方}

一意性を示すときには，決まった型がある．すなわち，「条件を満たすものが$2$つあったとして，それらが等しいことを示す」のである．\prref{逆行列の一意性}の証明も，逆行列が$B,C$の$2$つあるとして$B = C$を導いていた．
ここでは，最大値を例にとって，この型をもう一度確認しよう．

\phantomsection
\begin{definition}{最大値・最小値}{最大値・最小値}
  $A$を$\mathbb{R}$の部分集合とする．$M \in A$であって，任意の$a \in A$に対して$a \leqq M$を満たすものを，$A$の\emph{最大値}\index{さいだいち@最大値}という．
  同様に，$m \in A$であって，任意の$a \in A$に対して$m \leqq a$を満たすものを，$A$の\emph{最小値}\index{さいしょうち@最小値}という．
\end{definition}

\phantomsection
\begin{prop}{最大値の一意性}{最大値の一意性}
  $A$を$\mathbb{R}$の部分集合とする．$A$の最大値は，存在するとすればただひとつである．
\end{prop}

\begin{tleftbar}
  \begin{proof}
    $M,M'$がともに$A$の最大値であるとする．$M' \in A$であり，$M$は$A$の最大値であるから，$M' \leqq M$である．同様に，$M \in A$であり，$M'$は$A$の最大値であるから，$M \leqq M'$である．よって$M = M'$である．これが証明すべきことであった．
  \end{proof}
\end{tleftbar}

最小値についても同様である．
このように一意性が保証されているからこそ，$A$の最大値・最小値（が存在するとき）を，それぞれ$\max A,\min A$という記号で表すことができる．

\begin{remark}[一意性と記号]
  $\max A$や$A^{-1},\sqrt{a}$のように，ある条件を満たす対象をひとつの記号で表すときには，その対象がただひとつ存在することが前提となっている．存在しなければ記号は意味をもたず，ただひとつでなければ，記号がどれを指すのかが定まらない（$\sqrt{4}$は「$2$乗して$4$になる数」のうち非負のもの，と定めることで一意性を確保している）．
  記号を導入する前に存在と一意性を確かめることは，\partref{写像}で述べるwell-definedの考え方にもつながっている．
\end{remark}

なお，一意性を示すときに「$M \ne M'$と仮定して矛盾を導く」という背理法の形をとる必要はない．上の証明のように$M = M'$を直接導くほうが，見通しのよい証明になる．

最後に，\deref{最大値・最小値}が「$m \in A$」と「任意の$a \in A$に対して$m \leqq a$」の$2$つの条件からなることを，具体的な問題で確認しておこう．後者は「$m$以下の不等式がつねに成り立つ」という全称の主張であり，前者は「その値を実際にとる」という存在の主張である．最小値を求めるときに見落とされやすいのは，前者である．

\phantomsection
\begin{example}{不等式と最小値}{不等式と最小値}
  正の実数$x,y$について，$F(x, y) = \left( x + \dfrac{1}{y} \right) \left( y + \dfrac{9}{x} \right)$の値全体の集合を$S$とし，$S$の最小値を求めよう．

  まず，次の議論は誤りである．相加平均と相乗平均の関係（\prref{相加平均と相乗平均の関係}）より
  \[
    x + \frac{1}{y} \geqq 2 \sqrt{\frac{x}{y}},
    \qquad
    y + \frac{9}{x} \geqq 6 \sqrt{\frac{y}{x}}
  \]
  であり，両辺は正であるから$F(x, y) \geqq 12$である．ここまでは正しく，「任意の$s \in S$に対して$12 \leqq s$」が示されている．しかし，ここから「最小値は$12$」と結論することはできない．$12 \in S$であることが示されていないからである．実際，$2$つの不等式で等号が成り立つのはそれぞれ$xy = 1$，$xy = 9$のときであり，これらは同時には起こらないので，$F(x, y) = 12$となることはない．すなわち$12 \notin S$であり，$12$は$S$の最小値ではない．

  正しくは，次のようにする．展開すると
  \[
    F(x, y) = xy + \frac{9}{xy} + 10
  \]
  であり，$xy > 0$であるから，\prref{相加平均と相乗平均の関係}より$xy + \dfrac{9}{xy} \geqq 6$，したがって任意の$s \in S$に対して$16 \leqq s$である．また，$x = 3$，$y = 1$とすると
  \[
    F(3, 1) = (3 + 1)(1 + 3) = 16
  \]
  であるから，$16 \in S$である．以上より，$S$の最小値は$16$である．
\end{example}

\begin{remark}[等号成立条件は必要か]
  \exref{不等式と最小値}の後半で，$16 \in S$を示すために挙げたのは$(x, y) = (3, 1)$という例ひとつだけであった．$16 \in S$は「$F(x, y) = 16$となる正の実数$x,y$が存在する」という存在の主張であるから，構成的な存在証明として条件を満たす例をひとつ与えれば十分であり，等号が成り立つための必要十分条件（$xy = 3$）まで求める必要はない．
  逆に，「最小値を求めたら等号成立条件を必ず書く」と覚えてしまうと，なにを示すためにそれを書いているのかが見えなくなる．書くべきなのは，最小値の定義の$2$つの条件がそれぞれ満たされていることである．

  なお，「任意の$s \in S$に対して$16 \leqq s$」と「$S$の最小値は$16$」はどちらも，「$S = [16, \infty)$」，すなわち$F(x, y)$が$16$以上のすべての実数の値をとるという主張よりも真に弱い．$S = [16, \infty)$を示すには，さらに別の議論が必要となる（\exref{値域を求める}も参照）．
\end{remark}

\phantomsection
\part{集合}

\begin{abstract}
この章では，現代数学の基礎をなす「集合」を導入する．まず集合と元を定義し，集合の相等・部分集合・真部分集合といった基本概念を確認する．次に，$\mathbb{N}$，$\mathbb{Z}$，$\mathbb{Q}$，$\mathbb{R}$，$\mathbb{C}$ などのよく使う集合の記号と，外延的記法・内包的記法という集合の表記法を整理する．続いて，空集合とべき集合を導入する．最後に，共通部分・和集合・差集合・補集合といった集合の演算を定義し，交換法則・結合律・分配律・ド・モルガンの法則などの基本性質を証明したのち，無限個の集合を扱うための集合族の和集合・共通部分を導入しよう．
\end{abstract}

\phantomsection
\section{集合の定義と表記法}

\phantomsection
\subsection{集合の定義}

現代数学の基礎をなす概念に「集合」や「写像」がある．まず「集合」からみていこう．
\phantomsection
\begin{definition}{集合}{集合}
  もののあつまりを\emph{集合}\index{しゅうごう@集合}という\footnote{公理的集合論の立場では，集合とは「無定義語」であるが，ここで詳しくは触れない．}．集合を構成するものを\emph{元}\index{げん@元}または\emph{要素}\index{ようそ@要素}といい，集合$A$の元が$a$であることを$a \in A$，$A \ni a$などと表す．
\end{definition}

\begin{remark}[集合の表記に関する注意]
  集合について，いくつか注意しておこう．
  \begin{itemize}
    \item 要素の順序は問わない．たとえば，$\{1,2,3\}=\{3,2,1\}$である．
    \item 同じ要素を重複して書いても同じ集合を表す．たとえば，$\{1,1,2,2,2,3\}=\{1,2,3\}$である．
    \item 異なる種類のものを要素にすることもできる．たとえば，$\{4,\{3\}\}$は数$4$と集合$\{3\}$を要素に持つ．
  \end{itemize}
\end{remark}

次に，ある集合の要素の一部または全部からなる集合を考える．これを部分集合という．まずはその定義を確認し，具体例はのちほど紹介する．

\begin{definition}{集合の相等}{集合の相等}
  集合$A$，$B$のすべての要素が一致しているとき，$A$と$B$は\emph{相等}\index{そうとう@相等}であるといい，これを
  \[
    A = B
  \]
  と表す．

  また，集合$A$と$B$が相等でないとき，
  \[
    A \ne B
  \]
  と表す．
\end{definition}

\begin{remark}[相等を考える意義]
  \deref{集合の相等}を考える意義がぴんとこない読者もいるかもしれない．
  だが，たとえば$\pi$，$\int_{0}^{1} \frac{4}{1+x^2} \, dx$，$4 \sum_{n=0}^{\infty} \frac{(-1)^n}{2n+1}$は異なる表現であるが，すべて円周率を表している．
  このような例からわかるように，「相等」を考えることは数学的な主張を正確に表現するために重要なことなのだ．
\end{remark}

\phantomsection
\begin{definition}[
    sidebyside,
    righthand width=\mnsc{4.2cm},
    sidebyside gap=\mnsc{4mm},
  ]{部分集合}{部分集合}
  集合$A$の要素がすべて集合$B$の要素でもあるとき，$A$は$B$の\emph{部分集合}\index{ぶぶんしゅうごう@部分集合}であるといい，これを
  \[
    A \subset B,\quad B \supset A
  \]
  などと表す\footnote{このことを$ A \subseteq B$，$B \supseteq A$とかくこともある．一般に，$A \subset B$，$ B \supset A$と書いた場合には$A=B$の場合を含んで意味をとる．}．

  \tcblower

  \centering
  \begin{tikzpicture}[font=\small, line join=round]
    \draw[draw=black!60, fill=black!4, line width=0.7pt]
      (0,0) ellipse [x radius=1.7, y radius=1.05];
    \draw[draw=figaccent!70, fill=figaccent!12, line width=0.6pt]
      (-0.45,-0.05) ellipse [x radius=0.85, y radius=0.6];
    \node[text=figaccent] at (-0.45,-0.05) {$A$};
    \node[text=black!70] at (1.05,0.35) {$B$};
    % A の元は B の元でもある
    \foreach \p in {(-0.85,0.2), (-0.15,-0.35)} \fill[black!85] \p circle (1.3pt);
    \foreach \p in {(0.95,-0.4), (0.55,0.6)}   \fill[black!85] \p circle (1.3pt);
  \end{tikzpicture}
\end{definition}

\begin{definition}{真部分集合}{真部分集合}
  集合$A$の要素がすべて集合$B$の要素であり，なおかつ$A \ne B$であるとき，$A$は$B$の\emph{真部分集合}\index{しんぶぶんしゅうごう@真部分集合}であるといい，これを
  \[
    A \subsetneq B,\quad B \supsetneq A
  \]
  などと表す．
\end{definition}

集合の相等は，部分集合の言葉を用いて次のように特徴づけられる．

\phantomsection
\begin{prop}{相等と相互包含}{相等と相互包含}
  集合$A$，$B$に対して，
  \[
    A = B
    \lequiv
    (A \subset B) \land (B \subset A)
  \]
  である．また，量化子を用いれば，
  \[
    A = B
    \lequiv
    \lall{x} \lst ( x \in A \lequiv x \in B )
  \]
  と表すこともできる．
\end{prop}

\begin{tleftbar}
  \begin{proof}
    \deref{集合の相等}より，$A = B$とは$A$と$B$のすべての要素が一致すること，すなわち「任意の$x$について，$x \in A$と$x \in B$の真偽が一致する」ことである．一方，$A \subset B$は「任意の$x$について，$x \in A$ならば$x \in B$」を，$B \subset A$は「任意の$x$について，$x \in B$ならば$x \in A$」を意味する．同値の等価表現$P \lequiv Q \equiv (P \limp Q) \land (Q \limp P)$により，両者は一致する．これが証明すべきことであった．
  \end{proof}
\end{tleftbar}

この命題は，集合の等式$A = B$を証明するための標準的な方法を与えている．すなわち，$A \subset B$と$B \subset A$の両方向の包含をそれぞれ示せばよい．のちに集合の分配律を証明する際にも，実際にこの方法を用いる．

包含関係$\subset$については，次の基本的な性質も成り立つ．

\phantomsection
\begin{prop}{包含関係の基本性質}{包含関係の基本性質}
  $A$，$B$，$C$を集合とするとき，次が成り立つ：
  \begin{enumerate}[label=(\roman*),leftmargin=3.2em]
    \item $A \subset A$である．
    \item $A \subset B$かつ$B \subset C$ならば，$A \subset C$である．
  \end{enumerate}
\end{prop}

\begin{tleftbar}
  \begin{proof}
    (i)について，任意の$x$に対して「$x \in A$ならば$x \in A$」は真であるから，$A \subset A$である．

    (ii)について，$x$を任意にとる．仮定より「$x \in A$ならば$x \in B$」と「$x \in B$ならば$x \in C$」がともに真であるから，三段論法（\prref{推論規則}）により「$x \in A$ならば$x \in C$」も真である．よって$A \subset C$である．これが証明すべきことであった．
  \end{proof}
\end{tleftbar}

\phantomsection
\subsection{集合の表記}

集合を表す方法には，おもに次の2つがある．

\begin{description}
  \item[外延的記法\index{がいえんてききほう@外延的記法}（列挙形式）]
        $\{a,b,c\}$のように，要素を列挙して表す．
  \item[内包的記法\index{ないほうてききほう@内包的記法}（条件形式）]
        $\{x\in\mathbb{N}\mid x\leqq5\}$のように，
        要素が満たす条件を用いて表す．
\end{description}

\begin{example}{集合の表記}{集合の表記}
  集合は，外延的記法と内包的記法のどちらを用いても表すことができる．
  たとえば，次の等式が成り立つ．
  \begin{align*}
    \{1,2,3\}
      &=
      \{n\in\mathbb{N}\mid n^2-4n+3\leqq0\},\\
    \{-2,-1,0,1,2\}
      &=
      \{n\in\mathbb{Z}\mid n^2\leqq4\},\\
    \{2,4,6,8,\dots\}
      &=
      \{2n\mid n\in\mathbb{N}\},\\
    \{1,4,9,16,25,\dots\}
      &=
      \{n^2\mid n\in\mathbb{N}\},\\
      \{-1,1\} &= \{x\in\mathbb{R}\mid x^{2026}=1\}.
  \end{align*}
\end{example}

また数学では，よく使われる集合に固有の記号を用いることがある．

\begin{reviewbox}[数の集合]
  \small
  \begin{tabular}{@{}p{14\zw}@{\hspace{1.5\zw}}l@{}}
    自然数（natural numbers）&
    $\mathbb{N} = \{1,2,3,\ldots\}$\footnotemark\\[0.35\baselineskip]

    整数（integers）&
    $\mathbb{Z} = \{\ldots,-2,-1,0,1,2,\ldots\}$\\[0.35\baselineskip]

    有理数（rational numbers）&
    $\mathbb{Q} = \left\{ \frac{a}{b} \relmiddle| a \in \mathbb{Z},\ b \in \mathbb{N} \right\}$\\[0.35\baselineskip]

    実数（real numbers）&
    $\mathbb{R}$\\[0.35\baselineskip]

    複素数（complex numbers）&
    $\mathbb{C} = \{ a + bi \mid a,b \in \mathbb{R} \}$
  \end{tabular}
  \footnotetext{自然数に$0$を含める流儀もあるが，本書では自然数に$0$を含めないものとする．}
\end{reviewbox}

$\mathbb{N}$，$\mathbb{R}$，$\mathbb{C}$は，それぞれ natural number，real number，complex number に由来する記号である．
また，$\mathbb{Z}$はドイツ語の Zahl（数）に，$\mathbb{Q}$は quotient（商）に由来する．
なお，$\mathbb{C}$の表記に現れる$i$は，$i^2 = -1$を満たす虚数単位である．$i$は内包的記法の中で動く変数ではなく，あらかじめ固定された定数であることに注意したい．


$x$がどの集合の中を動くかは，文脈により省略されることがある．たとえば，以下のような問題があったとする．
\begin{quotation}
  方程式
  \[
    x^2 + 1 = 0
  \]
  を解け．
\end{quotation}
この問題では，$x$がどの集合の中を動くかが明記されていない．
もし$x \in \mathbb{R}$の範囲で解を探すなら，任意の実数$x$に対して$x^2 + 1 \geqq 1 > 0$なので，解は存在しない．
一方，$x \in \mathbb{C}$の範囲で解を探すなら，解は$x = \pm i$である．
このように，「方程式を解く」という問題の答えは，どの集合の中で解を探しているかによって変わる．
$x$が動く集合は，文脈から読み取れたり明記されていたりする場合が多いが，答えが変わりうる以上，つねに意識しておきたいことがらである．

数の集合以外にも，特定の集合を表すために固有の記号を用いることがある．

\phantomsection
\begin{definition}{空集合}{空集合}
  要素をひとつも持たない集合を\emph{空集合}\index{くうしゅうごう@空集合}といい，
  \[
    \varnothing
  \]
  で表す．
\end{definition}

空集合は$\emptyset$と表記することもある\footnote{空集合をギリシャ文字の$\phi$で表記するのは誤りである．空集合の記号$\varnothing$はAndré Weilによって導入されたもので，ノルウェー語やデンマーク語などで用いられる文字Øに着想を得ている．}．

\phantomsection
\begin{example}{空集合の例}{空集合の例}
  空集合は，外延的記法では
  \[
    \{\}
  \]
  と表される．一方，内包的記法では，たとえば
  \[
    \{x \in \mathbb{R} \mid x^2=-1\},
    \qquad
    \{n \in \mathbb{N} \mid n<0\},
    \qquad
    \{x \in \mathbb{R} \mid x<x\}
  \]
  は，いずれも条件を満たす要素をひとつも持たないため，空集合である．
  したがって，
  \[
    \{\}
    =
    \{x \in \mathbb{R} \mid x^2=-1\}
    =
    \{n \in \mathbb{N} \mid n<0\}
    =
    \{x \in \mathbb{R} \mid x<x\}
    =
    \varnothing
  \]
  である．
\end{example}

\begin{remark}[空集合に関する注意]
  空集合について，いくつか注意しておこう．
  \begin{itemize}
    \item 空集合は要素をひとつも持たないので，どのような$x$に対しても，$x \in \varnothing$は常に偽である．いいかえれば，任意の$x$に対して$x \notin \varnothing$が成り立つ．
    \item 集合
          \[
            \{\varnothing\}
          \]
          は空集合$\varnothing$をひとつの要素として持つ集合であり，空集合ではない．
    \item 空集合の要素の個数は$0$である．一般に，有限集合$A$の要素の個数を$|A|$や$n(A)$と表すことがあり，この記法を用いれば
          \[
            |\varnothing| = 0, \qquad |\{\varnothing\}| = 1
          \]
          である．
  \end{itemize}
\end{remark}

\phantomsection
\begin{example}{所属関係の例}{所属関係の例}
  \[
    x \in \mathbb{Q}
  \]
  とかくことで，$x$は有理数であることを表す．
\end{example}

\phantomsection
\begin{example}{部分集合の例}{部分集合の例}
  \[
    \mathbb{R} \subset \mathbb{C}
  \]
  である．つまり，実数全体の集合は複素数全体の集合の部分集合である．
\end{example}

\phantomsection
\begin{prop}{空集合は任意の集合の部分集合}{空集合は任意の集合の部分集合}
  $A$を任意の集合とするとき，
  \[
    \varnothing \subset A
  \]
  である．つまり，空集合はすべての集合の部分集合である．
\end{prop}

\begin{tleftbar}
  \begin{proof}
    部分集合の定義（\deref{部分集合}）より，$\varnothing \subset A$を示すには，
    \begin{quotation}
      任意の$x$に対して，$x \in \varnothing$ならば$x \in A$である
    \end{quotation}
    ことを示せばよい．ところが，空集合は要素をひとつも持たないので，$x \in \varnothing$を満たす$x$は存在しない．すなわち，仮定「$x \in \varnothing$」はどのような$x$に対しても偽である．仮定が偽である命題は結論によらず真であるから\footnote{このように，仮定を満たすものが存在しないために無条件で真となることを，空虚な真（vacuous truth）ということがある．}，「$x \in \varnothing$ならば$x \in A$」は任意の$x$に対して真である．よって$\varnothing \subset A$が成り立つ．
  \end{proof}
\end{tleftbar}

\begin{remark}[背理法による見方]
  次のように考えることもできる．もし$\varnothing \subset A$が成り立たないとすると，「$\varnothing$の要素であって$A$に属さないもの」が少なくともひとつ存在しなければならない．しかし，$\varnothing$はそもそも要素を持たないので，そのような要素は存在しえない．よって$\varnothing \subset A$である．
\end{remark}

集合を要素とする集合の重要な例として，べき集合がある．

\phantomsection
\begin{definition}{べき集合}{べき集合}
  集合$A$に対して，$A$の部分集合全体からなる集合
  \[
    \mathcal{P}(A) \coloneqq \{ B \mid B \subset A \}
  \]
  を$A$の\emph{べき集合}\index{べきしゅうごう@べき集合}という．
\end{definition}

\prref{空集合は任意の集合の部分集合}により，つねに$\varnothing \in \mathcal{P}(A)$である．また，$A \subset A$なので$A \in \mathcal{P}(A)$である．

\phantomsection
\begin{example}{べき集合}{べき集合}
  $A=\{ 1, 2\}$とすると，
  \[
    \mathcal{P}(A) = \{ \varnothing, \{1\}, \{2\}, \{1, 2\}\}
  \]
  である．
\end{example}

\exref{べき集合}では，要素が$2$個の集合のべき集合は$4 = 2^2$個の要素をもっていた．一般に，次が成り立つ．

\phantomsection
\begin{prop}{べき集合の要素の個数}{べき集合の要素の個数}
  $A$を要素の個数が$n$個の有限集合とするとき，
  \[
    \abs{\mathcal{P}(A)} = 2^n
  \]
  である．
\end{prop}

\begin{tleftbar}
  \begin{proof}
    $n$に関する数学的帰納法（$n = 0$を出発点とする）で示す．

    $n = 0$のとき，$A = \varnothing$であり，$\varnothing$の部分集合は$\varnothing$のみであるから，$\mathcal{P}(\varnothing) = \{ \varnothing \}$より$\abs{\mathcal{P}(A)} = 1 = 2^0$である．

    次に，$k \geqq 0$を満たす整数$k$を任意にとり，要素の個数が$k$個の任意の集合について主張が成り立つと仮定する．$A$を要素の個数が$k + 1$個の集合とし，$a \in A$をひとつ固定して$A' \coloneqq A \setminus \{a\}$とおく．$A'$の要素の個数は$k$個である．
    $A$の部分集合は，$a$を含まないものと$a$を含むものに分けられる．
    \begin{itemize}
      \item $a$を含まない$A$の部分集合は，$A'$の部分集合にほかならない．仮定より，その個数は$2^k$個である．
      \item $a$を含む$A$の部分集合は，$A'$の部分集合$B$を用いて$B \cup \{a\}$と表され，異なる$B$からは異なる$B \cup \{a\}$が得られる．よって，その個数も$2^k$個である．
    \end{itemize}
    したがって$\abs{\mathcal{P}(A)} = 2^k + 2^k = 2^{k+1}$である．
    以上と数学的帰納法により，任意の$n \geqq 0$について主張が成り立つ．これが証明すべきことであった．
  \end{proof}
\end{tleftbar}

直感的には，$A$の部分集合をひとつ作ることは，$A$の各要素について「含める」か「含めない」かの$2$通りを選ぶことにあたるので，その総数は$2^n$通りとなるのである．この事実にちなんで，べき集合$\mathcal{P}(A)$を$2^A$とかく流儀もある．



\phantomsection
\section{集合の演算}

\subsection{共通部分・和集合・差集合・補集合}
集合に対して，以下のような演算を定義することができる．

\phantomsection
\begin{definition}{共通部分・和集合・差集合・補集合}{共通部分・和集合・差集合・補集合}
  $A$と$B$を集合とするとき，以下の定義をする．
  \begin{description}
    \item[共通部分\index{きょうつうぶぶん@共通部分}]\mbox{}\\
          $A$と$B$の両方に属する要素全体の集合を$A \cap B$と表す．これを``$A$ intersection $B$''と読む．
    \item[和集合\index{わしゅうごう@和集合}]\mbox{}\\
          $A$または$B$のいずれかに属する要素全体の集合を$A \cup B$と表す．これを``$A$ union $B$''と読む．
    \item[差集合\index{さしゅうごう@差集合}]\mbox{}\\
          $A$の要素で$B$に属さないもの全体の集合を$A \setminus B$と表す．
    \item[補集合\index{ほしゅうごう@補集合}]\mbox{}\\
          全体集合$X$\footnote{考察の対象すべてを含む集合をあらかじめひとつ固定して考えるとき，その集合を\emph{全体集合}\index{ぜんたいしゅうごう@全体集合}という．}とその部分集合$A \subset X$に対して，$A^c \coloneqq X \setminus A$を$A$の補集合という．これを``$A$ complement''と読む．
  \end{description}
\end{definition}

\begin{remark}[論理記号による言い換え]
  集合の演算は，\partref{論理}で学んだ論理記号を用いて次のように言い換えられる：
  \begin{align*}
    x \in A \cap B & \lequiv (x \in A) \land (x \in B), \\
    x \in A \cup B & \lequiv (x \in A) \lor (x \in B), \\
    x \in A \setminus B & \lequiv (x \in A) \land \lnot (x \in B), \\
    x \in A^c & \lequiv \lnot (x \in A) \qquad (x \in X)
  \end{align*}
  すなわち，$\cap$は$\land$に，$\cup$は$\lor$に，補集合は$\lnot$に対応している．また，部分集合についても
  \[
    A \subset B \lequiv \lall{x} \lst ( x \in A \limp x \in B )
  \]
  である．この対応により，命題論理の法則から集合の法則を導くことができる（のちのド・モルガンの法則の証明で実際に用いる）．
\end{remark}


\begin{figure}[ht]
  \centering
  % --- ベン図の共通設定 ---
  \tikzset{
    venn frame/.style={
      draw=black!75,
      line width=0.7pt
    },
    venn circle/.style={
      draw=black!85,
      line width=0.8pt
    },
    venn fill/.style={
      fill=magenta!15
    },
    venn label/.style={
      font=\small
    }
  }
  % --- 共通部分 ---
  \begin{minipage}{0.45\linewidth}
    \centering
    \begin{tikzpicture}[scale=0.9]
      % 塗り
      \begin{scope}
        \clip (2.2,2) circle (1.45);
        \fill[venn fill] (3.8,2) circle (1.45);
      \end{scope}

      % 輪郭
      \draw[venn frame] (0,0) rectangle (6,4);
      \draw[venn circle] (2.2,2) circle (1.45);
      \draw[venn circle] (3.8,2) circle (1.45);

      % ラベル
      \node[venn label] at (1.35,2.9) {$A$};
      \node[venn label] at (4.65,2.9) {$B$};
      \node[venn label, anchor=north east] at (5.85,3.85) {$X$};
    \end{tikzpicture}
    \caption{共通部分 $A \cap B$}
  \end{minipage}
  \hfill
  % --- 和集合 ---
  \begin{minipage}{0.45\linewidth}
    \centering
    \begin{tikzpicture}[scale=0.9]
      % 塗り
      \fill[venn fill] (2.2,2) circle (1.45);
      \fill[venn fill] (3.8,2) circle (1.45);

      % 輪郭
      \draw[venn frame] (0,0) rectangle (6,4);
      \draw[venn circle] (2.2,2) circle (1.45);
      \draw[venn circle] (3.8,2) circle (1.45);

      % ラベル
      \node[venn label] at (1.35,2.9) {$A$};
      \node[venn label] at (4.65,2.9) {$B$};
      \node[venn label, anchor=north east] at (5.85,3.85) {$X$};
    \end{tikzpicture}
    \caption{和集合 $A \cup B$}
  \end{minipage}

  \vspace{3mm}

  % --- 差集合 ---
  \begin{minipage}{0.45\linewidth}
    \centering
    \begin{tikzpicture}[scale=0.9]
      % A を塗る
      \fill[venn fill] (2.2,2) circle (1.45);

      % A \cap B の部分を白く戻す
      \begin{scope}
        \clip (2.2,2) circle (1.45);
        \fill[white] (3.8,2) circle (1.45);
      \end{scope}

      % 輪郭
      \draw[venn frame] (0,0) rectangle (6,4);
      \draw[venn circle] (2.2,2) circle (1.45);
      \draw[venn circle] (3.8,2) circle (1.45);

      % ラベル
      \node[venn label] at (1.35,2.9) {$A$};
      \node[venn label] at (4.65,2.9) {$B$};
      \node[venn label, anchor=north east] at (5.85,3.85) {$X$};
    \end{tikzpicture}
    \caption{差集合 $A \setminus B$}
  \end{minipage}
  \hfill
  % --- 補集合 ---
  \begin{minipage}{0.45\linewidth}
    \centering
    \begin{tikzpicture}[scale=0.9]
      % X 全体を塗る
      \fill[venn fill] (0,0) rectangle (6,4);

      % A を白く抜く
      \fill[white] (2.2,2) circle (1.45);

      % 輪郭
      \draw[venn frame] (0,0) rectangle (6,4);
      \draw[venn circle] (2.2,2) circle (1.45);

      % ラベル
      \node[venn label] at (1.35,2.9) {$A$};
      \node[venn label, anchor=north east] at (5.85,3.85) {$X$};
    \end{tikzpicture}
    \caption{補集合 $A^c$}
  \end{minipage}
\end{figure}


\phantomsection
\begin{example}{集合の演算の例}{集合の演算の例}
  たとえば，$A = \{1, 2, 3\}$，$B = \{3, 4, 5\}$とすると，
  \begin{itemize*}
    \item $A \cap B = \{3\}$
    \item $A \cup B = \{1, 2, 3, 4, 5\}$
    \item $A \setminus B = \{1, 2\}$
    \item $X=\{1,2,3,4,5\}$とすると，$A^c = \{4, 5\}$
  \end{itemize*}
\end{example}

\begin{prop}{集合の交換法則}{集合の交換法則}
  $A$，$B$を集合とするとき，
  \[
    A \cap B = B \cap A , \quad A \cup B = B \cup A
  \]
  である．
\end{prop}

\begin{tleftbar}
  \begin{proof}
    共通部分および和集合の定義から従う．
  \end{proof}
\end{tleftbar}

\begin{remark}[交換法則を考える意義]
  ここで「なぜ交換法則を考えるのか」という疑問が生じるかもしれない．当たり前に成り立つことだと思える読者もいるかもしれないが，よく考えてみると「実数の減法」や「関数（写像）の合成」，「行列の積」など，
  数学において一般に交換法則が成り立たないものは数多く存在する．集合に対して交換法則を考える意義はそこにあるのである．
\end{remark}

\phantomsection
\begin{prop}{集合の結合律}{集合の結合律}
  $A$，$B$，$C$を集合とするとき，
  \[
    (A \cap B) \cap C = A \cap (B \cap C), \quad (A \cup B) \cup C = A \cup (B \cup C)
  \]
  である．
\end{prop}

\begin{tleftbar}
  \begin{proof}
    前半を示す．任意の$x$に対して，
    \begin{align*}
      x \in (A \cap B) \cap C
      & \lequiv \bigl( (x \in A) \land (x \in B) \bigr) \land (x \in C) \\
      & \lequiv (x \in A) \land \bigl( (x \in B) \land (x \in C) \bigr) \\
      & \lequiv x \in A \cap (B \cap C)
    \end{align*}
    である．途中，$2$つ目の同値で命題論理の結合律を用いた．\prref{相等と相互包含}の論理式による特徴づけにより，$(A \cap B) \cap C = A \cap (B \cap C)$である．後半も同様に示される．これが証明すべきことであった．
  \end{proof}
\end{tleftbar}

結合律が成り立つので，$3$つ以上の集合の共通部分や和集合は，括弧を省略して$A \cap B \cap C$や$A \cup B \cup C$とかいてよい．

\phantomsection
\begin{prop}{空集合との演算}{空集合との演算}
  $A$を任意の集合とするとき，
  \[
    A \cup \varnothing = A, \quad A \cap \varnothing = \varnothing
  \]
  である．
\end{prop}

\begin{tleftbar}
  \begin{proof}
    まず$A \cup \varnothing = A$を示す．$x \in A \cup \varnothing$とすると，$x \in A$または$x \in \varnothing$である．$x \in \varnothing$を満たす$x$は存在しないので，$x \in A$でなければならない．よって$A \cup \varnothing \subset A$である．逆に，$x \in A$ならば和集合の定義から$x \in A \cup \varnothing$なので，$A \subset A \cup \varnothing$である．以上より$A \cup \varnothing = A$が成り立つ．

    次に$A \cap \varnothing = \varnothing$を示す．$x \in A \cap \varnothing$とすると，特に$x \in \varnothing$であるが，そのような$x$は存在しない．すなわち$A \cap \varnothing$は要素をひとつも持たないので，$A \cap \varnothing = \varnothing$である．
  \end{proof}
\end{tleftbar}

\begin{remark}[空集合と差集合・直積]
  差集合についても，
  \[
    A \setminus \varnothing = A, \quad \varnothing \setminus A = \varnothing
  \]
  が成り立つ．これは差集合の定義から直ちに確かめられる．

  また，集合$A$，$B$に対して，$A$の要素と$B$の要素の順序対$(a,b)$全体の集合を直積$A \times B$という（詳しくは\partref{写像}で定義する）．直積についても，
  \[
    A \times \varnothing = \varnothing, \quad \varnothing \times A = \varnothing
  \]
  が成り立つ．たとえば$(a, b) \in A \times \varnothing$となる順序対が存在するとすれば，$b \in \varnothing$でなければならないが，そのような$b$は存在しないからである．
\end{remark}

\subsection{集合の分配律}

\phantomsection
\begin{prop}[
    sidebyside,
    righthand width=\mnsc{4.6cm},
    sidebyside gap=\mnsc{4mm},
  ]{集合の分配律}{集合の分配律}
  $A$，$B$，$C$を集合とするとき，
  \begin{align*}
    A \cap ( B \cup C) & = (A \cap B) \cup (A \cap C), \\
    A \cup ( B \cap C) & = (A \cup B) \cap (A \cup C)
  \end{align*}
  右図の塗られた部分が，1番目の等式の両辺が表す集合である．

  \tcblower

  \centering
  \begin{tikzpicture}[scale=0.82, font=\small]
    \def\cA{(0,0.55) circle (1.15)}
    \def\cB{(-0.65,-0.5) circle (1.15)}
    \def\cC{(0.65,-0.5) circle (1.15)}
    % A ∩ (B ∪ C) を塗る
    \begin{scope}
      \clip \cA;
      \fill[figaccent!22] \cB;
      \fill[figaccent!22] \cC;
    \end{scope}
    \draw[draw=black!60, line width=0.7pt] (-2.2,-1.95) rectangle (2.2,2.0);
    \foreach \venncirc in {\cA, \cB, \cC} \draw[draw=black!85, line width=0.7pt] \venncirc;
    \node at (0,1.35) {$A$};
    \node at (-1.3,-1.0) {$B$};
    \node at ( 1.3,-1.0) {$C$};
  \end{tikzpicture}
\end{prop}

\begin{tleftbar}
  \begin{proof}
    まず，
    \[
      A \cap (B \cup C) \subset (A \cap B) \cup (A \cap C)
    \]
    を示す：

    $x \in A \cap (B \cup C)$とすると，$x$は$A$に属し，かつ$B \cup C$に属する．
    $x \in B \cup C$より，$x \in B$または$x \in C$のいずれかが成り立つので，場合を分けて考える．
    \begin{itemize}
      \item $x \in B$の場合，$x$は$A$と$B$の両方に属するので，$x \in A \cap B$である．
      \item $x \in C$の場合，$x$は$A$と$C$の両方に属するので，$x \in A \cap C$である．
    \end{itemize}
    いずれの場合でも，$x$は$A \cap B$と$A \cap C$の少なくとも一方に属するから，$x \in (A \cap B) \cup (A \cap C)$となる．したがって
    \[
      A \cap (B \cup C) \subset (A \cap B) \cup (A \cap C)
    \]
    が示された．

    次に，
    \[
      (A \cap B) \cup (A \cap C) \subset A \cap (B \cup C)
    \]
    を示す：

    $x \in (A \cap B) \cup (A \cap C)$とすると，$x \in A \cap B$または$x \in A \cap C$のいずれかが成り立つので，場合を分けて考える．
    \begin{itemize}
      \item $x \in A \cap B$の場合，$x$は$A$に属し，かつ$B$に属する．$B \subset B \cup C$であるから，$x$は$B \cup C$にも属する．
      \item $x \in A \cap C$の場合，$x$は$A$に属し，かつ$C$に属する．$C \subset B \cup C$であるから，$x$は$B \cup C$にも属する．
    \end{itemize}
    いずれの場合でも，$x$は$A$と$B \cup C$の両方に属するから，$x \in A \cap (B \cup C)$となる．したがって
    \[
      (A \cap B) \cup (A \cap C) \subset A \cap (B \cup C)
    \]
    が示された．

    以上の2つの包含関係から，
    \[
      A \cap (B \cup C) =(A \cap B) \cup (A \cap C)
    \]
    が成り立つ．後半の等式も同様に示される．これが証明すべきことであった．
  \end{proof}
\end{tleftbar}

\subsection{集合のさまざまな法則}

\phantomsection
\begin{prop}{べき等律・全体集合との演算・補集合の性質}{べき等律・全体集合との演算・補集合の性質}
  $X$を全体集合，$A \subset X$とするとき，次が成り立つ：
  \begin{align*}
    & A \cup A = A, && A \cap A = A, \\
    & A \cup X = X, && A \cap X = A, \\
    & A \cup A^c = X, && A \cap A^c = \varnothing, \\
    & (A^c)^c = A
  \end{align*}
\end{prop}

\begin{tleftbar}
  \begin{proof}
    いずれも定義から確かめられる．例として$(A^c)^c = A$を示す．$x \in X$に対して，
    \[
      x \in (A^c)^c
      \lequiv \lnot ( x \in A^c )
      \lequiv \lnot \lnot ( x \in A )
      \lequiv x \in A
    \]
    である（最後の同値は\prref{二重否定}による）．よって$(A^c)^c = A$である．残りも同様に示される．
  \end{proof}
\end{tleftbar}

差集合は，補集合と共通部分を用いて表すこともできる．

\phantomsection
\begin{prop}{差集合と補集合}{差集合と補集合}
  $X$を全体集合，$A, B \subset X$とするとき，
  \[
    A \setminus B = A \cap B^c
  \]
  が成り立つ．
\end{prop}

\begin{tleftbar}
  \begin{proof}
    $x \in X$に対して，
    \[
      x \in A \setminus B
      \lequiv (x \in A) \land \lnot (x \in B)
      \lequiv (x \in A) \land (x \in B^c)
      \lequiv x \in A \cap B^c
    \]
    である．よって$A \setminus B = A \cap B^c$である．これが証明すべきことであった．
  \end{proof}
\end{tleftbar}

\phantomsection
\begin{prop}{集合のド・モルガンの法則}{集合のド・モルガンの法則}
  $X$を全体集合，$A, B \subset X$とするとき，
  \[
    (A \cup B)^c = A^c \cap B^c,
    \qquad
    (A \cap B)^c = A^c \cup B^c
  \]
  が成り立つ．
\end{prop}

\begin{tleftbar}
  \begin{proof}
    前半を示す．$x \in X$に対して，
    \begin{align*}
      x \in (A \cup B)^c
      & \lequiv \lnot ( x \in A \cup B ) \\
      & \lequiv \lnot \bigl( (x \in A) \lor (x \in B) \bigr) \\
      & \lequiv \lnot (x \in A) \land \lnot (x \in B) \\
      & \lequiv x \in A^c \cap B^c
    \end{align*}
    である．途中，3つ目の同値では，命題論理のド・モルガンの法則
    \[
      \lnot (P \lor Q) \equiv \lnot P \land \lnot Q
    \]
    を用いた．\prref{相等と相互包含}の論理式による特徴づけにより，$(A \cup B)^c = A^c \cap B^c$である．後半も同様に示される．これが証明すべきことであった．
  \end{proof}
\end{tleftbar}

\begin{remark}[論理の法則と集合の法則]
  この証明からわかるように，集合のド・モルガンの法則は，命題論理のド・モルガンの法則を$P = (x \in A)$，$Q = (x \in B)$として集合の言葉に翻訳したものにほかならない．同様に，集合の交換法則・結合律・分配律も，命題論理の対応する法則から導くことができる．論理と集合は，このように同じ構造を共有しているのである．
\end{remark}

\phantomsection
\subsection{集合族}

ここまでは$2$つや$3$つの集合の共通部分・和集合を考えてきたが，数学では無限個の集合を同時に扱うことも多い．そのための記法を導入しよう．

\phantomsection
\begin{definition}{集合族}{集合族}
  集合$\Lambda$の各元$\lambda$に対して，全体集合$X$の部分集合$A_\lambda$がひとつずつ与えられているとき，これらをまとめて$(A_\lambda)_{\lambda \in \Lambda}$とかき，$\Lambda$を添字集合とする\emph{集合族}\index{しゅうごうぞく@集合族}という．
  集合族$(A_\lambda)_{\lambda \in \Lambda}$に対して，
  \begin{align*}
    \bigcup_{\lambda \in \Lambda} A_\lambda & \coloneqq \{ x \in X \mid \lex{\lambda \in \Lambda} \lst x \in A_\lambda \}, \\
    \bigcap_{\lambda \in \Lambda} A_\lambda & \coloneqq \{ x \in X \mid \lall{\lambda \in \Lambda} \lst x \in A_\lambda \}
  \end{align*}
  をそれぞれ集合族の和集合，共通部分という．
\end{definition}

$\Lambda = \{1, 2\}$のときは，$\bigcup_{\lambda \in \Lambda} A_\lambda = A_1 \cup A_2$，$\bigcap_{\lambda \in \Lambda} A_\lambda = A_1 \cap A_2$であり，これまでの和集合・共通部分と一致する．また，$\Lambda = \mathbb{N}$のときは，$\bigcup_{n \in \mathbb{N}} A_n$を$\bigcup_{n=1}^{\infty} A_n$ともかく．

\phantomsection
\begin{example}{集合族の和集合と共通部分}{集合族の和集合と共通部分}
  全体集合を$\mathbb{R}$とする．
  \begin{enumerate}[(A)]
    \item $\displaystyle \bigcup_{n \in \mathbb{N}} \left[ \frac{1}{n}, 1 \right] = (0, 1]$である．
          実際，各$n$について$[1/n, 1] \subset (0, 1]$であるから，左辺は右辺に含まれる．逆に，$x \in (0, 1]$とすると，アルキメデスの原理より$n > 1/x$を満たす$n \in \mathbb{N}$が存在し，$1/n < x \leqq 1$より$x \in [1/n, 1]$である．
    \item $\displaystyle \bigcap_{n \in \mathbb{N}} \left( -\frac{1}{n}, \frac{1}{n} \right) = \{0\}$である．
          実際，$0$はすべての$(-1/n, 1/n)$に属する．逆に，$x \ne 0$とすると，$n > 1/\abs{x}$を満たす$n \in \mathbb{N}$が存在し，$\abs{x} > 1/n$より$x \notin (-1/n, 1/n)$である．
  \end{enumerate}
  (A)では閉区間の和集合が閉区間でない区間$(0, 1]$に，(B)では開区間の共通部分が開区間でない集合$\{0\}$になっている．無限個の集合を扱うと，有限個の場合には見られない現象が起こりうるのである．
\end{example}

ド・モルガンの法則は，集合族に対しても成り立つ．

\phantomsection
\begin{prop}{集合族のド・モルガンの法則}{集合族のド・モルガンの法則}
  $X$を全体集合，$(A_\lambda)_{\lambda \in \Lambda}$を$X$の部分集合からなる集合族とするとき，
  \[
    \Bigl( \bigcup_{\lambda \in \Lambda} A_\lambda \Bigr)^c = \bigcap_{\lambda \in \Lambda} A_\lambda^c,
    \qquad
    \Bigl( \bigcap_{\lambda \in \Lambda} A_\lambda \Bigr)^c = \bigcup_{\lambda \in \Lambda} A_\lambda^c
  \]
  が成り立つ．
\end{prop}

\begin{tleftbar}
  \begin{proof}
    前半を示す．$x \in X$に対して，
    \begin{align*}
      x \in \Bigl( \bigcup_{\lambda \in \Lambda} A_\lambda \Bigr)^c
      & \lequiv \lnot \bigl( \lex{\lambda \in \Lambda} \lst x \in A_\lambda \bigr) \\
      & \lequiv \lall{\lambda \in \Lambda} \lst \lnot ( x \in A_\lambda ) \\
      & \lequiv \lall{\lambda \in \Lambda} \lst x \in A_\lambda^c \\
      & \lequiv x \in \bigcap_{\lambda \in \Lambda} A_\lambda^c
    \end{align*}
    である．途中，$2$つ目の同値で\prref{量化子の否定}を用いた．後半も同様に示される．これが証明すべきことであった．
  \end{proof}
\end{tleftbar}

\begin{remark}[添字集合が空の場合]
  集合族の和集合は$\exists$で，共通部分は$\forall$で定義されている．$2$つの集合の場合に$\cup$が$\lor$に，$\cap$が$\land$に対応していたことの，無限個への拡張である．

  この定義のもとで，添字集合が$\Lambda = \varnothing$の場合を考えてみよう．$\lex{\lambda \in \varnothing} \lst x \in A_\lambda$はどの$x$についても偽であるから，$\bigcup_{\lambda \in \varnothing} A_\lambda = \varnothing$である．一方，$\lall{\lambda \in \varnothing} \lst x \in A_\lambda$は仮定を満たす$\lambda$が存在しないためにどの$x$についても真となり（空虚な真），$\bigcap_{\lambda \in \varnothing} A_\lambda = X$である．
  「なにも集めなければ空集合，なにも条件を課さなければ全体集合」と考えれば自然な結果であるが，共通部分の定義に全体集合$X$が必要となる理由のひとつはここにある．
\end{remark}

\phantomsection
\part{写像}

\begin{abstract}
この章では，集合とならんで現代数学の基礎をなす「写像」を導入する．まず写像を「各元にただひとつの元を対応させる規則」として定義し，自動販売機などの身近な例で直感的な理解を深める．続いて，直積と記号$\exists!$を用いた写像の厳密な定義を紹介し，恒等写像や写像の制限を導入する．次に，像・値域・逆像といった写像に付随する集合を定義する．そのうえで，単射・全射・全単射という写像の重要な性質を，図と具体例をもとに確認し，像・逆像と集合演算の関係を整理する．最後に，写像の合成と逆写像を扱い，\partref{用語の定義}で触れたwell-definedの概念を写像の言葉で述べ直そう．
\end{abstract}

\phantomsection
\section{写像の定義と例}

% --- この章の写像図で共通して使う TikZ スタイル ---
\tikzset{
  mapfig/.style     ={>=Latex, font=\small, line cap=round, line join=round},
  mapset/.style     ={draw=black!60, fill=black!4, line width=0.7pt},
  mapsub/.style     ={draw=black!50, fill=black!14, line width=0.5pt},
  mapelem/.style    ={circle, fill=black!85, inner sep=1.3pt},
  maparrow/.style   ={-{Latex[length=1.8mm,width=1.3mm]}, draw=black!70,
                      line width=0.55pt, shorten <=2.5pt, shorten >=3pt},
  mapleader/.style  ={draw=black!50, line width=0.4pt},
  mapsetname/.style ={font=\small, text=black!70},
  maplabL/.style    ={anchor=east, xshift=-2.5pt},
  maplabR/.style    ={anchor=west, xshift=2.5pt},
  % 色は黒（灰）系＋マゼンタ系のみ．写像そのもの（対応の矢印と写像の名前）を
  % マゼンタで描き，集合・元は灰と黒で描く．
  mapname/.style    ={font=\small, text=figaccent},
  mapacc/.style     ={maparrow, draw=figaccent},
  mapaccsub/.style  ={draw=figaccent!60, fill=figaccent!10, line width=0.5pt},
  mapaccleader/.style={draw=figaccent!70, line width=0.4pt},
  mapacclab/.style  ={font=\small, text=figaccent!85!black},
  % 像に入らない元など「対応の相手がない」元は白抜きで示す
  mapelemfree/.style={circle, draw=black!85, fill=white, line width=0.6pt, inner sep=1.1pt},
}

\phantomsection
\subsection{写像の定義}

いままでの数学の学習の過程で「関数\index{かんすう@関数}」という言葉に出会ったことがある読者もいるかもしれない．
「関数」と言われると，$y = \sin x$や$y = x^3 + 1$のような「$x$と$y$を結ぶ式」，あるいはそのグラフを思い浮かべる読者が多いであろう．
しかし，これから見るように，関数の本質は式やグラフそのものではなく，「元にただひとつの元を対応させる規則」である．
式やグラフは，この規則を表現し，視覚化するための手段にすぎない．

関数の一般的な表現として「写像」がある\footnote{現代数学では，「関数」と「写像」はほぼ同義に用いられることも多い．本書では，一般の集合の間の対応であることを強調するときに「写像」という語を用いる．なお，定義域や終域が$\mathbb{R}$など「数」の集合である写像を「関数」とよぶことが多いが，これは絶対的な区別ではない．}．写像の厳密な定義はのちほど（本章「写像の厳密な定義」の項で）述べるとして，まず簡潔な定義を提示しよう．

\begin{definition}{写像}{写像}
  集合$A$から集合$B$への\emph{写像}\index{しゃぞう@写像}$f$とは，任意の$a \in A$に対して，$b \in B$をただひとつ対応させる規則のことである．

  このとき，$A$を$f$の\emph{定義域}\index{ていぎいき@定義域}（domain），$B$を$f$の\emph{終域}\index{しゅういき@終域}（codomain）とよぶ\footnote{定義域を「始集合」，終域を「終集合」とよぶ文献もある．}．

  これを次のように表す：
  \[
    f \colon A \to B
  \]
  この写像$f$によって，$a \in A$が$b \in B$に対応するとき，
  \[
    b=f(a)
  \]
  または
  \[
    f \colon a \mapsto b
  \]
  とかく\footnote{これを$ f \colon A \ni a \mapsto b \in B$とかくときもある．}．
\end{definition}

\begin{wrapfigure}{r}{\mnsc{6.0cm}}
  \centering
  \begin{tikzpicture}[mapfig]
    % 集合
    \draw[mapset] (-1.75,0) ellipse [x radius=0.95, y radius=1.25];
    \draw[mapset] ( 1.75,0) ellipse [x radius=0.95, y radius=1.25];
    \node[mapsetname] at (-1.75,1.55) {$A$};
    \node[mapsetname] at ( 1.75,1.55) {$B$};
    \node[mapname] at (0,1.35) {$f$};
    % 要素
    \coordinate (A1) at (-1.75, 0.75);
    \coordinate (A2) at (-1.75, 0);
    \coordinate (A3) at (-1.75,-0.75);
    \coordinate (Ba) at ( 1.75, 0.75);
    \coordinate (Bb) at ( 1.75, 0);
    \coordinate (Bc) at ( 1.75,-0.75);
    \foreach \p in {A1,A2,A3,Ba,Bb} \node[mapelem] at (\p) {};
    \node[mapelemfree] at (Bc) {};% どの元にも対応しない
    \node[maplabL] at (A1) {$1$};
    \node[maplabL] at (A2) {$2$};
    \node[maplabL] at (A3) {$3$};
    \node[maplabR] at (Ba) {$a$};
    \node[maplabR] at (Bb) {$b$};
    \node[maplabR] at (Bc) {$c$};
    % 対応
    \draw[mapacc] (A1) to[bend left=10] (Ba);
    \draw[mapacc] (A2) -- (Bb);
    \draw[mapacc] (A3) to[bend left=12] (Ba);
  \end{tikzpicture}
  \caption{写像 $f \colon A \to B$}
\end{wrapfigure}

$A=\{ 1,2 ,3 \}$， $B=\{ a, b,c \}$とする．$f(1) = a$，$f(2) = b$，$f(3) = a$とすると，この写像$f$は右図のように表現できる．


この例では，$1$と$3$がともに$a$に対応しているが，上記の写像の定義ではこのような場合があってもよい．
また，$A$のどの元にも対応しない元$c$の存在も許容される．

学習が進むと，「任意の$ a_1 , a_2 \in A$に対して，$ a_1 \ne a_2$ならば$f(a_1) \ne f(a_2)$である写像」や
「任意の$ b \in B$に対して，ある$a \in A$が存在して，$b=f(a)$となる写像」を考える場合もあるが，この説明はもう少しあとにしよう．

\phantomsection
\subsection{写像の例}

ここでは直感的な理解を深めるために，具体例を挙げてみよう．

\begin{example}{自動販売機}{自動販売機}
  自動販売機のボタンの集合を$A$，提供される飲み物の集合を$B$とする．
  各ボタン$a \in A$に，そのボタンを押すと出てくる飲み物$f(a) \in B$を対応させる．
  どのボタンにも飲み物がただひとつ定まるから，$f \colon A \to B$は写像である．
\end{example}

\begin{example}{写像の例と写像でない例}{写像の例と写像でない例}
  \begin{itemize}
    \item $f \colon \mathbb{R} \ni x \mapsto x^2 \in \mathbb{R}$は写像である．
          各実数$x$に対して，$x^2$がただひとつ定まるからである．
    \item 各正の実数$x$に「2乗すると$x$になる実数」を対応させる規則は，写像ではない．
          たとえば$x = 4$には$2$と$-2$の2つが対応し，「ただひとつ対応させる」という条件を満たさないからである．
  \end{itemize}
\end{example}

\begin{remark}[「式」と「写像」の違い]
  「多項式$x^2 + x$は関数である」という言い方をすることがあるが，正確には，式$x^2 + x$それ自体は単なる多項式であって，写像ではない．定義域と終域を定め，
  \[
    f \colon \mathbb{R} \ni x \mapsto x^2 + x \in \mathbb{R}
  \]
  のように対応の規則を与えてはじめて写像となる．同じ式を用いても，定義域を$[0, 1]$に変えれば，写像としては別のものである．
  逆に，\exref{自動販売機}のように，ひとつの式では表されない写像も存在する．すなわち，「写像であること」と「式で表されること」は別のことがらである．

  また，関数をグラフでとらえる見方も有用であるが，グラフ（点$(x, f(x))$全体のなす集合）だけでは終域の情報は定まらないことに注意したい．
  同じ式・同じグラフでも，終域の取り方によって写像の性質は変わりうる（この違いは，\deref{全射}のあとの例で実際に現れる）．
\end{remark}


\phantomsection
\subsection{写像の厳密な定義}

\deref{写像}では，写像を「元にただひとつの元を対応させる規則」として定義した．この定義は直感的でわかりやすいが，「規則」という言葉自体は数学的に定義されていない，という弱点がある．そこで本項では，集合の言葉だけを用いた写像の厳密な定義を紹介する．準備として，順序対・直積と，記号$\exists!$を導入しよう．

\phantomsection
\begin{definition}{順序対と直積}{順序対と直積}
  $a \in A$と$b \in B$を，この順序を込めて組にしたもの$(a, b)$を\emph{順序対}\index{じゅんじょつい@順序対}という．順序対については，
  \[
    (a_1, b_1) = (a_2, b_2)
    \lequiv
    (a_1 = a_2) \land (b_1 = b_2)
  \]
  が成り立つものとする\footnote{順序対そのものも，$(a, b) \coloneqq \{ \{a\},\, \{a, b\} \}$のように集合として定義できる（クラトフスキーの定義）．本書では，上の性質を満たすものとして順序対を扱えば十分である．}．順序対全体からなる集合
  \[
    A \times B \coloneqq \{ (a, b) \mid a \in A,\ b \in B \}
  \]
  を$A$と$B$の\emph{直積}\index{ちょくせき@直積}という．
\end{definition}

集合においては$\{1, 2\} = \{2, 1\}$と要素の順序を問わなかったのに対し，順序対では$(1, 2) \ne (2, 1)$である．座標平面の点$(x, y)$は順序対の身近な例であり，実際，座標平面は直積$\mathbb{R} \times \mathbb{R}$とみなすことができる．

\phantomsection
\begin{definition}{記号$\exists!$}{一意存在の記号}
  $P(y)$を述語とする．「$P(y)$を満たす$y \in B$がただひとつ存在する」ことを
  \[
    \lexu{y \in B} \lst P(y)
  \]
  と表す\index{いちいそんざい@一意存在}．これは，
  \[
    \lex{y \in B} \lst \Bigl( P(y) \land \bigl( \lall{z \in B} \lst ( P(z) \limp z = y ) \bigr) \Bigr)
  \]
  の略記である．
\end{definition}

すなわち，$\lexu{y \in B} \lst P(y)$は，
\begin{itemize}
  \item \textbf{存在}：$P(y)$を満たす$y \in B$が（少なくともひとつ）存在する
  \item \textbf{一意性}：$P(z)$を満たす$z \in B$は，その$y$のほかにない
\end{itemize}
というふたつの主張をあわせたものである．\partref{証明}で見たように，「存在」と「一意性」は別々に証明すべきふたつの課題であった．記号$\exists!$は，まさにその両方を一度に主張する記号なのである．

準備が整ったので，写像の厳密な定義を述べよう．

\phantomsection
\begin{definition}{写像（厳密な定義）}{写像の厳密な定義}
  $A$，$B$を集合とする．$A$から$B$への\emph{写像}$f$とは，直積$A \times B$の部分集合$G \subset A \times B$であって，
  \[
    \lall{a \in A} \lsep \lexu{b \in B} \lst (a, b) \in G
  \]
  を満たすもののことである．このとき，各$a \in A$に対して，$(a, b) \in G$を満たすただひとつの$b \in B$を$f(a)$とかく．
\end{definition}

\begin{remark}[ふたつの定義の関係]
  \deref{写像}の「任意の$a \in A$に対して，$b \in B$をただひとつ対応させる」という条件は，厳密な定義では論理式
  \[
    \lall{a \in A} \lsep \lexu{b \in B} \lst (a, b) \in G
  \]
  として表現されている．そして，「規則」という無定義の言葉が，「対応する組$(a, b)$をすべて集めた集合$G$」に置き換えられている点が重要である．たとえば，$f \colon \mathbb{R} \ni x \mapsto x^2 \in \mathbb{R}$に対応する集合は
  \[
    G = \{ (x, x^2) \mid x \in \mathbb{R} \} \subset \mathbb{R} \times \mathbb{R}
  \]
  であり，これはまさに$y = x^2$のグラフである．すなわち，厳密な定義は「写像とは，そのグラフのことである」という見方を与えている\footnote{ただし，集合$G$だけからは終域$B$は復元できない（「式」と「写像」の違いに関する先の注意を参照）．そこで，終域の情報まで込めて，写像を3つ組$(A, B, G)$として定義する流儀もある．}．
\end{remark}

なお，本書では以後も\deref{写像}の直感的な定義に基づいて議論を進めるが，「規則」という言葉が集合の言葉で正当化できることは，記憶にとどめておいてほしい．

\phantomsection
\subsection{恒等写像と制限}

どの集合に対しても考えられる基本的な写像と，写像から新しい写像を作る操作をひとつ導入しておこう．

\phantomsection
\begin{definition}{恒等写像・包含写像・制限}{恒等写像・包含写像・制限}
  \begin{description}
    \item[恒等写像\index{こうとうしゃぞう@恒等写像}]\mbox{}\\
          集合$A$に対して，各$a \in A$に$a$自身を対応させる写像$A \ni a \mapsto a \in A$を，$A$上の恒等写像といい，$\mathrm{id}_A$とかく．
    \item[包含写像\index{ほうがんしゃぞう@包含写像}]\mbox{}\\
          集合$A$，$B$が$A \subset B$を満たすとき，写像$A \ni a \mapsto a \in B$を包含写像という．
    \item[制限\index{せいげん@制限}]\mbox{}\\
          写像$f \colon A \to B$と$A' \subset A$に対して，写像$A' \ni a \mapsto f(a) \in B$を$f$の$A'$への制限といい，$f|_{A'}$とかく．
  \end{description}
\end{definition}

たとえば，$f \colon \mathbb{R} \ni x \mapsto x^2 \in \mathbb{R}$の$[0, \infty)$への制限$f|_{[0, \infty)}$は，$[0, \infty) \ni x \mapsto x^2 \in \mathbb{R}$である．

\begin{remark}[恒等写像と包含写像の違い]
  恒等写像$\mathrm{id}_A$と，$A \subset B$のときの包含写像$A \to B$は，どちらも「各$a$に$a$自身を対応させる」写像であるが，$A \ne B$ならば終域が異なるので，写像としては別のものである．「式」と「写像」の違いに関する注意で述べたように，写像は対応の規則だけでなく，定義域と終域を込めて定まるのである．
  同様に，制限$f|_{A'}$は$f$と同じ規則で定まるが，定義域が異なるので$f$とは別の写像であり，のちに見るように，単射であるかどうかといった性質も変わりうる．
\end{remark}

\phantomsection
\section{像と逆像}

\subsection{像と値域，逆像}

\begin{definition}{像と値域}{像と値域}
  写像$f \colon A \to B$において，集合$A$の部分集合$A' \subset A$に対して，
  \[
    f(A') = \{ f(a) \mid a \in A' \}
  \]
  で定義される集合を，$f$による$A'$の\emph{像}という．

  特に，定義域$A$全体の像
  \[
    f(A) = \{ f(a) \mid a \in A \}
  \]
  を，写像$f$の\emph{値域}という\footnote{「値域」という語は文献によって流儀が異なり，終域$B$の意味で用いられることもある．本書では，$f(A)$を値域とよぶ．}．

  \index{ぞう@像} \index{ちいき@値域}
\end{definition}


\begin{definition}{逆像}{逆像}
  写像$f \colon A \to B$と集合$B' \subset B$に対して，$B'$の\emph{逆像}とは，
  \[
    f^{-1}(B') = \{ a \in A \mid f(a) \in B' \}
  \]
  で定義される集合である．

  \index{ぎゃくぞう@逆像}
\end{definition}

\begin{remark}[逆像と逆写像]
  \deref{逆像}で定義した逆像は「写像」ではなく，あくまで「集合」であることに注意したい．
  逆像と紛らわしい単語に「逆写像」があるが，逆写像は$f$が全単射であるときに定義される写像である（\deref{逆写像}）．その一方で，逆像は$f$が全単射でなくても定義される．
\end{remark}

\begin{figure}[ht]
  \centering
  \begin{minipage}[t]{0.48\linewidth}
    \centering
    \begin{tikzpicture}[mapfig]
      % 集合
      \draw[mapset] (-1.75,0) ellipse [x radius=0.95, y radius=1.25];
      \draw[mapset] ( 1.75,0) ellipse [x radius=0.95, y radius=1.25];
      \node[mapsetname] at (-1.75,-1.55) {$A$};
      \node[mapsetname] at ( 1.75,-1.55) {$B$};
      \node[mapname] at (0,1.35) {$f$};
      % A' = {1, 2} とその像 f(A') = {a, b}
      \draw[mapaccsub] (-1.75,0.375) ellipse [x radius=0.55, y radius=0.85];
      \draw[mapaccsub] ( 1.75,0.125) ellipse [x radius=0.50, y radius=0.85];
      \draw[mapaccleader] (-2.14,0.98) -- (-2.45,1.35)
        node[anchor=south, mapacclab] {$A'$};
      \draw[mapaccleader] ( 2.10,0.72) -- ( 2.45,1.35)
        node[anchor=south, mapacclab] {$f(A')$};
      % 要素
      \coordinate (A1) at (-1.75, 0.75);
      \coordinate (A2) at (-1.75, 0);
      \coordinate (A3) at (-1.75,-0.75);
      \coordinate (Ba) at ( 1.75, 0.5);
      \coordinate (Bb) at ( 1.75,-0.25);
      \foreach \p in {A1,A2,A3,Ba,Bb} \node[mapelem] at (\p) {};
      \node[maplabL] at (A1) {$1$};
      \node[maplabL] at (A2) {$2$};
      \node[maplabL] at (A3) {$3$};
      \node[maplabR] at (Ba) {$a$};
      \node[maplabR] at (Bb) {$b$};
      % 対応（A' の元からの矢印だけを強調）
      \draw[mapacc]   (A1) to[bend left=10] (Ba);
      \draw[mapacc]   (A2) -- (Bb);
      \draw[maparrow] (A3) to[bend left=12] (Ba);
    \end{tikzpicture}
    \caption{$A' = \{1, 2\}$ の像 $f(A') = \{a, b\}$}
  \end{minipage}
  \hfill
  \begin{minipage}[t]{0.48\linewidth}
    \centering
    \begin{tikzpicture}[mapfig]
      % 集合
      \draw[mapset] (-1.75,0) ellipse [x radius=0.95, y radius=1.35];
      \draw[mapset] ( 1.75,0) ellipse [x radius=0.95, y radius=1.35];
      \node[mapsetname] at (-1.75,-1.65) {$A$};
      \node[mapsetname] at ( 1.75,-1.65) {$B$};
      \node[mapname] at (0,1.35) {$f$};
      % B' = {a} とその逆像 f^{-1}(B') = {1, 2, 3}
      \draw[mapaccsub] ( 1.75,0.5) ellipse [x radius=0.45, y radius=0.42];
      \draw[mapaccsub] (-1.75,0.30) ellipse [x radius=0.60, y radius=1.00];
      \draw[mapaccleader] (-2.17,1.01) -- (-2.45,1.45)
        node[anchor=south, mapacclab] {$f^{-1}(B')$};
      \draw[mapaccleader] ( 2.07,0.80) -- ( 2.40,1.45)
        node[anchor=south, mapacclab] {$B'$};
      % 要素
      \coordinate (A1) at (-1.75, 0.95);
      \coordinate (A2) at (-1.75, 0.30);
      \coordinate (A3) at (-1.75,-0.35);
      \coordinate (A4) at (-1.75,-1.00);
      \coordinate (Ba) at ( 1.75, 0.5);
      \coordinate (Bb) at ( 1.75,-0.5);
      \foreach \p in {A1,A2,A3,A4,Ba,Bb} \node[mapelem] at (\p) {};
      \node[maplabL] at (A1) {$1$};
      \node[maplabL] at (A2) {$2$};
      \node[maplabL] at (A3) {$3$};
      \node[maplabL] at (A4) {$4$};
      \node[maplabR] at (Ba) {$a$};
      \node[maplabR] at (Bb) {$b$};
      % 対応（B' に入る矢印だけを強調）
      \draw[mapacc]   (A1) to[bend left=10] (Ba);
      \draw[mapacc]   (A2) -- (Ba);
      \draw[mapacc]   (A3) to[bend right=8] (Ba);
      \draw[maparrow] (A4) to[bend right=12] (Bb);
    \end{tikzpicture}
    \caption{$B' = \{a\}$ の逆像 $f^{-1}(B') = \{1, 2, 3\}$}
  \end{minipage}
\end{figure}



\phantomsection
\section{単射・全射・全単射}

\phantomsection
\subsection{単射\index{たんしゃ@単射}}

\begin{definition}{単射}{単射}
  写像$f \colon A \to B$が\emph{単射}であるとは，任意の$a_1, a_2 \in A$に対して，$a_1 \ne a_2$ならば$f(a_1) \ne f(a_2)$が成り立つことである\footnote{この定義は，対偶を考えると「$f(a_1) =f(a_2)$ならば$a_1=a_2$が成り立つ」とも言い換えられる．}．
\end{definition}

\phantomsection
\begin{example}[
    sidebyside,
    righthand width=\mnsc{6.2cm},
    sidebyside gap=\mnsc{4mm},
  ]{単射の図解}{単射の図解}
  集合$A = \{1, 2, 3\}$，$B = \{a, b, c, d\}$と写像
  \[
    f(1) = a, \quad f(2) = b, \quad f(3) = c
  \]
  を考える．
  $A$の異なる元は$B$の異なる元にうつるから，$f$は単射である．
  しかし，この$f$は全射でない．

  \tcblower

  \centering
  \begin{tikzpicture}[mapfig]
    % 集合
    \draw[mapset] (-1.75,0) ellipse [x radius=0.95, y radius=1.35];
    \draw[mapset] ( 1.75,0) ellipse [x radius=0.95, y radius=1.35];
    \node[mapsetname] at (-1.75,1.65) {$A$};
    \node[mapsetname] at ( 1.75,1.65) {$B$};
    \node[mapname] at (0,1.4) {$f$};
    % 要素
    \coordinate (A1) at (-1.75, 0.75);
    \coordinate (A2) at (-1.75, 0);
    \coordinate (A3) at (-1.75,-0.75);
    \coordinate (Ba) at ( 1.75, 0.95);
    \coordinate (Bb) at ( 1.75, 0.32);
    \coordinate (Bc) at ( 1.75,-0.32);
    \coordinate (Bd) at ( 1.75,-0.95);
    \foreach \p in {A1,A2,A3,Ba,Bb,Bc} \node[mapelem] at (\p) {};
    \node[mapelemfree] at (Bd) {};% 像に入らない（全射でない）
    \node[maplabL] at (A1) {$1$};
    \node[maplabL] at (A2) {$2$};
    \node[maplabL] at (A3) {$3$};
    \node[maplabR] at (Ba) {$a$};
    \node[maplabR] at (Bb) {$b$};
    \node[maplabR] at (Bc) {$c$};
    \node[maplabR] at (Bd) {$d$};
    % 対応
    \draw[mapacc] (A1) to[bend left=8] (Ba);
    \draw[mapacc] (A2) to[bend left=6] (Bb);
    \draw[mapacc] (A3) to[bend right=6] (Bc);
  \end{tikzpicture}
\end{example}

\begin{example}{単射の例}{単射の例}
  \[
    f \colon \mathbb{N} \ni n \mapsto 2n \in \mathbb{N}
  \]
  は単射である．実際，$2n_1 = 2n_2$ならば$n_1 = n_2$である．
\end{example}

\begin{example}{電話番号と電話機}{電話番号と電話機}
  電話番号の集合を$A$，電話機の集合を$B$とし，各電話番号$a \in A$に対応する電話機$f(a) \in B$を対応させる写像$f \colon A \to B$を考える．
  異なる電話番号が同じ電話機を指すことはないから，$f$は単射である．
\end{example}

\subsection{全射\index{ぜんしゃ@全射}}

\begin{definition}{全射}{全射}
  写像$f \colon A \to B$が\emph{全射}であるとは，任意の$b \in B$に対して，ある$a \in A$が存在して，$f(a) = b$が成り立つことである．
\end{definition}

\phantomsection
\begin{example}[
    sidebyside,
    righthand width=\mnsc{6.2cm},
    sidebyside gap=\mnsc{4mm},
  ]{全射の図解}{全射の図解}
  集合$A = \{1, 2, 3\}$，$B = \{a, b\}$と写像
  \[
    f(1) = a, \quad f(2) = a, \quad f(3) = b
  \]
  を考える．
  $B$のどの元も$A$のある元の像になっているから，$f$は全射である．
  しかし，$f(1) = f(2) = a$であるから，$f$は単射でない．

  \tcblower

  \centering
  \begin{tikzpicture}[mapfig]
    % 集合
    \draw[mapset] (-1.75,0) ellipse [x radius=0.95, y radius=1.25];
    \draw[mapset] ( 1.75,0) ellipse [x radius=0.95, y radius=1.25];
    \node[mapsetname] at (-1.75,1.55) {$A$};
    \node[mapsetname] at ( 1.75,1.55) {$B$};
    \node[mapname] at (0,1.35) {$f$};
    % 要素
    \coordinate (A1) at (-1.75, 0.75);
    \coordinate (A2) at (-1.75, 0);
    \coordinate (A3) at (-1.75,-0.75);
    \coordinate (Ba) at ( 1.75, 0.5);
    \coordinate (Bb) at ( 1.75,-0.5);
    \foreach \p in {A1,A2,A3,Ba,Bb} \node[mapelem] at (\p) {};
    \node[maplabL] at (A1) {$1$};
    \node[maplabL] at (A2) {$2$};
    \node[maplabL] at (A3) {$3$};
    \node[maplabR] at (Ba) {$a$};
    \node[maplabR] at (Bb) {$b$};
    % 対応
    \draw[mapacc] (A1) to[bend left=10] (Ba);
    \draw[mapacc] (A2) -- (Ba);
    \draw[mapacc] (A3) to[bend right=8] (Bb);
  \end{tikzpicture}
\end{example}

\begin{example}{全射の例}{全射の例}
  \[
    f \colon \mathbb{R} \ni x \mapsto x^2 \in [0, \infty)
  \]
  は全射である．実際，$b \in [0, \infty)$を任意にとるとき，$x = \sqrt{b}$とおけば$f(x) = b$となる．
  一方，$f(-1) = f(1) = 1$であるから，$f$は単射でない．

  なお，\exref{写像の例と写像でない例}の$x \mapsto x^2$とは式が同じで終域だけが異なるが，終域を$\mathbb{R}$とすると全射でなくなる．写像の性質は，式だけでなく定義域・終域を込めて定まるのである．
\end{example}

\begin{example}{分類コードと本}{分類コードと本}
  図書館にある本の集合を$A$，使用されている分類コードの集合を$B$とし，各本$a \in A$にその分類コード$f(a) \in B$を対応させる写像$f \colon A \to B$を考える．
  $B$は使用されているコード全体であるから，どのコード$b \in B$に対しても，それが割り当てられた本$a \in A$が存在する．よって$f$は全射である．
  一方，同じ分類コードを持つ本は複数ありうるから，$f$は一般に単射でない．
\end{example}


\subsection{全単射\index{ぜんたんしゃ@全単射}}

\begin{definition}{全単射}{全単射}
  写像$f \colon A \to B$が\emph{全単射}であるとは，$f$が単射かつ全射となることである．
\end{definition}

\phantomsection
\begin{example}[
    sidebyside,
    righthand width=\mnsc{6.2cm},
    sidebyside gap=\mnsc{4mm},
  ]{全単射の図解}{全単射の図解}
  集合$A = \{1, 2, 3\}$，$B = \{a, b, c\}$と写像
  \[
    f(1) = c, \quad f(2) = b, \quad f(3) = a
  \]
  を考える．
  この$f$は単射かつ全射，すなわち全単射である．
  $A$の元と$B$の元が，$f$によって過不足なく1対1に対応している．

  \tcblower

  \centering
  \begin{tikzpicture}[mapfig]
    % 集合
    \draw[mapset] (-1.75,0) ellipse [x radius=0.95, y radius=1.25];
    \draw[mapset] ( 1.75,0) ellipse [x radius=0.95, y radius=1.25];
    \node[mapsetname] at (-1.75,1.55) {$A$};
    \node[mapsetname] at ( 1.75,1.55) {$B$};
    \node[mapname] at (0,1.35) {$f$};
    % 要素
    \coordinate (A1) at (-1.75, 0.75);
    \coordinate (A2) at (-1.75, 0);
    \coordinate (A3) at (-1.75,-0.75);
    \coordinate (Ba) at ( 1.75, 0.75);
    \coordinate (Bb) at ( 1.75, 0);
    \coordinate (Bc) at ( 1.75,-0.75);
    \foreach \p in {A1,A2,A3,Ba,Bb,Bc} \node[mapelem] at (\p) {};
    \node[maplabL] at (A1) {$1$};
    \node[maplabL] at (A2) {$2$};
    \node[maplabL] at (A3) {$3$};
    \node[maplabR] at (Ba) {$a$};
    \node[maplabR] at (Bb) {$b$};
    \node[maplabR] at (Bc) {$c$};
    % 対応（1→c, 2→b, 3→a）
    \draw[mapacc] (A1) -- (Bc);
    \draw[mapacc] (A2) -- (Bb);
    \draw[mapacc] (A3) -- (Ba);
  \end{tikzpicture}
\end{example}

\begin{example}{学籍番号と学生}{学籍番号と学生}
  あるクラスの学生の集合を$A$，そのクラスで使われている学籍番号の集合を$B$とし，各学生$a \in A$にその学籍番号$f(a) \in B$を対応させる写像$f \colon A \to B$を考える．
  異なる学生が同じ学籍番号を持つことはないから$f$は単射であり，どの学籍番号もある学生に割り当てられているから$f$は全射である．よって$f$は全単射である．
\end{example}

\begin{example}{単射・全射・全単射}{単射・全射・全単射}
  \begin{enumerate}[(A)]
    \item
          \[
            f \colon [0,\infty) \ni x \mapsto x^2  \in \mathbb{R}
          \]
          は単射であるが全射でない．
    \item
          \[
            g \colon \mathbb{R} \ni x \mapsto x^3 -3x^2 -9 \in \mathbb{R}
          \]
          は全射であるが単射でない．
    \item
          \[
            h  \colon \mathbb{R} \ni x \mapsto  2 \sin x \in  \mathbb{R}
          \]
          は全射でも単射でもない．しかし，
          \[
            h'  \colon [-\pi/2 , \pi/2 ]  \ni x \mapsto  2 \sin  x  \in [-2,2]
          \]
          は全単射である．$h'$は，$h$を$[-\pi/2, \pi/2]$に制限し，さらに終域を$[-2, 2]$に取り替えたものである．
  \end{enumerate}
\end{example}

\phantomsection
\section{像・逆像と集合演算}

\phantomsection
\subsection{像・逆像の基本公式}

像と逆像は，\partref{集合}で学んだ集合の演算（$\cup$，$\cap$，$\setminus$）とどのように関わるのだろうか．まず，逆像については，集合の演算がすべてそのまま保たれる．

\phantomsection
\begin{prop}{逆像と集合演算}{逆像と集合演算}
  写像$f \colon A \to B$と$B_1, B_2 \subset B$に対して，次が成り立つ：
  \begin{align*}
    f^{-1}(B_1 \cup B_2) & = f^{-1}(B_1) \cup f^{-1}(B_2), \\
    f^{-1}(B_1 \cap B_2) & = f^{-1}(B_1) \cap f^{-1}(B_2), \\
    f^{-1}(B \setminus B_1) & = A \setminus f^{-1}(B_1)
  \end{align*}
\end{prop}

\begin{tleftbar}
  \begin{proof}
    2番目の等式を示す．$a \in A$に対して，
    \begin{align*}
      a \in f^{-1}(B_1 \cap B_2)
      & \lequiv f(a) \in B_1 \cap B_2 \\
      & \lequiv \bigl( f(a) \in B_1 \bigr) \land \bigl( f(a) \in B_2 \bigr) \\
      & \lequiv \bigl( a \in f^{-1}(B_1) \bigr) \land \bigl( a \in f^{-1}(B_2) \bigr) \\
      & \lequiv a \in f^{-1}(B_1) \cap f^{-1}(B_2)
    \end{align*}
    である．よって両者は等しい．残りの等式も，同様に逆像の定義に沿って書き換えれば示される．
  \end{proof}
\end{tleftbar}

一方，像については事情が異なる．和集合は保たれるが，共通部分については包含関係しか成り立たない．

\phantomsection
\begin{prop}{像と集合演算}{像と集合演算}
  写像$f \colon A \to B$と$A_1, A_2 \subset A$に対して，次が成り立つ：
  \begin{align*}
    f(A_1 \cup A_2) & = f(A_1) \cup f(A_2), \\
    f(A_1 \cap A_2) & \subset f(A_1) \cap f(A_2)
  \end{align*}
  ただし，2番目の包含は，一般には等号にならない．
\end{prop}

\begin{tleftbar}
  \begin{proof}
    2番目の包含を示す．$b \in f(A_1 \cap A_2)$とすると，ある$a \in A_1 \cap A_2$が存在して$b = f(a)$である．$a \in A_1$かつ$a \in A_2$なので，$b = f(a)$は$f(A_1)$にも$f(A_2)$にも属する．よって$b \in f(A_1) \cap f(A_2)$である．

    次に，等号が成り立たない例を挙げる．$f \colon \mathbb{R} \ni x \mapsto x^2 \in \mathbb{R}$，$A_1 = \{1\}$，$A_2 = \{-1\}$とすると，$A_1 \cap A_2 = \varnothing$より
    \[
      f(A_1 \cap A_2) = f(\varnothing) = \varnothing
    \]
    である一方，
    \[
      f(A_1) \cap f(A_2) = \{1\} \cap \{1\} = \{1\} \ne \varnothing
    \]
    である．なお，1番目の等式も，像の定義に沿って書き換えれば同様に示される．これが証明すべきことであった．
  \end{proof}
\end{tleftbar}

なぜ等号が崩れるのかというと，この$f$が異なる元$1$と$-1$を同じ元$1$にうつしてしまうからである．実際，$f$が単射であれば，この現象は起こらない．

\phantomsection
\begin{prop}{単射と像の共通部分}{単射と像の共通部分}
  写像$f \colon A \to B$が単射ならば，任意の$A_1, A_2 \subset A$に対して，
  \[
    f(A_1 \cap A_2) = f(A_1) \cap f(A_2)
  \]
  が成り立つ．
\end{prop}

\begin{tleftbar}
  \begin{proof}
    包含$\subset$は\prref{像と集合演算}で示したので，$\supset$を示す．$b \in f(A_1) \cap f(A_2)$とすると，ある$a_1 \in A_1$と$a_2 \in A_2$が存在して，
    \[
      b = f(a_1) = f(a_2)
    \]
    である．ここで$f$は単射であるから，この等式から$a_1 = a_2$が従う．よって$a_1$は$A_1$と$A_2$の両方に属し，$a_1 \in A_1 \cap A_2$である．したがって$b = f(a_1)$は$f(A_1 \cap A_2)$に属する．これが証明すべきことであった．
  \end{proof}
\end{tleftbar}

この命題は，単射という性質がなぜ重要なのかを示すひとつの例である．逆像がつねに集合演算と相性がよいのに対して，像のふるまいは写像の性質（単射かどうか）に応じて変わるのである．

\phantomsection
\subsection{像と逆像の往復}

部分集合$A' \subset A$の像をとってから逆像をとると，$A'$に戻るのだろうか．また，$B' \subset B$の逆像をとってから像をとるとどうなるだろうか．

\phantomsection
\begin{prop}{像と逆像の往復}{像と逆像の往復}
  写像$f \colon A \to B$と$A' \subset A$，$B' \subset B$に対して，次が成り立つ：
  \begin{enumerate}[label=(\roman*),leftmargin=3.2em]
    \item $A' \subset f^{-1}(f(A'))$であり，$f$が単射ならば等号が成り立つ．
    \item $f(f^{-1}(B')) \subset B'$であり，$f$が全射ならば等号が成り立つ．
  \end{enumerate}
\end{prop}

\begin{tleftbar}
  \begin{proof}
    (i)について，$a \in A'$とすると$f(a) \in f(A')$であるから，逆像の定義より$a \in f^{-1}(f(A'))$である．よって$A' \subset f^{-1}(f(A'))$である．
    次に$f$が単射であるとし，$a \in f^{-1}(f(A'))$とする．$f(a) \in f(A')$より，ある$a' \in A'$が存在して$f(a) = f(a')$である．$f$は単射であるから$a = a'$であり，$a \in A'$である．よって等号が成り立つ．

    (ii)について，$b \in f(f^{-1}(B'))$とすると，ある$a \in f^{-1}(B')$が存在して$b = f(a)$である．$a \in f^{-1}(B')$より$f(a) \in B'$であるから，$b \in B'$である．よって$f(f^{-1}(B')) \subset B'$である．
    次に$f$が全射であるとし，$b \in B'$とする．$f$は全射であるから，ある$a \in A$が存在して$f(a) = b$である．$f(a) = b \in B'$より$a \in f^{-1}(B')$であるから，$b = f(a) \in f(f^{-1}(B'))$である．よって等号が成り立つ．これが証明すべきことであった．
  \end{proof}
\end{tleftbar}

\phantomsection
\begin{example}{像と逆像の往復で等号が成り立たない例}{像と逆像の往復で等号が成り立たない例}
  $f \colon \mathbb{R} \ni x \mapsto x^2 \in \mathbb{R}$とする．この$f$は単射でも全射でもない．
  \begin{enumerate}[(A)]
    \item $A' = [0, 1]$とすると，$f(A') = [0, 1]$であり，$f^{-1}(f(A')) = [-1, 1] \supsetneq A'$である．
    \item $B' = [-1, 1]$とすると，$f^{-1}(B') = [-1, 1]$であり，$f(f^{-1}(B')) = [0, 1] \subsetneq B'$である．
  \end{enumerate}
\end{example}

\begin{remark}[逆像は像の「逆」ではない]
  \exref{像と逆像の往復で等号が成り立たない例}からわかるように，像をとる操作と逆像をとる操作は，一般には互いに打ち消しあわない．(A)で$A'$が大きくなったのは，$f$が単射でないために，$A'$の外の元$-1$なども同じ像にうつるからである．(B)で$B'$が小さくなったのは，$f$が全射でないために，負の数が$f$の像に現れないからである．
  記号$f^{-1}$は「逆」を連想させるが，逆像$f^{-1}(B')$は，$f$でうつして$B'$に入る元を集めた集合にすぎないことに注意しよう．
\end{remark}

\phantomsection
\section{写像の合成}

\phantomsection
\subsection{合成写像}

写像を$2$つ続けて施すと，新しい写像が得られる．

\phantomsection
\begin{definition}[
    sidebyside,
    righthand width=\mnsc{6.2cm},
    sidebyside gap=\mnsc{4mm},
  ]{合成写像}{合成写像}
  写像$h \colon X \to Y$と$g \colon Y \to Z$に対して，各$x \in X$に$g(h(x)) \in Z$を対応させる写像を$h$と$g$の\emph{合成写像}\index{ごうせいしゃぞう@合成写像}といい，$g \circ h \colon X \to Z$とかく：
  \[
    (g \circ h)(x) \coloneqq g(h(x)) \qquad (x \in X)
  \]

  \tcblower

  \centering
  \begin{tikzpicture}[mapfig]
    % 集合
    \foreach \c/\n in {-2.05/X, 0/Y, 2.05/Z} {
      \draw[mapset] (\c,0) ellipse [x radius=0.62, y radius=1.05];
      \node[mapsetname] at (\c,-1.35) {$\n$};
    }
    % 元
    \coordinate (x1) at (-2.05, 0.45);
    \coordinate (x2) at (-2.05,-0.45);
    \coordinate (y1) at ( 0, 0.6);
    \coordinate (y2) at ( 0, 0);
    \coordinate (y3) at ( 0,-0.6);
    \coordinate (z1) at ( 2.05, 0.4);
    \coordinate (z2) at ( 2.05,-0.4);
    \foreach \p in {x1,x2,y1,y3,z1,z2} \node[mapelem] at (\p) {};
    \node[mapelemfree] at (y2) {};
    % h と g（灰）
    \draw[maparrow] (x1) -- (y1);
    \draw[maparrow] (x2) -- (y3);
    \draw[maparrow] (y1) -- (z1);
    \draw[maparrow] (y2) -- (z1);
    \draw[maparrow] (y3) -- (z2);
    \node[mapsetname] at (-1.02,0.98) {$h$};
    \node[mapsetname] at ( 1.02,0.98) {$g$};
    % 合成 g∘h（マゼンタ）：X から Z へ直接
    \draw[mapacc] (-2.05,1.12) to[bend left=28]
      node[midway, above, mapname] {$g \circ h$} (2.05,1.12);
  \end{tikzpicture}
\end{definition}

「先に$h$，次に$g$」を施す写像が$g \circ h$であり，記号の順序が施す順序と逆になることに注意しよう．

 
また，高校数学で学んだ「合成関数」は，合成写像の特別な場合である．身近な合成写像の例をいくつか見ておこう．
 
\begin{example}{加法の合成}{加法の合成}
  「$2$を足す」写像と「$5$を足す」写像
  \[
    h \colon \mathbb{R} \ni x \mapsto x + 2 \in \mathbb{R}, \qquad
    g \colon \mathbb{R} \ni x \mapsto x + 5 \in \mathbb{R}
  \]
  を考える．合成写像$g \circ h$は，
  \[
    (g \circ h)(x) = g(h(x)) = g(x + 2) = (x + 2) + 5 = x + 7
  \]
  であり，たとえば$(g \circ h)(3) = (3 + 2) + 5 = 10$である．
  すなわち，$3 + 2 + 5$という計算は，$3$に写像$h$，$g$を順に施したものと見ることができる．
  この例では$(h \circ g)(x) = (x + 5) + 2 = x + 7$でもあり，$g \circ h = h \circ g$が成り立つ．
\end{example}
 
\begin{example}{数列の合成}{数列の合成}
  数列$a_1, a_2, a_3, \ldots$は，番号$n$に第$n$項$a_n$を対応させる写像
  \[
    a \colon \mathbb{N} \ni n \mapsto a_n \in \mathbb{R}
  \]
  とみなすことができる\footnote{本書では$\mathbb{N} = \{1, 2, 3, \ldots\}$とする．}．たとえば$a_n = n^2$とし，$k \colon \mathbb{N} \ni n \mapsto 2n \in \mathbb{N}$とすると，
  \[
    (a \circ k)(n) = a(k(n)) = a_{2n} = (2n)^2 = 4n^2
  \]
  である．すなわち，$a \circ k$は元の数列から偶数番目の項$a_2, a_4, a_6, \ldots$だけを取り出した数列$4, 16, 36, \ldots$である．
  一方，$g \colon \mathbb{R} \ni x \mapsto 2x + 1 \in \mathbb{R}$との合成$g \circ a$は，
  \[
    (g \circ a)(n) = g(a_n) = 2n^2 + 1
  \]
  であり，各項を$2x + 1$で変換した数列$3, 9, 19, \ldots$である．
\end{example}
 
\begin{example}{関数の合成}{関数の合成}
  $f \colon \mathbb{R} \ni x \mapsto \sin x \in \mathbb{R}$，$g \colon \mathbb{R} \ni x \mapsto x^2 \in \mathbb{R}$とすると，
  \[
    (g \circ f)(x) = (\sin x)^2 = \sin^2 x, \qquad
    (f \circ g)(x) = \sin(x^2)
  \]
  である．たとえば$x = \sqrt{\pi}$とすると，$(f \circ g)(x) = \sin \pi = 0$であるのに対し，$(g \circ f)(x) = \sin^2 \sqrt{\pi}$は$0$でない．よってこれらは異なる関数であり，$g \circ f \ne f \circ g$である．
  \exref{加法の合成}とは対照的に，合成の順序を入れ替えると一般には別の写像になる．
\end{example}

このように合成の交換法則は一般には成り立たないが，結合律は成り立つ．

\phantomsection
\begin{prop}{合成の結合律と恒等写像}{合成の結合律と恒等写像}
  写像$f \colon A \to B$，$g \colon B \to C$，$h \colon C \to D$に対して，次が成り立つ：
  \begin{enumerate}[label=(\roman*),leftmargin=3.2em]
    \item $(h \circ g) \circ f = h \circ (g \circ f)$である．
    \item $f \circ \mathrm{id}_A = f$かつ$\mathrm{id}_B \circ f = f$である．
  \end{enumerate}
\end{prop}

\begin{tleftbar}
  \begin{proof}
    (i)について，両辺はともに$A$から$D$への写像であり，任意の$a \in A$に対して，
    \[
      \bigl( (h \circ g) \circ f \bigr)(a) = (h \circ g)(f(a)) = h(g(f(a))),
      \qquad
      \bigl( h \circ (g \circ f) \bigr)(a) = h((g \circ f)(a)) = h(g(f(a)))
    \]
    であるから，両者は等しい．

    (ii)について，任意の$a \in A$に対して$(f \circ \mathrm{id}_A)(a) = f(a)$，$(\mathrm{id}_B \circ f)(a) = f(a)$である．これが証明すべきことであった．
  \end{proof}
\end{tleftbar}

結合律が成り立つので，$3$つ以上の写像の合成は，括弧を省略して$h \circ g \circ f$とかいてよい．

単射・全射という性質は，合成とどのように関わるだろうか．

\phantomsection
\begin{prop}{単射・全射と合成}{単射・全射と合成}
  写像$f \colon A \to B$，$g \colon B \to C$に対して，次が成り立つ：
  \begin{enumerate}[label=(\roman*),leftmargin=3.2em]
    \item $f$と$g$がともに単射ならば，$g \circ f$は単射である．
    \item $f$と$g$がともに全射ならば，$g \circ f$は全射である．
    \item $g \circ f$が単射ならば，$f$は単射である．
    \item $g \circ f$が全射ならば，$g$は全射である．
  \end{enumerate}
\end{prop}

\begin{tleftbar}
  \begin{proof}
    (i)について，$a_1, a_2 \in A$が$g(f(a_1)) = g(f(a_2))$を満たすとする．$g$は単射であるから$f(a_1) = f(a_2)$であり，$f$は単射であるから$a_1 = a_2$である．

    (ii)について，$c \in C$を任意にとる．$g$は全射であるから，ある$b \in B$が存在して$g(b) = c$である．さらに$f$は全射であるから，ある$a \in A$が存在して$f(a) = b$である．このとき$(g \circ f)(a) = c$である．

    (iii)について，$a_1, a_2 \in A$が$f(a_1) = f(a_2)$を満たすとする．両辺を$g$でうつせば$(g \circ f)(a_1) = (g \circ f)(a_2)$であり，$g \circ f$は単射であるから$a_1 = a_2$である．

    (iv)について，$c \in C$を任意にとる．$g \circ f$は全射であるから，ある$a \in A$が存在して$g(f(a)) = c$である．$b \coloneqq f(a) \in B$とおけば$g(b) = c$である．これが証明すべきことであった．
  \end{proof}
\end{tleftbar}

\begin{remark}[(iii)，(iv)で結論されるのはどちらか]
  \prref{単射・全射と合成}の(iii)で単射と結論されるのは先に施す$f$であり，(iv)で全射と結論されるのは後に施す$g$である．逆の組合せは一般には成り立たない．
  たとえば，$A = \{1\}$，$B = \{1, 2\}$，$C = \{1\}$とし，$f(1) = 1$，$g(1) = g(2) = 1$とすると，$g \circ f = \mathrm{id}_{\{1\}}$は単射かつ全射であるが，$g$は単射でなく，$f$は全射でない．
  「$g \circ f$が単射ならば，$A$の異なる元は最後まで区別されるのだから，最初の$f$の段階で区別されていなければならない」「$g \circ f$が全射ならば，$C$のどの元にも届くのだから，最後の$g$の段階で届いていなければならない」と考えると覚えやすい．
\end{remark}

\phantomsection
\subsection{逆写像\index{ぎゃくしゃぞう@逆写像}}

写像$f \colon A \to B$が全単射であるとき，どの$b \in B$に対しても，$f(a) = b$を満たす$a \in A$は，全射性より存在し，単射性よりただひとつである．このことを命題の形で述べておこう．

\phantomsection
\begin{prop}{全単射の言い換え}{全単射の言い換え}
  写像$f \colon A \to B$について，$f$が全単射であることと，
  \[
    \lall{b \in B} \lsep \lexu{a \in A} \lst f(a) = b
  \]
  が成り立つことは同値である．
\end{prop}

\begin{tleftbar}
  \begin{proof}
    \deref{一意存在の記号}より，上の論理式は，任意の$b \in B$について次の$2$つが成り立つことを意味する：
    \begin{itemize}
      \item 存在：$f(a) = b$を満たす$a \in A$が存在する．
      \item 一意性：$f(a) = b$かつ$f(a') = b$ならば$a = a'$である．
    \end{itemize}
    前者がすべての$b \in B$について成り立つことは，$f$が全射であることにほかならない．後者がすべての$b \in B$について成り立つことは，「$f(a) = f(a')$ならば$a = a'$」，すなわち$f$が単射であることにほかならない．これが証明すべきことであった．
  \end{proof}
\end{tleftbar}

\prref{全単射の言い換え}の論理式は，写像の厳密な定義（\deref{写像の厳密な定義}）の条件で$A$と$B$の役割を入れ替えた形をしている．したがって，全単射$f$に対しては，「逆向きの対応」がふたたび写像を定める．

\phantomsection
\begin{definition}[
    sidebyside,
    righthand width=\mnsc{6.2cm},
    sidebyside gap=\mnsc{4mm},
  ]{逆写像}{逆写像}
  写像$f \colon A \to B$が全単射であるとき，各$b \in B$に対して，$f(a) = b$を満たすただひとつの$a \in A$を対応させる写像を，$f$の\emph{逆写像}といい，$f^{-1} \colon B \to A$とかく．すなわち，$a \in A$，$b \in B$に対して
  \[
    f^{-1}(b) = a \lequiv f(a) = b
  \]
  である．

  \tcblower

  \centering
  \begin{tikzpicture}[mapfig]
    % 集合
    \draw[mapset] (-1.75,0) ellipse [x radius=0.75, y radius=1.3];
    \draw[mapset] ( 1.75,0) ellipse [x radius=0.75, y radius=1.3];
    \node[mapsetname] at (-1.75,-1.6) {$A$};
    \node[mapsetname] at ( 1.75,-1.6) {$B$};
    % 元
    \coordinate (A1) at (-1.75, 0.8);
    \coordinate (A2) at (-1.75, 0);
    \coordinate (A3) at (-1.75,-0.8);
    \coordinate (Ba) at ( 1.75, 0.8);
    \coordinate (Bb) at ( 1.75, 0);
    \coordinate (Bc) at ( 1.75,-0.8);
    \foreach \p in {A1,A2,A3,Ba,Bb,Bc} \node[mapelem] at (\p) {};
    \node[maplabL] at (A1) {$1$};
    \node[maplabL] at (A2) {$2$};
    \node[maplabL] at (A3) {$3$};
    \node[maplabR] at (Ba) {$a$};
    \node[maplabR] at (Bb) {$b$};
    \node[maplabR] at (Bc) {$c$};
    % f（灰）と f^{-1}（マゼンタ）：同じ対応を逆向きにたどる
    \foreach \a/\b in {A1/Ba, A2/Bb, A3/Bc} {
      \draw[maparrow] (\a) to[bend left=10] (\b);
      \draw[mapacc]   (\b) to[bend left=10] (\a);
    }
    \node[mapsetname] at (0,1.3) {$f$};
    \node[mapname]    at (0,-1.3) {$f^{-1}$};
  \end{tikzpicture}
\end{definition}

\phantomsection
\begin{example}{逆写像の例}{逆写像の例}
  \begin{enumerate}[(A)]
    \item $f \colon \mathbb{R} \ni x \mapsto 2x + 1 \in \mathbb{R}$は全単射であり，$y = 2x + 1$を$x$について解くと$x = (y - 1)/2$であるから，
          \[
            f^{-1} \colon \mathbb{R} \ni y \mapsto \frac{y - 1}{2} \in \mathbb{R}
          \]
          である．
    \item \exref{単射・全射・全単射}の全単射$h' \colon [-\pi/2, \pi/2] \ni x \mapsto 2 \sin x \in [-2, 2]$の逆写像は，逆正弦関数$\arcsin$を用いて$[-2, 2] \ni y \mapsto \arcsin (y/2) \in [-\pi/2, \pi/2]$と表される．
          $h \colon \mathbb{R} \to \mathbb{R}$のままでは全単射でないので逆写像は存在しないが，定義域を制限し終域を取り替えることで，逆写像を考えられるようになるのである．
  \end{enumerate}
\end{example}

逆写像は，合成と恒等写像を用いて特徴づけることができる．

\phantomsection
\begin{theorem}{全単射と逆写像}{全単射と逆写像}
  写像$f \colon A \to B$について，次の$2$つは同値である：
  \begin{enumerate}[label=(\roman*),leftmargin=3.2em]
    \item $f$は全単射である．
    \item ある写像$g \colon B \to A$が存在して，$g \circ f = \mathrm{id}_A$かつ$f \circ g = \mathrm{id}_B$が成り立つ．
  \end{enumerate}
  さらに，このとき(ii)の$g$はただひとつであり，$g = f^{-1}$である．
\end{theorem}

\begin{tleftbar}
  \begin{proof}
    (i)$\limp$(ii)：$g = f^{-1}$とすればよい．実際，$a \in A$に対して$b = f(a)$とおくと，\deref{逆写像}より$f^{-1}(b) = a$であるから$(f^{-1} \circ f)(a) = a$である．また，$b \in B$に対して$a = f^{-1}(b)$とおくと，$f(a) = b$であるから$(f \circ f^{-1})(b) = b$である．

    (ii)$\limp$(i)：$\mathrm{id}_A$は単射であるから，$g \circ f = \mathrm{id}_A$と\prref{単射・全射と合成}(iii)より$f$は単射である．また，$\mathrm{id}_B$は全射であるから，$f \circ g = \mathrm{id}_B$と\prref{単射・全射と合成}(iv)より$f$は全射である．

    最後に一意性を示す．$g, g' \colon B \to A$がともに(ii)を満たすとすると，\prref{合成の結合律と恒等写像}より
    \[
      g = g \circ \mathrm{id}_B = g \circ (f \circ g') = (g \circ f) \circ g' = \mathrm{id}_A \circ g' = g'
    \]
    である．(i)$\limp$(ii)の証明より$f^{-1}$は(ii)を満たすから，$g = f^{-1}$である．これが証明すべきことであった．
  \end{proof}
\end{tleftbar}

一意性の証明は，\prref{逆行列の一意性}の証明$B = BE = B(AC) = (BA)C = EC = C$とまったく同じ形をしていることに注意しよう．\thref{全単射と逆写像}より，$f$が全単射ならば$f^{-1}$も全単射であり，$(f^{-1})^{-1} = f$である．

\phantomsection
\begin{cor}{合成写像の逆写像}{合成写像の逆写像}
  写像$f \colon A \to B$，$g \colon B \to C$がともに全単射ならば，$g \circ f$も全単射であり，
  \[
    (g \circ f)^{-1} = f^{-1} \circ g^{-1}
  \]
  が成り立つ．
\end{cor}

\begin{tleftbar}
  \begin{proof}
    \prref{合成の結合律と恒等写像}を用いると，
    \[
      (f^{-1} \circ g^{-1}) \circ (g \circ f)
      = f^{-1} \circ (g^{-1} \circ g) \circ f
      = f^{-1} \circ \mathrm{id}_B \circ f
      = f^{-1} \circ f
      = \mathrm{id}_A
    \]
    であり，同様に$(g \circ f) \circ (f^{-1} \circ g^{-1}) = \mathrm{id}_C$である．よって\thref{全単射と逆写像}より，$g \circ f$は全単射であり，その逆写像は$f^{-1} \circ g^{-1}$である．これが証明すべきことであった．
  \end{proof}
\end{tleftbar}

逆写像をとると合成の順序が入れ替わることに注意しよう．「先に$f$，次に$g$」を施したものを元に戻すには，「先に$g$を戻し，次に$f$を戻す」必要があるからである．線型代数学で学ぶ行列の積の逆行列の公式$(AB)^{-1} = B^{-1} A^{-1}$も，同じ構造をもっている．

\begin{remark}[記号$f^{-1}$のふたつの意味]
  記号$f^{-1}$は，逆像$f^{-1}(B')$（\deref{逆像}）と逆写像$f^{-1}(b)$（\deref{逆写像}）の両方に用いられる．逆像は任意の写像に対して定義される「$B$の部分集合から$A$の部分集合を作る操作」であり，逆写像は全単射に対してのみ定義される「$B$の元を$A$の元にうつす写像」である．
  $f$が全単射のときには，$B' \subset B$の$f$による逆像と$f^{-1}$による像は一致する：
  \[
    \{ a \in A \mid f(a) \in B' \} = \{ f^{-1}(b) \mid b \in B' \}
  \]
  よって，この場合にはどちらの意味に読んでも$f^{-1}(B')$は同じ集合を表し，記号の衝突は起こらない．
\end{remark}

\phantomsection
\subsection{well-defined}
\label{subsec:well-defined}

\partref{用語の定義}では，「対象の表し方によらず値が定まること」をwell-definedとよんだ．写像の合成が使えるようになったいま，この概念をもう少し正確に述べ直しておこう．

次のような状況を考える．写像$f \colon X \to Z$と全射$h \colon X \to Y$が与えられていて，$Y$の元$y$に対して「$h(x) = y$を満たす$x$をひとつ選び，その$f(x)$を$y$の行き先とする」という規則で，写像$g \colon Y \to Z$をつくりたい．$h$は全射だから，どの$y \in Y$に対しても$h(x) = y$を満たす$x \in X$は少なくともひとつ存在する．しかし，そのような$x$は一般にはひとつとは限らない．$h(x_1) = h(x_2) = y$を満たす$x_1 \ne x_2$があったとき，$f(x_1) \ne f(x_2)$であれば，$g(y)$の値は「どちらの$x$を選んだか」によって変わってしまい，$g$は写像の定義（\deref{写像}）の「ただひとつ対応させる」という条件を満たさない．この状況を整理したものが，次の定義である．

\phantomsection
\begin{definition}[
    sidebyside,
    righthand width=\mnsc{4.6cm},
    sidebyside gap=\mnsc{4mm},
  ]{well-defined}{well-defined}
  写像$f \colon X \to Z$と全射$h \colon X \to Y$を考える．任意の$x_1, x_2 \in X$に対して，
  \[
    h(x_1) = h(x_2) \limp f(x_1) = f(x_2)
  \]
  が成り立つとき，すなわち「$h$で同じ元にうつる$x$たちは，$f$でも同じ元にうつる」とき，$f(x)$の値は\emph{$y = h(x)$を満たす$x \in X$のとり方によらない}という．

  このとき，各$y \in Y$に対して，$h(x) = y$を満たす$x$をひとつ選んで
  \[
    g(y) \coloneqq f(x)
  \]
  と定めると，$g \colon Y \to Z$は写像であり，$f = g \circ h$を満たす．この$g$は\emph{well-defined}\index{welldefined@well-defined}（よく定義されている）であるという．

  \tcblower

  \centering
  \begin{tikzpicture}[mapfig, >=Latex]
    \node (X) at (0, 0)    {$X$};
    \node (Y) at (0, -2.4) {$Y$};
    \node (Z) at (2.6, -1.2) {$Z$};
    \draw[->, draw=black!70, line width=0.55pt, shorten <=3pt, shorten >=3pt]
      (X) -- (Y) node[midway, left=1pt] {$h$};
    \draw[->, draw=black!70, line width=0.55pt, shorten <=3pt, shorten >=3pt]
      (X) -- (Z) node[midway, above right=-1pt] {$f$};
    \draw[->, dashed, draw=figaccent, line width=0.7pt, shorten <=3pt, shorten >=3pt]
      (Y) -- (Z) node[midway, below right=-1pt, text=figaccent] {$g$};
  \end{tikzpicture}
  \par\smallskip
  {\small $f = g \circ h$}
\end{definition}

右の図は\emph{可換図式}\index{かかんずしき@可換図式}とよばれるもので，「$X$から$Z$へ，どの矢印をたどって行っても結果が同じ（$f = g \circ h$）」ことを表している．破線の$g$が「これから定めたい写像」であり，well-definedであるとは，この破線を実線に置き換えてよい，ということにほかならない．

\begin{remark}[「表し方によらない」との関係]
  \partref{用語の定義}で述べた「対象の表し方によらず値が定まる」という説明と，\deref{well-defined}とのつながりを確認しよう．$X$を「表し方」の集合，$Y$を「本来の対象」の集合とし，$h \colon X \to Y$を「表し方から対象を読み取る」写像だと思えばよい．同じ対象$y$は，$h(x) = y$を満たす$x$の個数だけ表し方をもつ．$f$は表し方$x$に対して定められた計算であり，
  「$f(x)$が表し方$x$のとり方によらない」ならば，$f$は本来の対象$y$に対する写像$g$を定めている\,---\,これが\deref{well-defined}のいうことである．
\end{remark}

\begin{example}{有理数の関数のwell-defined性}{有理数の関数のwell-defined性}
  $X = \{ (a, b) \in \mathbb{Z} \times \mathbb{Z} \mid b \ne 0 \}$を「分数の表し方」の集合，$Y = \mathbb{Q}$とし，
  \[
    h \colon X \ni (a, b) \mapsto a/b \in \mathbb{Q}
  \]
  とする．どの有理数も分数の形に表せるから，$h$は全射である．
  また，$h(1, 2) = h(2, 4)$のように，同じ有理数にうつる組は複数ある．
  \begin{enumerate}[(A)]
    \item $f \colon X \ni (a, b) \mapsto a + b \in \mathbb{Z}$とする．$h(1, 2) = h(2, 4)$であるが，$f(1, 2) = 3 \ne 6 = f(2, 4)$である．よって$f(a, b)$は$(a, b)$のとり方によるから，$g(a/b) = a + b$は写像$g \colon \mathbb{Q} \to \mathbb{Z}$を定めない（well-definedでない）．
    \item $f \colon X \ni (a, b) \mapsto 2a/b \in \mathbb{Q}$とする．$h(a, b) = h(a', b')$，すなわち$a/b = a'/b'$ならば，$ab' = a'b$より$2a/b = 2a'/b'$，すなわち$f(a, b) = f(a', b')$である．よって$g (a/b ) = 2a/b$はwell-definedであり，写像$g \colon \mathbb{Q} \to \mathbb{Q}$を定める．
  \end{enumerate}
  \partref{用語の定義}で述べた例が，まさにこの形をしていたことを確認してほしい．
\end{example}

\begin{remark}[「細かい」という言い回し]
  \deref{well-defined}の条件「$h(x_1) = h(x_2)$ならば$f(x_1) = f(x_2)$」は，$X$を「$h$の値が等しい元どうし」で分類したとき，その分類が「$f$の値が等しい元どうし」の分類より細かい，と言い表されることがある．同じ$h$の値をもつ元たちは，かならず同じ$f$の値をもつ（が，逆は要求しない）からである．同値関係や商集合の言葉を学ぶと，この定義はさらに簡潔に述べ直すことができる（\partref{関係}）．
\end{remark}

\phantomsection
\part{関係}

\begin{abstract}
この章では，「$=$」「$\leqq$」「$\subset$」のように$2$つの元の間に成り立つかどうかが定まる「関係」を，集合の言葉で定式化する．まず二項関係を直積の部分集合として定義し，反射律・対称律・反対称律・推移律といった性質を整理する．次に，「同じとみなす」ための同値関係を導入し，同値類と商集合を定義して，商集合が集合の分割を与えることを確認する．そのうえで，\partref{写像}で述べたwell-definedの概念を商集合上の写像として述べ直し，剰余類の演算を例として扱う．最後に，大小関係を一般化した順序関係に触れよう．
\end{abstract}

\phantomsection
\section{二項関係}

\phantomsection
\subsection{関係の定義}

実数どうしの「$x \leqq y$」，集合どうしの「$A \subset B$」，自然数どうしの「$a$は$b$を割り切る」のように，数学には$2$つの元の間に成り立つかどうかが定まる「関係」が数多く現れる．
写像を直積の部分集合として定義した（\deref{写像の厳密な定義}）のと同じように，関係も直積の部分集合として定義することができる．

\phantomsection
\begin{definition}{二項関係}{二項関係}
  $A$を集合とする．直積$A \times A$の部分集合$R$を，$A$上の\emph{二項関係}\index{にこうかんけい@二項関係}（または単に関係）という．
  $(a, b) \in R$であることを$a \mathrel{R} b$とかき，$a$と$b$は関係$R$にあるという．
\end{definition}

すなわち，関係$R$とは「関係にある組$(a, b)$をすべて集めた集合」のことである．実際の場面では，$R$を集合として書き下すことは少なく，$\leqq$や$\sim$などの記号を用いて$a \leqq b$，$a \sim b$のようにかくことが多い．

\phantomsection
\begin{example}{二項関係の例}{二項関係の例}
  \begin{enumerate}[(A)]
    \item 実数の大小関係$\leqq$は，$\{ (x, y) \in \mathbb{R} \times \mathbb{R} \mid x \leqq y \}$という$\mathbb{R}$上の関係である．
    \item 自然数$a$，$b$について，ある$k \in \mathbb{N}$が存在して$b = ka$となるとき，$a$は$b$を\emph{割り切る}\index{わりきる@割り切る}といい，$a \mid b$とかく．$\mid$は$\mathbb{N}$上の関係である．
    \item 集合$X$に対して，包含関係$\subset$はべき集合$\mathcal{P}(X)$上の関係である．
    \item $n \in \mathbb{N}$を固定する．整数$a$，$b$について，$a - b$が$n$の倍数であるとき，すなわちある整数$k$が存在して$a - b = nk$となるとき，$a$と$b$は$n$を法として\emph{合同}\index{ごうどう@合同}であるといい，
          \[
            a \equiv b \pmod{n}
          \]
          とかく．これは$\mathbb{Z}$上の関係である．$a \equiv b \pmod{n}$は，「$a$と$b$を$n$で割った余りが等しい」ことと同値である．
  \end{enumerate}
\end{example}

\begin{remark}[関係と写像]
  \deref{二項関係}を少し一般化して，直積$A \times B$の部分集合を「$A$から$B$への関係」とよぶことにすれば，写像の厳密な定義（\deref{写像の厳密な定義}）は，写像とは
  \[
    \lall{a \in A} \lsep \lexu{b \in B} \lst (a, b) \in G
  \]
  を満たす$A$から$B$への関係$G$のことである，と読める．すなわち写像は，関係のうち「各$a$に対して，関係にある$b$がただひとつ定まる」という特別なものである．関係のほうでは，ひとつの$a$に対して関係にある$b$がいくつあってもよく，ひとつもなくてもよい．
\end{remark}

\phantomsection
\subsection{関係の性質}

関係がもちうる性質のうち，基本的なものに名前をつけておこう．

\phantomsection
\begin{definition}{反射律・対称律・反対称律・推移律}{関係の性質}
  $R$を集合$A$上の関係とする．
  \begin{description}
    \item[反射律\index{はんしゃりつ@反射律}]
          任意の$a \in A$に対して，$a \mathrel{R} a$である．
    \item[対称律\index{たいしょうりつ@対称律}]
          任意の$a, b \in A$に対して，$a \mathrel{R} b$ならば$b \mathrel{R} a$である．
    \item[反対称律\index{はんたいしょうりつ@反対称律}]
          任意の$a, b \in A$に対して，$a \mathrel{R} b$かつ$b \mathrel{R} a$ならば$a = b$である．
    \item[推移律\index{すいいりつ@推移律}]
          任意の$a, b, c \in A$に対して，$a \mathrel{R} b$かつ$b \mathrel{R} c$ならば$a \mathrel{R} c$である．
  \end{description}
\end{definition}

\exref{二項関係の例}の関係について，これらの性質が成り立つかどうかをまとめると，次のようになる．

\begin{center}
  \begin{tabular}{|l|c|c|c|c|} \hline
    \textbf{関係}                      & \textbf{反射律} & \textbf{対称律} & \textbf{反対称律} & \textbf{推移律} \\ \hline
    $\mathbb{R}$上の$=$               & ○      & ○      & ○        & ○      \\ \hline
    $\mathbb{R}$上の$\leqq$           & ○      & ×      & ○        & ○      \\ \hline
    $\mathbb{R}$上の$<$               & ×      & ×      & ○        & ○      \\ \hline
    $\mathbb{N}$上の$\mid$            & ○      & ×      & ○        & ○      \\ \hline
    $\mathcal{P}(X)$上の$\subset$     & ○      & ×\footnotemark & ○        & ○      \\ \hline
    $\mathbb{Z}$上の$\equiv \pmod{n}$ & ○      & ○      & ×\footnotemark & ○      \\ \hline
  \end{tabular}
\end{center}
\addtocounter{footnote}{-1}
\footnotetext{$X \ne \varnothing$の場合．$X = \varnothing$ならば$\mathcal{P}(X) = \{\varnothing\}$であり，対称律も成り立つ．}
\stepcounter{footnote}
\footnotetext{$n \geqq 2$の場合．たとえば$0 \equiv n \pmod{n}$かつ$n \equiv 0 \pmod{n}$であるが$0 \ne n$である．}

$<$が反対称律を満たすのは，「$a < b$かつ$b < a$」を満たす$a$，$b$が存在せず，仮定が偽となるためである（空虚な真）．
このように，$4$つの性質は互いに独立であり，どの組合せを満たすかによって関係の性格が決まる．とくに重要なのが，次節で扱う同値関係と，最後の節で扱う順序関係である．

\phantomsection
\section{同値関係}

\phantomsection
\subsection{同値関係と同値類}

数学では，異なる対象を「ある観点からは同じもの」とみなしたい場面がしばしばある．たとえば，$1/2$と$2/4$は分数の表し方としては異なるが，有理数としては同じである．また，整数を$3$で割った余りだけに注目するならば，$1$と$4$と$7$は同じとみなせる．
このような「同じとみなす」関係がもつべき性質を抽出したものが，同値関係である．

\phantomsection
\begin{definition}{同値関係}{同値関係}
  集合$A$上の関係$\sim$が反射律・対称律・推移律を満たすとき，$\sim$を$A$上の\emph{同値関係}\index{どうちかんけい@同値関係}という．すなわち，任意の$a, b, c \in A$に対して，次が成り立つ：
  \begin{enumerate}[label=(\roman*),leftmargin=3.2em]
    \item $a \sim a$である．
    \item $a \sim b$ならば$b \sim a$である．
    \item $a \sim b$かつ$b \sim c$ならば$a \sim c$である．
  \end{enumerate}
\end{definition}

\phantomsection
\begin{prop}{合同は同値関係}{合同は同値関係}
  $n \in \mathbb{N}$とする．$n$を法とする合同$\equiv \pmod{n}$は，$\mathbb{Z}$上の同値関係である．
\end{prop}

\begin{tleftbar}
  \begin{proof}
    $a, b, c \in \mathbb{Z}$とする．
    反射律について，$a - a = n \cdot 0$であるから$a \equiv a \pmod{n}$である．
    対称律について，$a - b = nk$（$k \in \mathbb{Z}$）ならば$b - a = n(-k)$であるから，$b \equiv a \pmod{n}$である．
    推移律について，$a - b = nk$，$b - c = nl$（$k, l \in \mathbb{Z}$）ならば，$a - c = (a - b) + (b - c) = n(k + l)$であるから，$a \equiv c \pmod{n}$である．これが証明すべきことであった．
  \end{proof}
\end{tleftbar}

\phantomsection
\begin{example}{同値関係の例}{同値関係の例}
  \begin{enumerate}[(A)]
    \item 任意の集合$A$上の等号$=$は，同値関係である．
    \item ある大学の学生全体の集合上で，「$x$と$y$は誕生月が同じである」という関係は同値関係である．
    \item $X = \{ (a, b) \in \mathbb{Z} \times \mathbb{Z} \mid b \ne 0 \}$上の関係$\sim$を
          \[
            (a, b) \sim (c, d) \lequiv ad = bc
          \]
          で定める．これは「分数$a/b$と$c/d$が有理数として等しい」ことを，割り算を用いずに表したものであり，同値関係である．
          反射律と対称律は明らかである．推移律について，$(a, b) \sim (c, d)$かつ$(c, d) \sim (e, f)$とすると，$ad = bc$，$cf = de$であるから，
          \[
            adf = bcf = bde
          \]
          となり，$d(af - be) = 0$である．$d \ne 0$より$af = be$，すなわち$(a, b) \sim (e, f)$である．
  \end{enumerate}
  一方，$\mathbb{R}$上の$\leqq$は対称律を満たさないので，同値関係ではない．
\end{example}

同値関係が与えられると，互いに同値な元をひとまとめにすることができる．

\phantomsection
\begin{definition}{同値類}{同値類}
  $\sim$を集合$A$上の同値関係とする．$a \in A$に対して，$a$と同値な元全体の集合
  \[
    [a] \coloneqq \{ x \in A \mid x \sim a \}
  \]
  を，$a$の\emph{同値類}\index{どうちるい@同値類}という．また，$C$がある元の同値類であるとき，$a \in C$を満たす$a$を$C$の\emph{代表元}\index{だいひょうげん@代表元}という．
\end{definition}

\phantomsection
\begin{example}{剰余類}{剰余類}
  $\mathbb{Z}$上の同値関係$\equiv \pmod{3}$について，
  \begin{align*}
    [0] & = \{ \ldots, -6, -3, 0, 3, 6, \ldots \}, \\
    [1] & = \{ \ldots, -5, -2, 1, 4, 7, \ldots \}, \\
    [2] & = \{ \ldots, -4, -1, 2, 5, 8, \ldots \}
  \end{align*}
  である．$[3] = [0]$，$[4] = [1]$，$[-1] = [2]$のように，ほかの整数の同値類はこの$3$つのいずれかと一致する．$n$を法とする合同についての同値類を，$n$を法とする\emph{剰余類}\index{じょうよるい@剰余類}という．
\end{example}

\exref{剰余類}では，$3$つの同値類は互いに共通部分をもたず，あわせると$\mathbb{Z}$全体になっていた．これは一般に成り立つ．

\phantomsection
\begin{prop}{同値類の性質}{同値類の性質}
  $\sim$を集合$A$上の同値関係とし，$a, b \in A$とする．このとき，次が成り立つ：
  \begin{enumerate}[label=(\roman*),leftmargin=3.2em]
    \item $a \in [a]$である．
    \item $a \sim b$であることと，$[a] = [b]$であることは同値である．
    \item $[a] \ne [b]$ならば，$[a] \cap [b] = \varnothing$である．
  \end{enumerate}
\end{prop}

\begin{tleftbar}
  \begin{proof}
    (i)は反射律$a \sim a$から従う．

    (ii)について，$a \sim b$とする．$x \in [a]$ならば$x \sim a$であり，推移律より$x \sim b$，すなわち$x \in [b]$である．よって$[a] \subset [b]$である．対称律より$b \sim a$でもあるから，同様に$[b] \subset [a]$であり，$[a] = [b]$である．
    逆に$[a] = [b]$ならば，(i)より$a \in [a] = [b]$であるから，$a \sim b$である．

    (iii)について，対偶「$[a] \cap [b] \ne \varnothing$ならば$[a] = [b]$」を示す．$x \in [a] \cap [b]$とすると，$x \sim a$かつ$x \sim b$である．対称律より$a \sim x$であり，推移律より$a \sim b$である．よって(ii)より$[a] = [b]$である．これが証明すべきことであった．
  \end{proof}
\end{tleftbar}

(ii)より，同値類$[a]$は代表元$a$の選び方によらず定まる．たとえば\exref{剰余類}の$[1]$と$[4]$は，同じ集合を異なる代表元で表したものにすぎない．

\phantomsection
\subsection{商集合と分割}

同値類そのものを元とする集合を考えよう．

\phantomsection
\begin{definition}{商集合}{商集合}
  $\sim$を集合$A$上の同値関係とする．同値類全体の集合
  \[
    A / {\sim} \coloneqq \{ [a] \mid a \in A \}
  \]
  を，$A$の$\sim$による\emph{商集合}\index{しょうしゅうごう@商集合}という．また，写像
  \[
    \pi \colon A \ni a \mapsto [a] \in A / {\sim}
  \]
  を\emph{自然な射影}\index{しぜんなしゃえい@自然な射影}という．
\end{definition}

商集合の定義より，$A / {\sim}$のどの元も，ある$a \in A$を用いて$[a] = \pi(a)$と表せるから，自然な射影$\pi$は全射である．
また，\prref{同値類の性質}(ii)より，$a, a' \in A$に対して
\[
  \pi(a) = \pi(a') \lequiv a \sim a'
\]
である．つまり，$\pi$は「同値な元を同じ元にうつし，同値でない元を異なる元にうつす」写像である．

たとえば，$\mathbb{Z}$の$\equiv \pmod{3}$による商集合は$\{ [0], [1], [2] \}$であり，これを$\mathbb{Z}/3\mathbb{Z}$とかく．一般に，$n$を法とする合同による商集合を$\mathbb{Z}/n\mathbb{Z}$とかく．
\figref{剰余類と自然な射影}のように，自然な射影$\pi$は，ひとつの同値類に属する無数の整数をまとめて，$\mathbb{Z}/3\mathbb{Z}$のただひとつの元にうつす．

\begin{figure}[htbp]
  \centering
  \begin{tikzpicture}[
      >=Latex, font=\small, line cap=round, line join=round,
      cls/.style    ={draw=black!50, fill=black!4, line width=0.5pt,
                      rounded corners=6pt, minimum width=\mnsc{6.2cm},
                      minimum height=\mnsc{0.62cm}, inner sep=0pt},
      clsacc/.style ={cls, draw=figaccent!60, fill=figaccent!10},
      num/.style    ={font=\small, text=black!85},
      elem/.style   ={circle, fill=black!85, inner sep=1.3pt},
      proj/.style   ={-{Latex[length=1.8mm,width=1.3mm]}, draw=black!70,
                      line width=0.55pt, shorten <=3pt, shorten >=3pt},
    ]
    % Z を 3 つの同値類（剰余類）に切り分ける
    \draw[draw=black!60, line width=0.7pt, rounded corners=10pt]
      (-3.45,-1.55) rectangle (3.45,1.55);
    \node[font=\small, text=black!70, anchor=south west] at (-3.45,1.58) {$\mathbb{Z}$};
    \foreach \r/\y/\st in {0/0.9/cls, 1/0/clsacc, 2/-0.9/cls} {
      \node[\st] (C\r) at (0,\y) {};
      \foreach \k [count=\i from 0] in {-2,-1,0,1,2} {
        \pgfmathtruncatemacro{\v}{3*\k+\r}
        \node[num] at (-2.2+1.1*\i,\y) {$\v$};
      }
      \node[num] at (-2.85,\y) {$\cdots$};
      \node[num] at ( 2.85,\y) {$\cdots$};
    }
    % 商集合 Z/3Z
    \draw[draw=black!60, fill=black!4, line width=0.7pt]
      (6.1,0) ellipse [x radius=0.85, y radius=1.45];
    \node[font=\small, text=black!70] at (6.1,1.8) {$\mathbb{Z}/3\mathbb{Z}$};
    \foreach \r/\y in {0/0.9, 2/-0.9} {
      \node[elem] (P\r) at (6.1,\y) {};
      \node[anchor=west, xshift=2.5pt] at (P\r) {$[\r]$};
    }
    \node[elem, fill=figaccent] (P1) at (6.1,0) {};
    \node[anchor=west, xshift=2.5pt, text=figaccent] at (P1) {$[1]$};
    % 自然な射影 π
    \draw[proj] (C0.east) -- (P0);
    \draw[proj, draw=figaccent] (C1.east) -- (P1);
    \draw[proj] (C2.east) -- (P2);
    \node[font=\small, text=black!70] at (4.4,1.3) {$\pi$};
  \end{tikzpicture}
  \caption{$\mathbb{Z}$を$3$を法とする剰余類に分割し，各類を$\mathbb{Z}/3\mathbb{Z}$の$1$点にうつす}
  \label{fig:剰余類と自然な射影}
\end{figure}

商集合が$A$をどのように「切り分けて」いるかを述べるために，次の用語を導入する．

\phantomsection
\begin{definition}{分割}{分割}
  集合$A$の部分集合からなる集合$\mathcal{C}$が次の$3$つを満たすとき，$\mathcal{C}$を$A$の\emph{分割}\index{ぶんかつ@分割}という：
  \begin{enumerate}[label=(\roman*),leftmargin=3.2em]
    \item $\mathcal{C}$の各元は空集合でない．
    \item $C, C' \in \mathcal{C}$が$C \ne C'$ならば，$C \cap C' = \varnothing$である．
    \item $A$の各元は，$\mathcal{C}$のいずれかの元に属する．
  \end{enumerate}
\end{definition}

\phantomsection
\begin{prop}{商集合は分割を与える}{商集合は分割を与える}
  $\sim$を集合$A$上の同値関係とすると，商集合$A / {\sim}$は$A$の分割である．
\end{prop}

\begin{tleftbar}
  \begin{proof}
    \prref{同値類の性質}(i)より，各同値類$[a]$は$a$を含むから空集合でなく，また$A$の各元$a$は同値類$[a]$に属する．よって分割の条件(i)，(iii)が成り立つ．条件(ii)は\prref{同値類の性質}(iii)そのものである．これが証明すべきことであった．
  \end{proof}
\end{tleftbar}

\begin{remark}[同値関係と分割は同じものの$2$つの見方]
  \prref{商集合は分割を与える}の逆も成り立つ．すなわち，$\mathcal{C}$を$A$の分割とし，$a, b \in A$に対して
  \[
    a \sim b \lequiv \text{$a$と$b$は$\mathcal{C}$の同じ元に属する}
  \]
  と定めると，$\sim$は$A$上の同値関係であり，その商集合$A / {\sim}$は$\mathcal{C}$に一致する．
  したがって，「同値関係を与えること」と「分割を与えること」は，同じことを異なる言葉で述べたものにほかならない．同値関係は「どの元とどの元を同じとみなすか」という局所的な見方であり，分割は「全体をどのように切り分けるか」という大域的な見方である．
\end{remark}

\phantomsection
\section{商集合上の写像}

\phantomsection
\subsection{well-definedふたたび}

商集合$A / {\sim}$上の写像を定めるとき，もっとも自然な方法は，同値類$[a]$の行き先を代表元$a$を用いて定めることである．しかし，同値類は複数の代表元をもちうるので，こうして定めた行き先が代表元の選び方によらないことを確かめなければならない．
これは，\subsecref{well-defined}で扱った状況において，全射$h$を自然な射影$\pi$としたものにほかならない．

\phantomsection
\begin{theorem}{商集合上の写像}{商集合上の写像}
  $\sim$を集合$A$上の同値関係，$\pi \colon A \to A / {\sim}$を自然な射影とし，写像$f \colon A \to Z$が
  \[
    a \sim a' \limp f(a) = f(a') \qquad (a, a' \in A)
  \]
  を満たすとする．このとき，写像$g \colon A / {\sim} \to Z$であって$f = g \circ \pi$を満たすものが，ただひとつ存在する．すなわち，
  \[
    g([a]) \coloneqq f(a) \qquad (a \in A)
  \]
  はwell-definedである．
\end{theorem}

\begin{tleftbar}
  \begin{proof}
    自然な射影$\pi$は全射であり，$a, a' \in A$に対して，
    \[
      \pi(a) = \pi(a') \limp a \sim a' \limp f(a) = f(a')
    \]
    が成り立つ．よって\deref{well-defined}より，$g([a]) \coloneqq f(a)$は写像$g \colon A / {\sim} \to Z$を定め，$f = g \circ \pi$を満たす．

    一意性を示す．$g, g' \colon A / {\sim} \to Z$がともに$f = g \circ \pi = g' \circ \pi$を満たすとする．任意の$C \in A / {\sim}$に対して，$\pi$は全射であるから$C = \pi(a)$を満たす$a \in A$が存在し，
    \[
      g(C) = g(\pi(a)) = f(a) = g'(\pi(a)) = g'(C)
    \]
    である．よって$g = g'$である．これが証明すべきことであった．
  \end{proof}
\end{tleftbar}

\phantomsection
\begin{example}{剰余類の演算}{剰余類の演算}
  $n \in \mathbb{N}$とし，$\mathbb{Z}/n\mathbb{Z}$における和を
  \[
    [a] + [b] \coloneqq [a + b] \qquad (a, b \in \mathbb{Z})
  \]
  で定めたい．これがwell-definedであること，すなわち右辺が代表元$a$，$b$の選び方によらないことを確かめよう．
  $[a] = [a']$，$[b] = [b']$とすると，\prref{同値類の性質}(ii)より$a - a' = nk$，$b - b' = nl$（$k, l \in \mathbb{Z}$）と表せる．このとき
  \[
    (a + b) - (a' + b') = n(k + l)
  \]
  であるから，$a + b \equiv a' + b' \pmod{n}$，すなわち$[a + b] = [a' + b']$である．よって和はwell-definedである．

  たとえば$\mathbb{Z}/3\mathbb{Z}$では，$[1] + [2] = [3] = [0]$，$[2] + [2] = [4] = [1]$である．これは「$3$で割った余り」の足し算にほかならない．
\end{example}

一方，代表元を用いた定義がいつもうまくいくとは限らない．
たとえば，$\mathbb{Z}/3\mathbb{Z}$から$\mathbb{Z}$への写像を「$g([a]) \coloneqq a$」で定めようとしても，$[0] = [3]$であるのに$0 \ne 3$であるから，$g([0])$の値が代表元の選び方によって変わってしまう．よってこの$g$はwell-definedでない．

\begin{remark}[代表元を用いた定義の確認]
  商集合上の写像や演算を代表元を用いて定めるときには，それが代表元の選び方によらないことを，必ず確認しなければならない．\thref{商集合上の写像}は，この確認さえ済めば写像がただひとつ定まることを保証している．
  \exref{有理数の関数のwell-defined性}は，まさにこの構造をしていた．実際，\exref{同値関係の例}(C)の同値関係$\sim$について，$X / {\sim}$の各元$[(a, b)]$を有理数$a/b$と同一視すれば，$h \colon X \to \mathbb{Q}$は自然な射影$\pi$にあたる．
\end{remark}

\begin{mycolumn}
  同値関係と商集合は，数の体系を構成する際にも用いられる．
  たとえば整数は，自然数の組$(a, b)$（$a, b \in \mathbb{N}$）を「$a - b$」の表し方とみなし，$(a, b) \sim (c, d) \lequiv a + d = b + c$という同値関係による商集合として構成できる．引き算を用いずに「差が等しい」ことを表しているのである．
  同様に，有理数は\exref{同値関係の例}(C)の同値関係による商集合として構成でき，実数も有理数の列を用いた同値関係（コーシー列の同値類）によって構成できる．
  このように，「新しい数を，古い数の組の同値類として作る」という考え方は，現代数学の基本的な手法のひとつである．
\end{mycolumn}

\phantomsection
\section{順序関係}

最後に，大小関係を一般化した関係を扱う．

\phantomsection
\begin{definition}{半順序・全順序}{半順序・全順序}
  集合$A$上の関係$\preceq$が反射律・反対称律・推移律を満たすとき，$\preceq$を$A$上の\emph{半順序}\index{はんじゅんじょ@半順序}（または順序）といい，組$(A, \preceq)$を\emph{半順序集合}\index{はんじゅんじょしゅうごう@半順序集合}という．
  さらに，任意の$a, b \in A$に対して$a \preceq b$または$b \preceq a$が成り立つとき，$\preceq$を\emph{全順序}\index{ぜんじゅんじょ@全順序}という．
\end{definition}

\phantomsection
\begin{example}{半順序の例}{半順序の例}
  \begin{enumerate}[(A)]
    \item $\mathbb{R}$上の$\leqq$は全順序である．
    \item 集合$X$のべき集合$\mathcal{P}(X)$上の$\subset$は半順序である．しかし，$X = \{1, 2\}$のとき，$\{1\} \subset \{2\}$も$\{2\} \subset \{1\}$も成り立たないので，全順序ではない．
    \item $\mathbb{N}$上の$\mid$は半順序である．しかし，$2 \mid 3$も$3 \mid 2$も成り立たないので，全順序ではない．
  \end{enumerate}
\end{example}

\begin{wrapfigure}[9]{r}{\mnsc{5.6cm}}
  \centering
  \begin{tikzpicture}[
      x=1.3cm, y=0.95cm, font=\small,
      ps/.style     ={inner sep=2pt, text=black!85},
      psacc/.style  ={ps, text=figaccent},
      cover/.style  ={draw=black!45, line width=0.5pt, shorten <=1pt, shorten >=1pt},
      chain/.style  ={draw=figaccent, line width=1.4pt, shorten <=1pt, shorten >=1pt},
    ]
    \node[psacc] (e)   at ( 0,0) {$\varnothing$};
    \node[psacc] (1)   at (-1,1) {$\{1\}$};
    \node[ps]    (2)   at ( 0,1) {$\{2\}$};
    \node[ps]    (3)   at ( 1,1) {$\{3\}$};
    \node[psacc] (12)  at (-1,2) {$\{1,2\}$};
    \node[ps]    (13)  at ( 0,2) {$\{1,3\}$};
    \node[ps]    (23)  at ( 1,2) {$\{2,3\}$};
    \node[psacc] (123) at ( 0,3) {$\{1,2,3\}$};
    \foreach \a/\b in {e/2, e/3, 1/13, 2/12, 2/23, 3/13, 3/23, 13/123, 23/123}
      \draw[cover] (\a) -- (\b);
    % 全順序をなす部分（鎖）を太線で示す
    \foreach \a/\b in {e/1, 1/12, 12/123}
      \draw[chain] (\a) -- (\b);
  \end{tikzpicture}
  \caption{$(\mathcal{P}(\{1,2,3\}), \subset)$のハッセ図}
  \label{fig:ハッセ図}
\end{wrapfigure}
半順序では，$2$つの元が「比べられない」ことがありうる．日常的な「大小」の感覚は全順序にもとづいているので，半順序を扱うときにはこの点に注意が必要である．

有限の半順序集合は，$a \preceq b$かつ$a \ne b$で，その間にほかの元がないときに，$a$を下，$b$を上に置いて線で結んだ図で表すと見通しがよい．このような図を\emph{ハッセ図}\index{はっせず@ハッセ図}という．
\figref{ハッセ図}は$\mathcal{P}(\{1, 2, 3\})$の包含関係のハッセ図である．線を下から上へたどって行ける$2$元は比べることができ，$\{1\}$と$\{2, 3\}$のようにたどって行けない$2$元は比べられない．
また，太線で示した$\varnothing \subset \{1\} \subset \{1, 2\} \subset \{1, 2, 3\}$のように，一列に並ぶ元だけを取り出せば，その範囲では全順序になっている．

\begin{remark}[最大元の一意性と反対称律]
  \deref{最大値・最小値}にならって，半順序集合$(A, \preceq)$の元$M$が「任意の$a \in A$に対して$a \preceq M$」を満たすとき，$M$を$A$の最大元という．
  \prref{最大値の一意性}の証明では，「$M' \leqq M$かつ$M \leqq M'$ならば$M = M'$」という推論を用いたが，これはまさに反対称律である．したがって，同じ証明により，半順序集合の最大元も存在するとすればただひとつである．
  このように，ある性質の証明が実際にはどの仮定を用いているのかを見極めると，その性質が成り立つ範囲を広げることができる．
\end{remark}

\clearpage


\phantomsection

\part*{参考文献}

\noindent
本書は\cite{sugiura1}，\cite{nakajima}，\cite{kaneko}，\cite{arai}，\cite{saito}，\cite{akasetsuya}，\cite{sakiyama2023}を参考にした．

\medskip

\printbibliography[heading=none]

% 索引
% 参考文献の次のページ（偶数ページ）から索引を開始する．
% 通常の \part は右ページ開始だが，索引だけは右ページ送りを無効にする．

\clearpage
\thispagestyle{empty}
\setcounter{section}{6}
\renewcommand{\indexname}{索引}

\begingroup
  \makeatletter
  \def\part@newpage{\clearpage\thispagestyle{empty}}
  \makeatother
  \let\section\part
  \pagestyle{empty}
  \printindex
\endgroup

% 裏表紙が \BookTotalPages ページ目（偶数ページ）に来るよう白紙を補充する．
% 末尾 2 ページ＝裏表紙の裏（白紙）＋裏表紙．
% \ReadonlyShipoutCounter は表紙を含めた出力済みの物理ページ数．
\clearpage
\pagestyle{empty}
\makeatletter
\@tempcnta=\numexpr\BookTotalPages-2-\ReadonlyShipoutCounter\relax
\ifnum\@tempcnta<0
  \PackageWarningNoLine{製本}{本文が \BookTotalPages\space ページを超えています}
\fi
\loop\ifnum\@tempcnta>0
  \thispagestyle{empty}\null\clearpage
  \advance\@tempcnta\m@ne
\repeat
\makeatother
\MathNoteBackCoverBlankfalse

\MathNoteBackCover{Foundations of Mathematics}{2026 Edition}

\end{document}
